% Leçon 24 — Expression orale : démocratie — Chinois Tle A — Cours signé OnBuch+
% Programme officiel MINESEC. Compiler avec : tectonic ZH24-oral-democratie.tex
\newcommand{\DOCMATIERE}{Chinois}
\newcommand{\DOCNIVEAU}{Tle A}
\newcommand{\DOCMODULE}{Module 3 : Citoyenneté et ouverture au monde}
\newcommand{\DOCLECON}{ZH24}
\newcommand{\DOCTITRE}{Leçon 24 — Expression orale : démocratie}
% OnBuch+ — préambule « cours signés » ENRICHI (compilé avec tectonic / XeLaTeX)
% Système visuel « Pop » d'OnBuch+ : fond crème (jamais blanc), bordures encre
% épaisses, ombres « dures » décalées, titres Archivo Black en capitales.
% Version MATHÉMATIQUES : graphes (pgfplots/tikz), boîtes pédagogiques
% (définition, propriété, méthode, attention, le savais-tu, exercice résolu,
% exercice, TP, bilan), page de garde et signature OnBuch+ ancrée en bas.
%
% Macros attendues dans main.tex AVANT \input{preamble} :
%   \DOCMATIERE  \DOCNIVEAU  \DOCMODULE  \DOCLECON (numéro)  \DOCTITRE
\documentclass[11pt]{article}

\usepackage{fontspec}
\usepackage[french]{babel} % français = langue principale ; le chinois via \zh{..} / textebox (xeCJK)
\usepackage{xeCJK}
\usepackage{pifont} % \ding{51} \ding{55} utilisés par les modèles
\usepackage[a4paper,margin=2cm,top=3.4cm,bottom=2.8cm,headheight=1.4cm,footskip=1.3cm]{geometry}
\usepackage{amsmath,amssymb,amsfonts}
\usepackage{mathtools} % \xmapsto, \coloneqq... souvent utilisés par les modèles
\usepackage{cancel} % simplifications barrées (bilans de forces)
% \abs{z} (module, valeur absolue) et \norm{u} (norme), utilisés par les modèles
\providecommand{\abs}{}\let\abs\relax\DeclarePairedDelimiter{\abs}{\lvert}{\rvert}
\providecommand{\norm}{}\let\norm\relax\DeclarePairedDelimiter{\norm}{\lVert}{\rVert}
\usepackage[table]{xcolor} % [table] : fournit aussi \rowcolors (alternance de lignes)
\usepackage{pagecolor}
\usepackage{tikz}
\usetikzlibrary{arrows.meta,positioning,calc,shapes.geometric,shapes.misc,shapes.arrows,decorations.pathmorphing,decorations.pathreplacing,decorations.markings,patterns,shadows,mindmap,trees,fit,backgrounds,matrix,intersections,angles,quotes}
\usepackage{pgfplots}
\usepackage[version=4]{mhchem} % équations nucléaires \ce{^{235}_{92}U}
\usepackage[european,siunitx]{circuitikz} % schémas électriques
\pgfplotsset{compat=1.18}
\usepgfplotslibrary{fillbetween,groupplots} % aires entre courbes (calcul intégral)
\usepackage{enumitem}
\usepackage{fancyhdr}
% [most] charge toutes les bibliothèques tcolorbox (listings, minted, raster,
% documentation...) et gonfle inutilement la mémoire TeX sur un document long
% (~300 boîtes) : on ne charge que ce qui est réellement utilisé.
\usepackage[skins,breakable]{tcolorbox}
\usepackage{graphicx}
\usepackage{colortbl}
\usepackage{booktabs}
\usepackage{tabularx}
\usepackage{diagbox} % case d'en-tête diagonale des tableaux à double entrée
% Colonnes extensibles centrée / gauche / droite (C, L, R), souvent utilisées par les modèles
\newcolumntype{C}{>{\centering\arraybackslash}X}
\newcolumntype{L}{>{\raggedright\arraybackslash}X}
\newcolumntype{R}{>{\raggedleft\arraybackslash}X}
\usepackage{array}
\usepackage{multicol}
\usepackage{multirow}
\usepackage{siunitx}
% parse-numbers=false : accepte aussi \SI{2,5 \times 10^{-3}}{...}
\sisetup{locale=FR,output-decimal-marker={,},group-minimum-digits=4,parse-numbers=false}
\DeclareSIUnit\ohm{\ensuremath{\symup{\Omega}}} % U+2126 absent de la police
\DeclareSIUnit\division{div} % réglages d'oscilloscope (ms/div, V/div)
\DeclareSIUnit\tr{tr} % tours (vitesse de rotation, ex. \SI{1500}{\tr\per\minute})
\DeclareSIUnit\molar{M} % concentration molaire (ex. \SI{3}{\molar}, \SI{1}{\micro\molar})
\DeclareSIUnit\an{an} % année (le modèle confond parfois avec \year, un compteur TeX) — ex. \SI{5}{\kilo\meter\per\an}
\DeclareSIUnit\FCFA{FCFA} % franc CFA (ex. \SI{500}{\FCFA\per\kilo\gram})
\DeclareSIUnit\degreeBrix{\textdegree Brix} % teneur en sucre d'un sirop/jus (réfractométrie)
\DeclareSIUnit\month{mois} % le modèle utilise parfois \month, un compteur TeX, comme unité de durée
\usepackage{qrcode}
% \degree : commande fréquemment utilisée par les modèles pour un angle ou une
% température en dehors de \SI{}{} (non fournie par les paquets chargés).
\providecommand{\degree}{^{\circ}}
% \qed (fin de démonstration), utilisé par les modèles sans amsthm
\providecommand{\qed}{\hfill$\square$}
% \barre : double barre des valeurs interdites dans un tableau de variations (tkz-tab)
\providecommand{\barre}{\ensuremath{\|}}
% environnement proof (amsthm non chargé) utilisé par les modèles
\ifdefined\proof\else\newenvironment{proof}[1][Démonstration]{\par\noindent\textit{#1.}\ }{\qed\par}\fi
\usepackage{lastpage}
\usepackage{hyperref}
\hypersetup{hidelinks}

% ============================================================
% Palette « Pop » OnBuch+ — mêmes hex que la classe P (plus_ui.dart)
% ============================================================
\definecolor{poporange}{HTML}{F59321}
\definecolor{popdark}{HTML}{E07A0C}
\definecolor{popcream}{HTML}{FFF6E8}
\definecolor{popink}{HTML}{1C1714}
\definecolor{poppurple}{HTML}{7A5AE0}
\definecolor{popgreen}{HTML}{2E9E6A}
\definecolor{poppink}{HTML}{E0457B}
\definecolor{popblue}{HTML}{2F7DE1}
\definecolor{popgold}{HTML}{C98A12}
\definecolor{popmuted}{HTML}{8A7F76}
\definecolor{popline}{HTML}{EADFD2}
\definecolor{popgray}{HTML}{8A7F76}
\colorlet{popgrey}{popgray}
% Alias tolérants (le modèle nomme parfois les couleurs différemment)
\colorlet{popwhite}{white}
\colorlet{popblack}{popink}
\colorlet{popred}{poppink}
\colorlet{popyellow}{popgold}
% Teintes claires (fonds de tableaux/encadrés) — le modèle nomme parfois
% les couleurs avec un suffixe L (ex. popblueL) sans les avoir définies.
\colorlet{poporangeL}{poporange!20}
\colorlet{popdarkL}{popdark!20}
\colorlet{poppurpleL}{poppurple!20}
\colorlet{popgreenL}{popgreen!20}
\colorlet{poppinkL}{poppink!20}
\colorlet{popblueL}{popblue!20}
\colorlet{popgoldL}{popgold!20}
\colorlet{popredL}{popred!20}
\colorlet{popgrayL}{popgray!20}
% Teintes claires pour fonds de tableaux / zones
\colorlet{popblueL}{popblue!10!white}
\colorlet{poporangeL}{poporange!14!white}
\colorlet{popgreenL}{popgreen!12!white}
\colorlet{poppurpleL}{poppurple!12!white}
\colorlet{poppinkL}{poppink!12!white}
\colorlet{popgoldL}{popgold!14!white}

\pagecolor{popcream}

% ============================================================
% Typographie
% ============================================================
\defaultfontfeatures{Ligatures=TeX}
\setmainfont{PlusJakartaSans-Medium.ttf}[Path=../../../fonts/,BoldFont=PlusJakartaSans-Bold.ttf]
% Symboles, API et emojis absents de la police principale : repli sur DejaVu Sans (caractères actifs)
\newfontface\symfont{DejaVuSans.ttf}[Path=../../../fonts/]
\makeatletter
\newcommand\symactive[1]{\catcode#1=13 \begingroup\lccode`\~=#1 \lowercase{\endgroup\def~}{{\symfont\char#1}}}
\newcommand\symblank[1]{\catcode#1=13 \begingroup\lccode`\~=#1 \lowercase{\endgroup\def~}{}}
\newcommand\symhyph[1]{\catcode#1=13 \begingroup\lccode`\~=#1 \lowercase{\endgroup\def~}{\mbox{-}}}
\newcount\symc
\newcommand\symrange[2]{\symc=#1 \@whilenum\symc<#2 \do{\expandafter\symactive\expandafter{\the\symc}\advance\symc by 1 }}
\newcommand\symblankrange[2]{\symc=#1 \@whilenum\symc<#2 \do{\expandafter\symblank\expandafter{\the\symc}\advance\symc by 1 }}
\makeatother
\symrange{"0250}{"02FF}\symactive{"2713}\symactive{"2714}\symactive{"2717}\symactive{"2718}\symactive{"2610}\symactive{"2611}\symactive{"2612}
\symrange{"2600}{"27BF}
\catcode"2705=13 \begingroup\lccode`\~="2705 \lowercase{\endgroup\def~}{{\symfont\char"2713}}\symblank{"26BD}\symblank{"26A1}
\symrange{"2460}{"257F}\symactive{"223C}\symactive{"2248}\symactive{"2260}
\symactive{"01F8}\symactive{"01F9}\symactive{"1E3E}\symactive{"1E3F}
\symhyph{"2011}\symhyph{"2E1A}\symblankrange{"1F300}{"1FAFF}\symblank{"FE0F}
% Caractères chinois : WenQuanYi Zen Hei (sous-ensemble GB2312 + pinyin), gras/italique simulés
\setCJKmainfont{WenQuanYiZenHei-subset.ttf}[Path=../../../fonts/,AutoFakeBold=3,AutoFakeSlant=0.2]
\xeCJKsetup{CJKspace=false}
% Pinyin : voyelles à caron (ǎ ǐ ǒ ǔ ǖ ǘ ǚ ǜ) absentes de la police principale -> caractères actifs
\newfontface\pyfont{WenQuanYiZenHei-subset.ttf}[Path=../../../fonts/]
\newcommand\pyactive[1]{\catcode#1=13 \begingroup\lccode`\~=#1 \lowercase{\endgroup\def~}{{\pyfont\char#1}}}
\pyactive{"01CD}\pyactive{"01CE}\pyactive{"01CF}\pyactive{"01D0}\pyactive{"01D1}\pyactive{"01D2}\pyactive{"01D3}\pyactive{"01D4}\pyactive{"01D5}\pyactive{"01D6}\pyactive{"01D7}\pyactive{"01D8}\pyactive{"01D9}\pyactive{"01DA}\pyactive{"01DB}\pyactive{"01DC}
% Maths en sans-serif (Fira Math) avec chiffres et lettres droites de Plus
% Jakarta, pour que les formules se fondent dans le texte.
\usepackage[math-style=ISO,bold-style=ISO]{unicode-math}
\setmathfont{FiraMath-Regular.otf}[Path=../../../fonts/]
\setmathfont{PlusJakartaSans-Medium.ttf}[Path=../../../fonts/,range={up/num,up/{Latin,latin},\mathup}]
\usepackage{icomma} % virgule décimale sans espace en mode math
% Caractères mathématiques Unicode absents des polices de texte (Plus Jakarta,
% Archivo Black) : rendus par leur équivalent mathématique, en texte comme en
% formule — sinon ils s'affichent en carré vide (ex. « Arithmétique dans ℤ »).
\usepackage{newunicodechar}
\newunicodechar{ℝ}{\ensuremath{\mathbb{R}}}
\newunicodechar{ℤ}{\ensuremath{\mathbb{Z}}}
\newunicodechar{ℂ}{\ensuremath{\mathbb{C}}}
\newunicodechar{ℕ}{\ensuremath{\mathbb{N}}}
\newunicodechar{ℚ}{\ensuremath{\mathbb{Q}}}
\newunicodechar{𝔻}{\ensuremath{\mathbb{D}}}
\newunicodechar{α}{\ensuremath{\alpha}}
\newunicodechar{β}{\ensuremath{\beta}}
\newunicodechar{γ}{\ensuremath{\gamma}}
\newunicodechar{δ}{\ensuremath{\delta}}
\newunicodechar{ε}{\ensuremath{\varepsilon}}
\newunicodechar{θ}{\ensuremath{\theta}}
\newunicodechar{λ}{\ensuremath{\lambda}}
\newunicodechar{μ}{\ensuremath{\mu}}
\newunicodechar{σ}{\ensuremath{\sigma}}
\newunicodechar{τ}{\ensuremath{\tau}}
\newunicodechar{φ}{\ensuremath{\varphi}}
\newunicodechar{ψ}{\ensuremath{\psi}}
\newunicodechar{ω}{\ensuremath{\omega}}
\newunicodechar{Σ}{\ensuremath{\Sigma}}
% majuscules produites par \MakeUppercase dans les titres (α-aminés -> Α)
\newunicodechar{Α}{\ensuremath{\alpha}}
\newunicodechar{Β}{\ensuremath{\beta}}
\newunicodechar{Γ}{\ensuremath{\Gamma}}
\newunicodechar{Δ}{\ensuremath{\Delta}}
\newunicodechar{Ε}{\ensuremath{\varepsilon}}
\newunicodechar{Θ}{\ensuremath{\Theta}}
\newunicodechar{Λ}{\ensuremath{\Lambda}}
\newunicodechar{Μ}{\ensuremath{\mu}}
\newunicodechar{Τ}{\ensuremath{\tau}}
\newunicodechar{Φ}{\ensuremath{\Phi}}
\newunicodechar{Ψ}{\ensuremath{\Psi}}
\newunicodechar{Ω}{\ensuremath{\Omega}}
\newunicodechar{↦}{\ensuremath{\mapsto}}
\newunicodechar{∈}{\ensuremath{\in}}
\newunicodechar{∉}{\ensuremath{\notin}}
\newunicodechar{∧}{\ensuremath{\wedge}}
\newunicodechar{⇔}{\ensuremath{\Leftrightarrow}}
\newunicodechar{⟺}{\ensuremath{\Longleftrightarrow}}
\newunicodechar{⇒}{\ensuremath{\Rightarrow}}
\newunicodechar{⟹}{\ensuremath{\Longrightarrow}}
\newunicodechar{∘}{\ensuremath{\circ}}
\newunicodechar{∪}{\ensuremath{\cup}}
\newunicodechar{∩}{\ensuremath{\cap}}
\newunicodechar{≡}{\ensuremath{\equiv}}
\newunicodechar{⊂}{\ensuremath{\subset}}
\newunicodechar{⊥}{\ensuremath{\perp}}
\newunicodechar{∃}{\ensuremath{\exists}}
\newunicodechar{∀}{\ensuremath{\forall}}
\newunicodechar{∛}{\ensuremath{\sqrt[3]{\,}}}
\newunicodechar{∜}{\ensuremath{\sqrt[4]{\,}}}
\newunicodechar{⃗}{\ensuremath{{}^{\to}}}
\newunicodechar{ⁿ}{\ifmmode^{n}\else\textsuperscript{\ensuremath{n}}\fi}
\newunicodechar{ˣ}{\ifmmode^{x}\else\textsuperscript{\ensuremath{x}}\fi}
\newunicodechar{ᵇ}{\ifmmode^{b}\else\textsuperscript{\ensuremath{b}}\fi}
\newunicodechar{ʳ}{\ifmmode^{r}\else\textsuperscript{\ensuremath{r}}\fi}
\newunicodechar{ᵘ}{\ifmmode^{u}\else\textsuperscript{\ensuremath{u}}\fi}
\newunicodechar{ʸ}{\ifmmode^{y}\else\textsuperscript{\ensuremath{y}}\fi}
\newunicodechar{ᵃ}{\ifmmode^{a}\else\textsuperscript{\ensuremath{a}}\fi}
\newunicodechar{ᵉ}{\ifmmode^{e}\else\textsuperscript{\ensuremath{e}}\fi}
\newunicodechar{ᵏ}{\ifmmode^{k}\else\textsuperscript{\ensuremath{k}}\fi}
\newunicodechar{ᵖ}{\ifmmode^{p}\else\textsuperscript{\ensuremath{p}}\fi}
\newunicodechar{ᴺ}{\ifmmode^{N}\else\textsuperscript{\ensuremath{N}}\fi}
\newunicodechar{ᶜ}{\ifmmode^{c}\else\textsuperscript{\ensuremath{c}}\fi}
\newunicodechar{ʰ}{\ifmmode^{h}\else\textsuperscript{\ensuremath{h}}\fi}
\newunicodechar{ᵠ}{\ifmmode^{q}\else\textsuperscript{\ensuremath{q}}\fi}
\newunicodechar{ₙ}{\ifmmode_{n}\else\textsubscript{\ensuremath{n}}\fi}
\newunicodechar{ₐ}{\ifmmode_{a}\else\textsubscript{\ensuremath{a}}\fi}
\newunicodechar{ᵢ}{\ifmmode_{i}\else\textsubscript{\ensuremath{i}}\fi}
\newunicodechar{ᵣ}{\ifmmode_{r}\else\textsubscript{\ensuremath{r}}\fi}
\newunicodechar{ₖ}{\ifmmode_{k}\else\textsubscript{\ensuremath{k}}\fi}
\newunicodechar{ᵦ}{\ifmmode_{\beta}\else\textsubscript{\ensuremath{\beta}}\fi}
\newunicodechar{ₓ}{\ifmmode_{x}\else\textsubscript{\ensuremath{x}}\fi}
\newfontfamily\popdisplay{ArchivoBlack-Regular.ttf}[Path=../../../fonts/]
\newfontfamily\poplogo{SpaceGrotesk-Bold.ttf}[Path=../../../fonts/]
\newfontfamily\popsemi{PlusJakartaSans-SemiBold.ttf}[Path=../../../fonts/]
\newfontfamily\popxbold{PlusJakartaSans-ExtraBold.ttf}[Path=../../../fonts/]
\newfontfamily\popxboldls{PlusJakartaSans-ExtraBold.ttf}[Path=../../../fonts/,LetterSpace=3.0]

\setlist[enumerate]{leftmargin=1.7em,itemsep=5pt,topsep=4pt}
\setlist[itemize]{leftmargin=1.7em,itemsep=4pt,topsep=4pt,label={\color{poporange}\textbullet}}
\setlength{\parindent}{0pt}
\setlength{\parskip}{5pt}
\renewcommand{\arraystretch}{1.25}

% pgfplots : style Pop par défaut
\pgfplotsset{
  every axis/.append style={axis line style={popink,line width=1pt},tick style={popink},
    grid style={popline,line width=0.5pt},label style={font=\small},tick label style={font=\footnotesize}},
  popcurve/.style={poporange,line width=1.8pt,no markers,smooth},
}
% Même style disponible dans un \draw TikZ simple (hors pgfplots)
\tikzset{popcurve/.style={poporange,line width=1.8pt}}
\tikzset{popgrid/.style={popline,very thin}}
\colorlet{popcurve}{poporange} % le nom du style est parfois utilisé comme couleur

% ============================================================
% Pill / squiggle
% ============================================================
\newcommand{\poppill}[3]{%
  \tikz[baseline=-0.55ex]\node[draw=popink,line width=1.2pt,rounded corners=2.6pt,%
    fill=#2,inner xsep=6.5pt,inner ysep=3pt]{%
    \popxboldls\fontsize{7}{7}\selectfont\color{#3}\MakeUppercase{#1}};%
}
\newcommand{\popsquiggle}[1]{%
  \tikz[baseline=-0.4ex]{
    \draw[#1,line width=2.6pt,line cap=round]
      plot [smooth] coordinates {(0,0.09) (0.34,0.01) (0.68,0.21) (1.02,0.01) (1.36,0.10)};
  }%
}
% Mot-clé surligné (marqueur orange sous le texte)
\newcommand{\cle}[1]{{\popxbold\color{popdark}#1}}

% ============================================================
% En-tête / pied
% ============================================================
\pagestyle{fancy}
\fancyhf{}
\renewcommand{\headrulewidth}{0pt}
\renewcommand{\footrulewidth}{0pt}
\lhead{\raisebox{-0.15ex}{\poplogo\fontsize{13}{13}\selectfont\color{popink}OnBuch\color{poporange}+}}
\rhead{\poppill{\DOCMATIERE\ \textbullet\ \DOCNIVEAU\ \textbullet\ Leçon \DOCLECON}{white}{popink}}
\lfoot{\popxbold\fontsize{7.6}{7.6}\selectfont\color{popmuted}\MakeUppercase{Cours signé OnBuch+}\ \textbullet\ {\popsemi\DOCTITRE}}
\rfoot{\tikz[baseline=-0.6ex]\node[rounded corners=8.5pt,fill=popink,minimum height=17pt,minimum width=17pt,inner xsep=5pt,inner sep=0]{\popxbold\fontsize{8}{8}\selectfont\color{popcream}\thepage\,/\,\pageref*{LastPage}};}

% ============================================================
% Titres
% ============================================================
\newcommand{\chaptitre}[2]{%
  \begin{tcolorbox}[enhanced,colback=poporange,colframe=popink,boxrule=2.6pt,arc=11pt,
      drop shadow=popink,left=15pt,right=15pt,top=12pt,bottom=11pt,before skip=3pt,after skip=18pt]
    \poppill{#2}{popink}{popcream}\par\vspace{9pt}
    {\raggedright\hyphenpenalty=10000\exhyphenpenalty=10000\popdisplay\fontsize{22}{24}\selectfont\color{popcream}\MakeUppercase{#1}\par}
  \end{tcolorbox}
}
% Section numérotée : pastille orange avec le numéro + titre Archivo
\newcounter{popsec}
\newcommand{\coursec}[1]{%
  \stepcounter{popsec}\setcounter{popsub}{0}%
  \par\vspace{16pt}\noindent
  \begin{minipage}{\linewidth}
  \tikz[baseline=-0.7ex]\node[draw=popink,line width=1.6pt,fill=poporange,rounded corners=4pt,minimum size=22pt,inner sep=0]{\popdisplay\fontsize{12}{12}\selectfont\color{popcream}\thepopsec};%
  \hspace{8pt}{\popdisplay\fontsize{15}{17}\selectfont\color{popink}\MakeUppercase{#1}}\par
  \vspace{4pt}\noindent{\color{poporange}\rule{36pt}{3.6pt}}
  \end{minipage}\par\nopagebreak\vspace{8pt}
}
\newcounter{popsub}
\newcommand{\courssub}[1]{%
  \stepcounter{popsub}%
  \par\vspace{9pt}\noindent{\popxbold\fontsize{11.5}{13}\selectfont\color{popink}{\color{poporange}\thepopsec.\thepopsub}\hspace{6pt}#1}\par\nopagebreak\vspace{4pt}
}

% ============================================================
% Boîtes pédagogiques — même recette : fond blanc, bordure épaisse colorée,
% ombre dure encre, badge plein. Titre optionnel [#1].
% ============================================================
\tcbset{popbase/.style={breakable,enhanced,colback=white,boxrule=2.3pt,arc=9pt,
  shadow={3pt}{-3pt}{0pt}{popink},left=14pt,right=14pt,top=11pt,bottom=11pt,before skip=11pt,after skip=15pt,
  fontupper=\popsemi\fontsize{10.6}{14.5}\selectfont\color{popink},
  fontlower=\popsemi\fontsize{10.6}{14.5}\selectfont\color{popink}}}
% Coche dessinée (le glyphe ✓ n'existe pas dans les polices du document)
\newcommand{\popcheck}[1]{\tikz[baseline=-0.5ex]\draw[#1,line width=1.6pt,line cap=round,line join=round](-0.13,0.0)--(-0.04,-0.09)--(0.13,0.1);}
\providecommand{\coche}{\popcheck{popgreen}}
\providecommand{\resultat}[1]{\par\smallskip\noindent\fcolorbox{popgreen}{popgreenL}{\parbox{\dimexpr\linewidth-2\fboxsep-2\fboxrule\relax}{\textbf{Résultat.} #1}}\par\smallskip}
\newcommand{\popboxhead}[3]{\poppill{#1}{#2}{white}\ifx\relax#3\relax\else\hspace{6pt}{\popxbold\fontsize{10.6}{12}\selectfont\color{popink}#3}\fi\par\vspace{7pt}}

\newtcolorbox{aretenir}[1][]{popbase,colframe=popgreen,before upper={\popboxhead{À retenir}{popgreen}{#1}}}
% Texte en chinois (document de lecture, dialogue, extrait) : cadre
% d'étude. Un titre optionnel (source, auteur) s'affiche dans l'en-tête.
\newtcolorbox{textebox}[1][]{popbase,colframe=popblue,before upper={\popboxhead{文本 · Texte}{popblue}{#1}},after upper={}}
% Marques ✓ / ✗ (forme correcte / fautive) souvent utilisées par les modèles
\providecommand{\cross}{\textcolor{poppink}{\ensuremath{\times}}}
\providecommand{\xmark}{\cross}\providecommand{\crossmark}{\cross}\providecommand{\cmark}{\ensuremath{\checkmark}}
% Ligne de réponse / justification à compléter (inventée par les modèles)
\providecommand{\justification}[1]{\par\noindent\textit{Justification~:} \dotfill\par}
% \textipa (package tipa non chargé) : la police ne couvre pas l'API -> simple passage
\providecommand{\textipa}[1]{#1}
% Soulignement (ulem non chargé) : \uline, \ul
\providecommand{\uline}[1]{\underline{#1}}\providecommand{\ul}[1]{\underline{#1}}
% \glossaire{\item ...} inventée par les modèles : liste de glossaire
\providecommand{\glossaire}[1]{\begin{itemize}#1\end{itemize}}
% \alert (beamer) inventée par les modèles : texte mis en évidence
\providecommand{\alert}[1]{\textbf{\textcolor{poppink}{#1}}}
% variantes ASCII d'ordinaux écrites en macro
\providecommand{\iere}{\textsuperscript{re}}\providecommand{\ieme}{\textsuperscript{e}}\providecommand{\ere}{\textsuperscript{re}}\providecommand{\eme}{\textsuperscript{e}}\providecommand{\er}{\textsuperscript{er}}\providecommand{\ie}{\textsuperscript{e}}
% Ordinaux français écrits en macro par les modèles : 1\ière, 3\ième
\newcommand{\ière}{\textsuperscript{re}}\newcommand{\ième}{\textsuperscript{e}}
% \exemple (inventée par les modèles) : étiquette « Exemple : »
\providecommand{\exemple}{\textbf{Exemple~:}\ }
% \square utilisable aussi hors mode math (cases à cocher)
\AtBeginDocument{\let\squareMath\square\renewcommand{\square}{\ensuremath{\squareMath}}}
% Mot ou phrase en chinois à l'intérieur d'un paragraphe français
\newcommand{\zh}[1]{\ifmmode\text{#1}\else{#1}\fi}
\let\eng\zh \let\esp\zh \let\all\zh % alias tolérés
% flèches d'intonation inventées par les modèles
\providecommand{\textsearrow}{}\renewcommand{\textsearrow}{\ensuremath{\searrow}}\providecommand{\textnearrow}{}\renewcommand{\textnearrow}{\ensuremath{\nearrow}}
% \fr{..} : mot français (inventé par les modèles)
\providecommand{\fr}[1]{\textit{#1}}
% zhbox : encadré de texte chinois inventé par les modèles
\newenvironment{zhbox}{\begin{quote}}{\end{quote}}
% \courssubsub : titre de niveau 3 inventé par les modèles
\providecommand{\courssubsub}[1]{\par\medskip\noindent\textbf{#1}\par\smallskip}
% \vs : « versus » inventé par les modèles
\providecommand{\vs}{\textit{vs}\ }
% \blank : case à compléter inventée par les modèles
\providecommand{\blank}[1][]{\underline{\hspace{1.6em}}}
\newtcolorbox{exemplebox}[1][]{popbase,colframe=poporange,before upper={\popboxhead{Exemple}{poporange}{#1}}}
\newtcolorbox{definition}[1][]{popbase,colframe=popblue,before upper={\popboxhead{Définition}{popblue}{#1}}}
\newtcolorbox{propriete}[1][]{popbase,colframe=poppurple,before upper={\popboxhead{Propriété}{poppurple}{#1}}}
\newtcolorbox{methode}[1][]{popbase,colframe=popgold,before upper={\popboxhead{Méthode}{popgold}{#1}}}
\newtcolorbox{attention}[1][]{popbase,colframe=poppink,before upper={\popboxhead{Attention}{poppink}{#1}}}
\newtcolorbox{experience}[1][]{popbase,colframe=popblue,colback=popblueL,before upper={\popboxhead{Expérience}{popblue}{#1}}}
% « Le savais-tu ? » : carte violette pleine (contexte, culture, Cameroun)
\newtcolorbox{savaistu}[1][]{popbase,colback=poppurple,colframe=popink,
  fontupper=\popsemi\fontsize{10.4}{14.2}\selectfont\color{white},
  before upper={\poppill{Le savais-tu\,?}{white}{poppurple}\ifx\relax#1\relax\else\hspace{6pt}{\popxbold\color{white}#1}\fi\par\vspace{7pt}}}
% Exercice résolu : énoncé en haut, \tcblower puis la solution sur fond crème
\newcounter{popexo}\newcounter{popexr}
\newtcolorbox{exoresolu}[1][]{popbase,colframe=poporange,colbacklower=popcream,
  segmentation style={popink,line width=1.2pt,dashed},
  before upper={\stepcounter{popexr}\popboxhead{Exercice résolu \thepopexr}{poporange}{#1}},
  before lower={\poppill{Solution}{popink}{popcream}\par\vspace{6pt}}}
% Exercice (non résolu en place) — la correction est dans la partie Corrigés
\newtcolorbox{exercice}[1][]{popbase,colframe=popink,
  before upper={\stepcounter{popexo}\popboxhead{Exercice \thepopexo}{popink}{#1}}}
% Corrigé
\newtcolorbox{corrige}[1][]{popbase,colframe=popgreen,colback=popgreenL,
  before upper={\popboxhead{Corrigé}{popgreen}{#1}}}
% Objectifs / prérequis / bilan
\newtcolorbox{objectifs}{popbase,colframe=popink,colback=poporangeL,
  before upper={\popboxhead{Objectifs de la leçon}{popink}{}}}
\newtcolorbox{prerequis}{popbase,colframe=popink,colback=popblueL,
  before upper={\popboxhead{Prérequis}{popblue}{}}}
\newtcolorbox{bilan}{popbase,colframe=popink,colback=white,boxrule=2.8pt,
  before upper={\popboxhead{Fiche bilan}{poporange}{L'essentiel à connaître pour l'examen}}}
% Figure encadrée avec légende
\newtcolorbox{popfigure}[1][]{enhanced,breakable=false,colback=white,colframe=popline,boxrule=1.4pt,arc=8pt,
  left=10pt,right=10pt,top=10pt,bottom=8pt,before skip=10pt,after skip=12pt,halign=center,
  lower separated=false,halign lower=center,
  fontlower=\popsemi\fontsize{9}{11.5}\selectfont\color{popmuted},
  #1}
\newcommand{\legende}[1]{\par\vspace{6pt}{\popsemi\fontsize{9}{11.5}\selectfont\color{popmuted}#1}}

% ============================================================
% Signature OnBuch+
% ============================================================
\newcommand{\poplogoinline}[1]{{\poplogo\fontsize{#1}{#1}\selectfont\color{popink}OnBuch\color{poporange}+}}
\newcommand{\signatureonbuch}{%
  \begin{tcolorbox}[enhanced,colback=popink,colframe=popink,boxrule=2.2pt,arc=10pt,
      drop shadow=poporange,left=14pt,right=14pt,top=10pt,bottom=10pt,before skip=0pt,after skip=0pt]
    \tikz[baseline=-0.7ex]\node[circle,fill=popgreen,draw=popcream,line width=1.3pt,minimum size=22pt,inner sep=0]{\popcheck{white}};%
    \hspace{9pt}\begin{minipage}[c]{0.62\linewidth}
      {\popxbold\fontsize{12}{13}\selectfont\color{popcream}Signé\ \textbullet\ {\poplogo OnBuch\color{poporange}+}}\par\vspace{2pt}
      {\raggedright\popsemi\fontsize{8}{10.5}\selectfont\color{popline}\MakeUppercase{Équipe pédagogique OnBuch+}\\\MakeUppercase{Conforme au programme officiel MINESEC}\par}
    \end{minipage}\hfill
    {\poplogo\fontsize{22}{22}\selectfont\color{popcream}OnBuch\color{poporange}+}
  \end{tcolorbox}%
}

% ============================================================
% Page de garde
% #1 titre · #2 matière · #3 niveau · #4 module · #5 n° leçon · #6 durée ·
% #7 description / objectifs (texte ou itemize)
% ============================================================
\newcommand{\pagegarde}[7]{%
  \thispagestyle{empty}
  \newgeometry{a4paper,margin=1.7cm,top=1.6cm,bottom=1.4cm}%
  \begin{tikzpicture}[remember picture,overlay]
    \fill[poporange!18] ([xshift=-3.2cm,yshift=-3.6cm]current page.north east) circle (4.6cm);
    \fill[poppurple!14] ([xshift=2.4cm,yshift=7.6cm]current page.south west) circle (3.8cm);
  \end{tikzpicture}%
  {\poplogo\fontsize{30}{30}\selectfont\color{popink}OnBuch\color{poporange}+}\hfill
  \raisebox{0.6ex}{\poppill{Cours signé \textbullet\ Premium}{popink}{popcream}}\par
  \vspace{1pt}\popsquiggle{poppurple}\par
  \vspace{8pt}
  \begin{tcolorbox}[enhanced,colback=poporange,colframe=popink,boxrule=2.8pt,arc=13pt,
      drop shadow=popink,left=18pt,right=18pt,top=12pt,bottom=13pt,after skip=10pt]
    \poppill{#2 \textbullet\ #3}{popink}{popcream}\hspace{5pt}\poppill{#4}{white}{popink}\par\vspace{8pt}
    {\popdisplay\fontsize{14}{14}\selectfont\color{popink}LEÇON #5}\par\vspace{4pt}
    {\raggedright\hyphenpenalty=10000\exhyphenpenalty=10000\popdisplay\fontsize{25}{27}\selectfont\color{popcream}\MakeUppercase{#1}\par}
    \vspace{8pt}
    {\popxbold\fontsize{9.5}{9.5}\selectfont\color{popink}\MakeUppercase{Durée conseillée : #6}}
  \end{tcolorbox}
  \begin{tcolorbox}[enhanced,colback=white,colframe=popink,boxrule=2.2pt,arc=10pt,
      drop shadow=popink,left=16pt,right=16pt,top=10pt,bottom=10pt,after skip=8pt]
    \poppill{Dans cette leçon}{poppurple}{white}\par\vspace{5pt}
    \popsemi\fontsize{9.3}{12.6}\selectfont\color{popink}#7
  \end{tcolorbox}
  \noindent\begin{minipage}[t]{0.487\linewidth}
    \begin{tcolorbox}[enhanced,colback=popgreen,colframe=popink,boxrule=2.2pt,arc=10pt,
        drop shadow=popink,left=11pt,right=11pt,top=10pt,bottom=10pt,height=3.1cm,valign=center]
      \tcbox[colback=white,colframe=popink,boxrule=1.3pt,arc=5pt,boxsep=0pt,nobeforeafter,
        left=3pt,right=3pt,top=3pt,bottom=3pt,valign=center]{\qrcode[height=1.9cm]{https://play.google.com/store/apps/details?id=com.luvvix.onbuch&hl=fr}}%
      \hspace{8pt}\begin{minipage}[c]{0.5\linewidth}
        {\popdisplay\fontsize{11}{12.5}\selectfont\color{white}\MakeUppercase{Télécharge OnBuch}}\par\vspace{4pt}
        {\raggedright\popsemi\fontsize{8.4}{11}\selectfont\color{white}Gratuit \textbullet\ Google Play\\Tuteur IA, annales, concours, résultats\par}
      \end{minipage}
    \end{tcolorbox}
  \end{minipage}\hfill
  \begin{minipage}[t]{0.487\linewidth}
    \begin{tcolorbox}[enhanced,colback=poppurple,colframe=popink,boxrule=2.2pt,arc=10pt,
        drop shadow=popink,left=11pt,right=11pt,top=10pt,bottom=10pt,height=3.1cm,valign=center]
      \tikz[baseline=-0.7ex]\node[circle,fill=white,draw=popink,line width=1.3pt,minimum size=1.6cm,inner sep=0]{\popdisplay\fontsize{24}{24}\selectfont\color{poppurple}+};%
      \hspace{8pt}\begin{minipage}[c]{0.6\linewidth}
        {\popdisplay\fontsize{11}{12.5}\selectfont\color{white}\MakeUppercase{Passe à OnBuch+}}\par\vspace{4pt}
        {\raggedright\popsemi\fontsize{8.4}{11}\selectfont\color{white}Tous les cours signés, Tuteur IA illimité, fiches PDF — onglet OnBuch+ de l'app\par}
      \end{minipage}
    \end{tcolorbox}
  \end{minipage}
  \vspace{8pt}
  \popcontenu
  \vfill
  \signatureonbuch
  \restoregeometry
  \newpage
}

% Rangée « Ce cours contient » (page de garde)
\newcommand{\poptile}[3]{% #1 couleur #2 gros texte #3 légende
  \begin{tcolorbox}[enhanced,colback=#1,colframe=popink,boxrule=2pt,arc=9pt,shadow={2.5pt}{-2.5pt}{0pt}{popink},
      left=6pt,right=6pt,top=9pt,bottom=9pt,halign=center,valign=center,height=1.75cm,width=\linewidth,nobeforeafter]
    {\popdisplay\fontsize{12.5}{13}\selectfont\color{white}\MakeUppercase{#2}}\par\vspace{3pt}
    {\centering\popsemi\fontsize{7.8}{9.5}\selectfont\color{white}#3\par}
  \end{tcolorbox}}
\newcommand{\popcontenu}{%
  \par\noindent{\popxboldls\fontsize{7.5}{7.5}\selectfont\color{popmuted}CE COURS SIGNÉ CONTIENT}\par\vspace{7pt}
  \noindent\begin{minipage}[t]{0.235\linewidth}\poptile{popblue}{Cours}{complet, illustré, pas à pas}\end{minipage}\hfill
  \begin{minipage}[t]{0.235\linewidth}\poptile{poporange}{Méthodes}{et exercices résolus}\end{minipage}\hfill
  \begin{minipage}[t]{0.235\linewidth}\poptile{poppink}{Exercices}{type examen + corrigés}\end{minipage}\hfill
  \begin{minipage}[t]{0.235\linewidth}\poptile{popgreen}{Fiche bilan}{l'essentiel à retenir}\end{minipage}\par}

% Sommaire
\newtcolorbox{sommaire}{enhanced,colback=white,colframe=popink,boxrule=2.2pt,arc=10pt,
  drop shadow=popink,left=18pt,right=18pt,top=13pt,bottom=13pt,before skip=6pt,after skip=20pt,
  before upper={\poppill{Sommaire}{poppurple}{white}\par\vspace{9pt}\popsemi\fontsize{10.6}{14.5}\selectfont\color{popink}}}

% Signature de fin de cours — poussée tout en bas de la dernière page
\newcommand{\findecours}{%
  \par\vfill\nopagebreak
  \begin{center}\popsquiggle{poporange}\end{center}\vspace{4pt}
  \signatureonbuch
}
\AtBeginDocument{\def\llbracket{\mathopen{[\mkern-3.5mu[}}\def\rrbracket{\mathclose{]\mkern-3.5mu]}}} % intervalles d'entiers (glyphe ⟦ absent de la police)
% Glyphes absents de Fira Math : redéfinis à partir de glyphes disponibles
\AtBeginDocument{%
  \def\setminus{\mathbin{\backslash}}\let\smallsetminus\setminus
  \def\vdots{\mathord{\vbox{\baselineskip3.2pt\lineskiplimit0pt\kern4pt\hbox{.}\hbox{.}\hbox{.}}}}%
  \def\ddots{\mathinner{\mkern1mu\raise6.5pt\hbox{.}\mkern2mu\raise3.6pt\hbox{.}\mkern2mu\raise0.7pt\hbox{.}\mkern1mu}}%
  \def\triangle{\mathord{\scalebox{0.9}{$\symup{\Delta}$}}}%
  \def\nabla{\mathord{\raisebox{\depth}{\scalebox{1}[-1]{$\symup{\Delta}$}}}}%
  \def\checkmark{\popcheck{popdark}}%
  \def\star{\mathbin{\ast}}\def\bigstar{\mathord{\ast}}%
  \def\diamond{\mathbin{\scalebox{0.6}{$\Diamond$}}}%
  \def\bigcup{\mathop{\mathchoice{\raisebox{-0.3em}{\scalebox{1.7}{$\cup$}}}{\scalebox{1.25}{$\cup$}}{\cup}{\cup}}\displaylimits}%
  \def\bigcap{\mathop{\mathchoice{\raisebox{-0.3em}{\scalebox{1.7}{$\cap$}}}{\scalebox{1.25}{$\cap$}}{\cap}{\cap}}\displaylimits}%
  \def\lesseqqgtr{\mathrel{\lessgtr}}\def\bigoplus{\mathop{\oplus}\displaylimits}%
}
\providecommand{\ssi}{si et seulement si} % abréviation fréquente
\AtBeginDocument{\providecommand{\wideparen}{\overparen}} % arc de cercle

%% FIGFIX-BEGIN v3 — correctifs globaux des figures (docs/onbuch-generation/tools/figfix)
\usepackage{adjustbox}\usepackage{etoolbox}
% toute figure TikZ/pgfplots plus large que la ligne est réduite à la largeur disponible
\BeforeBeginEnvironment{tikzpicture}{\begin{adjustbox}{max width=\linewidth}}
\AfterEndEnvironment{tikzpicture}{\end{adjustbox}}
% légendes pgfplots lisibles par-dessus les courbes
\pgfplotsset{every axis/.append style={legend style={fill=white,fill opacity=0.92,text opacity=1,draw=black!15,inner sep=3pt}}}
% étiquettes d'arêtes (syntaxe edge["..."]) avec fond blanc
\tikzset{every edge quotes/.append style={fill=white,inner sep=1.2pt}}
% bulles de cartes mentales : texte à la ligne dans la bulle, bulle un peu plus grande
\tikzset{every concept/.append style={text width=2.0cm,align=center,minimum size=2.3cm,inner sep=2pt}}
%% FIGFIX-END
\begin{document}

\pagegarde{Leçon 24 — Expression orale : démocratie}{Chinois}{Tle A}{Module 3 : Citoyenneté et ouverture au monde}{ZH24}{2 h}{Cette leçon d'expression orale permet aux élèves de Terminale A de présenter un exposé simple et de participer à un débat guidé en chinois sur le thème de la démocratie. Elle fournit les formules fonctionnelles, les dialogues modèles et le travail des tons nécessaires à une prise de parole claire et correcte.
\vspace{4pt}
\begin{multicols}{2}\small
\begin{itemize}
\item À la fin de cette leçon, je sais utiliser le vocabulaire de base lié à la démocratie (民主, 选举, 投票, 意见, 自由...)
\item À la fin de cette leçon, je sais présenter oralement un exposé simple et structuré en chinois
\item À la fin de cette leçon, je sais donner mon avis avec les formules 我觉得 / 我认为 / 在我看来
\item À la fin de cette leçon, je sais exprimer l'accord et le désaccord poliment (我同意 / 我不同意 / 你说得对)
\item À la fin de cette leçon, je sais relancer un débat et poser des questions à mes camarades (你觉得呢？ 为什么？)
\item À la fin de cette leçon, je sais prononcer correctement les tons des mots clés du thème
\end{itemize}\end{multicols}}

\begin{sommaire}
\begin{multicols}{2}
\begin{enumerate}
\item Mise en route et vocabulaire du thème
\item Dialogue modèle 1 : donner son avis
\item Dialogue modèle 2 : accord, désaccord et relance
\item Structurer un exposé simple
\item Travail spécifique des tons
\item Jeux de rôle et débat guidé
\item Activité d'intégration
\item Méthodes et exercices résolus
\item Exercices
\item Corrigés des exercices
\item Fiche bilan
\end{enumerate}
\end{multicols}
\end{sommaire}

\begin{objectifs}
\begin{itemize}
\item À la fin de cette leçon, je sais utiliser le vocabulaire de base lié à la démocratie (民主, 选举, 投票, 意见, 自由...)
\item À la fin de cette leçon, je sais présenter oralement un exposé simple et structuré en chinois
\item À la fin de cette leçon, je sais donner mon avis avec les formules 我觉得 / 我认为 / 在我看来
\item À la fin de cette leçon, je sais exprimer l'accord et le désaccord poliment (我同意 / 我不同意 / 你说得对)
\item À la fin de cette leçon, je sais relancer un débat et poser des questions à mes camarades (你觉得呢？ 为什么？)
\item À la fin de cette leçon, je sais prononcer correctement les tons des mots clés du thème
\item À la fin de cette leçon, je sais conclure un exposé avec une formule de synthèse (总之, 所以)
\item À la fin de cette leçon, je sais comparer brièvement des situations citoyennes du Cameroun et de la Chine à l'oral
\end{itemize}
\end{objectifs}

\begin{prerequis}
\begin{itemize}
\item Structures de phrase de base sujet-verbe-objet (Première)
\item Pronoms personnels et verbes modaux 想, 会, 可以, 应该
\item Formules de présentation personnelle et d'opinion simple vues en Première (我喜欢, 我觉得)
\item Maîtrise des quatre tons et du pinyin
\item Vocabulaire de la vie scolaire et de la citoyenneté (学校, 国家, 人民)
\end{itemize}
\end{prerequis}

\newpage

\coursec{Mise en route et vocabulaire du thème}

\courssub{Situation d'accroche : un vote contesté en classe de Terminale}

Imaginez la scène. Lundi matin, lycée de Yaoundé, classe de Terminale A. Le professeur principal vient d'organiser l'élection des délégués. Les élèves votent à main levée : « Qui est pour Amina ? Qui est pour Junior ? » Les mains se lèvent, on compte, on annonce le résultat. Amina est élue déléguée avec vingt-trois voix contre dix-neuf. Mais aussitôt, les discussions s'animent : certains élèves contestent le résultat (« le comptage était trop rapide ! »), d'autres défendent le choix de la majorité (« la majorité a parlé, il faut respecter le vote »). Chacun veut donner son avis, s'exprimer, débattre.

Cette petite scène de la vie scolaire camerounaise est en réalité une leçon de vie citoyenne. Voter, donner son opinion, accepter le choix de la majorité, débattre dans le respect : voilà des gestes que l'on retrouve au Cameroun lors des élections, en Chine dans les assemblées locales, et partout où les citoyens participent à la vie collective.

\textbf{Problématique de la leçon :} comment dire en chinois « je suis d'accord », « je ne suis pas d'accord », « à mon avis » ? Comment présenter un petit exposé et débattre sur le thème de la démocratie ? Pour y arriver, il faut d'abord posséder le vocabulaire du thème. C'est l'objectif de cette première étape : découvrir, comprendre, prononcer et mémoriser les huit mots clés qui nous permettront de parler de démocratie en chinois.

\courssub{Brainstorming : qu'est-ce que la démocratie ?}

Avant d'ouvrir le dictionnaire chinois, réfléchissons en français. Posez-vous les questions suivantes, en groupe ou en classe entière :

\begin{itemize}
\item Qu'est-ce que la démocratie ? Qui décide, et comment ?
\item Avez-vous déjà voté ? (élection des délégués de classe, du chef de salle, du capitaine de l'équipe de football\ldots)
\item Quels droits possède un citoyen camerounais ?
\item Pourquoi est-il important de pouvoir donner son opinion librement ?
\end{itemize}

En classe, le professeur note au tableau toutes les idées proposées : \emph{élections}, \emph{vote}, \emph{liberté d'expression}, \emph{droits}, \emph{citoyens}, \emph{majorité}, \emph{respect des opinions}. Remarquez que ces mots français correspondent presque exactement aux huit mots chinois que nous allons apprendre. Le thème est universel ; seule la langue change.

\begin{experience}[Brainstorming en classe]
Par groupes de quatre, listez en deux minutes le maximum de mots français liés au vote et à la vie citoyenne. Le groupe qui trouve le plus de mots gagne. Ensuite, chaque groupe choisit ses trois mots les plus importants et les compare avec le tableau de vocabulaire chinois ci-dessous : quels mots chinois correspondent à vos choix ?
\end{experience}

\courssub{Première observation : un texte d'ancrage}

Lisons d'abord ce court texte. Ne cherchez pas encore à tout comprendre mot par mot : observez, repérez les mots qui reviennent, écoutez le professeur le lire à voix haute.

\begin{textebox}[Texte d'ancrage : la démocratie et le vote]
民主很重要。每个公民都有投票的权利。\\
Mínzhǔ hěn zhòngyào. Měi ge gōngmín dōu yǒu tóupiào de quánlì.\\
« La démocratie est importante. Chaque citoyen a le droit de voter. »\\[4pt]
在选举中，公民可以自由地表达自己的意见。\\
Zài xuǎnjǔ zhōng, gōngmín kěyǐ zìyóu de biǎodá zìjǐ de yìjiàn.\\
« Lors des élections, les citoyens peuvent exprimer librement leur opinion. »\\[4pt]
每个国家都有自己的选举制度。\\
Měi ge guójiā dōu yǒu zìjǐ de xuǎnjǔ zhìdù.\\
« Chaque pays a son propre système électoral. »
\end{textebox}

Observez bien : les mots \zh{民主}, \zh{公民}, \zh{投票}, \zh{权利}, \zh{选举}, \zh{自由}, \zh{意见}, \zh{国家} apparaissent tous dans ce petit texte. Ce sont précisément nos huit mots clés. Remarquez aussi la structure très régulière des phrases chinoises : Sujet + verbe + complément, sans conjugaison ni article. Le verbe \zh{有} (yǒu, avoir) exprime la possession d'un droit, exactement comme « avoir le droit de » en français. Notez enfin la particule \zh{地} (de) dans \zh{自由地表达} : placée après un adjectif ou un nom utilisé comme adverbe, elle introduit la manière dont on fait l'action, comme le suffixe français « -ment » (libre $\rightarrow$ libre\emph{ment}).

\courssub{Le mot central : 民主 mínzhǔ}

\begin{definition}[民主 mínzhǔ : la démocratie]
Le mot \zh{民主} (mínzhǔ) se compose de deux caractères : \zh{民} (mín), qui signifie « le peuple », et \zh{主} (zhǔ), qui signifie « maître, diriger, être maître de ». Littéralement : \cle{« le peuple est maître »}, c'est-à-dire la démocratie. Cette étymologie transparente est une chance pour vous : en retenant les deux caractères séparément, vous retenez le mot entier et vous comprenez sa logique interne.
\end{definition}

Cette construction par composition est typique du chinois moderne : la plupart des mots sont formés de deux caractères dont le sens se combine. Comparez avec le français : « démocratie » vient du grec \emph{dêmos} (peuple) et \emph{kratos} (pouvoir). Le chinois et le français ont donc construit le même mot avec la même idée : le pouvoir appartient au peuple.

\begin{savaistu}[Le caractère 民, une brique très productive]
Le caractère \zh{民} (mín, peuple) apparaît dans de nombreux mots de la vie civique : \zh{人民} (rénmín, le peuple), \zh{民主} (mínzhǔ, la démocratie), \zh{公民} (gōngmín, le citoyen, littéralement « peuple public »). Apprendre un caractère, c'est donc ouvrir la porte à toute une famille de mots. De même, \zh{国} (guó, pays) se retrouve dans \zh{国家} (guójiā, le pays, l'État), \zh{中国} (Zhōngguó, la Chine), \zh{外国} (wàiguó, l'étranger). Observez les familles de caractères : c'est le secret d'une mémorisation rapide.
\end{savaistu}

\courssub{Tableau de vocabulaire du thème}

Voici les huit mots clés de la leçon. Lisez le pinyin à voix haute, en marquant bien les tons.

\begin{center}
\begin{tabularx}{\linewidth}{|X|X|X|X|}
\hline
\rowcolor{popblueL} Caractères & Pinyin & Nature & Traduction \\
\hline
民主 & mínzhǔ & nom & la démocratie \\
\hline
选举 & xuǎnjǔ & nom / verbe & l'élection ; élire \\
\hline
投票 & tóupiào & verbe & voter, déposer un bulletin \\
\hline
意见 & yìjiàn & nom & l'opinion, l'avis \\
\hline
自由 & zìyóu & nom / adjectif & la liberté ; libre \\
\hline
权利 & quánlì & nom & le droit (d'un citoyen) \\
\hline
公民 & gōngmín & nom & le citoyen \\
\hline
国家 & guójiā & nom & le pays, l'État \\
\hline
\end{tabularx}
\end{center}

Quelques remarques de précision pour bien utiliser ces mots :

\begin{itemize}
\item \zh{意见} (yìjiàn) désigne l'opinion, l'avis que l'on exprime. Pour dire « à mon avis », on dira \zh{我的意见是} (wǒ de yìjiàn shì) ou, plus courant à l'oral, \zh{我觉得} (wǒ juéde, je pense que).
\item \zh{权利} (quánlì) est le droit reconnu au citoyen (droit de vote, droit d'expression). Ne le confondez pas avec \zh{权力} (quánlì, le pouvoir), qui se prononce exactement pareil mais s'écrit différemment : un vrai piège d'homophone !
\item \zh{投票} (tóupiào) : \zh{投} signifie « jeter, lancer » et \zh{票} « bulletin, billet ». Voter, c'est littéralement « lancer son bulletin ». Belle image, facile à retenir.
\item \zh{自由} (zìyóu) : \zh{自} « soi-même » + \zh{由} « provenir de, suivre » : « venir de soi-même », donc être libre.
\item \zh{选举} (xuǎnjǔ) peut être un nom (\zh{在选举中}, lors des élections) ou un verbe (\zh{我们选举代表}, nous élisons nos représentants).
\end{itemize}

\begin{popfigure}
\begin{tikzpicture}[mindmap, grow cyclic, every node/.style={concept}, concept color=popblue!60, text=popdark, level 1/.append style={level distance=3.2cm, sibling angle=60}, level 1 concept/.append style={concept color=poporange!50, font=\small}]
\node[align=center]{民主\\ mínzhǔ\\ démocratie}
  child { node[align=center]{选举\\ xuǎnjǔ\\ élection} }
  child { node[align=center]{投票\\ tóupiào\\ voter} }
  child { node[align=center]{意见\\ yìjiàn\\ opinion} }
  child { node[align=center]{自由\\ zìyóu\\ liberté} }
  child { node[align=center]{权利\\ quánlì\\ droit} }
  child { node[align=center]{公民\\ gōngmín\\ citoyen} };
\end{tikzpicture}
\legende{Carte mentale du lexique : autour du mot central 民主 (mínzhǔ), six mots satellites à mémoriser ensemble ; le huitième mot, 国家 (guójiā), sera ajouté par vos soins lors de la recopie de mémoire.}
\end{popfigure}

\courssub{Prononciation guidée : travailler les tons}

Le vocabulaire ne sert à rien s'il est mal prononcé : en chinois, un ton faux change le sens du mot. Procédons en trois étapes.

\textbf{Étape 1 : répétition chorale.} Le professeur lit chaque mot deux fois, lentement, puis la classe répète en chœur trois fois. Le professeur exagère volontairement les tons pour les rendre perceptibles.

\textbf{Étape 2 : analyse des tons difficiles.}

\begin{itemize}
\item \zh{民主} mínzhǔ : 2\textsuperscript{e} ton (qui monte, comme quand on demande « hein ? ») puis 3\textsuperscript{e} ton (qui descend puis remonte). Enchaîner une montée puis une descente-remontée demande de l'entraînement : dites \emph{mín} en montant, puis \emph{zhǔ} en creusant la voix.
\item \zh{投票} tóupiào : 2\textsuperscript{e} ton puis 4\textsuperscript{e} ton (qui tombe brutalement, comme un ordre : « non ! »). Montez sur \emph{tóu}, laissez tomber la voix sur \emph{piào}.
\item \zh{选举} xuǎnjǔ : deux 3\textsuperscript{e} tons de suite. Règle de sandhi : le premier 3\textsuperscript{e} ton devient un 2\textsuperscript{e} ton à la prononciation. On prononce donc presque \emph{xuánjǔ}, mais on écrit toujours xuǎnjǔ en pinyin.
\item \zh{自由} zìyóu : 4\textsuperscript{e} ton puis 2\textsuperscript{e} ton : tombez sur \emph{zì}, remontez sur \emph{yóu}.
\item \zh{公民} gōngmín : deux 1\textsuperscript{ers} tons, hauts et plats, comme une note de musique tenue. Ne laissez pas la voix descendre entre les deux syllabes.
\end{itemize}

\textbf{Étape 3 : vérification individuelle.} Le professeur interroge quelques élèves isolément et corrige les tons.

\begin{attention}[Erreur fréquente : 意见 yìjiàn]
Les francophones prononcent souvent \zh{意见} avec deux syllabes plates et égales, sans aucune chute de voix : \emph{yi-jian}. C'est une erreur. Les deux caractères portent le 4\textsuperscript{e} ton : la voix doit \cle{chuter franchement} sur \zh{意} (yì), comme un ordre bref, puis repartir d'en haut et \cle{chuter à nouveau} sur \zh{见} (jiàn), avec un peu moins de force sur la deuxième syllabe. Pensez à deux marches d'escalier que l'on descend d'un pas sec : haut $\rightarrow$ bas, puis haut $\rightarrow$ bas, plus léger. Écoutez le professeur, répétez en ralentissant : \emph{yì\ldots jiàn}.
\end{attention}

\courssub{Mini-phrases d'ancrage}

Un mot ne se mémorise vraiment que dans une phrase. Voici des mini-phrases courtes, à répéter et à recopier. Elles réutilisent toutes le vocabulaire du tableau.

\begin{exemplebox}[Mini-phrases d'ancrage]
\zh{公民有投票的权利。}\\
\emph{Gōngmín yǒu tóupiào de quánlì.}\\
« Les citoyens ont le droit de voter. »\\[4pt]
\zh{民主很重要。}\\
\emph{Mínzhǔ hěn zhòngyào.}\\
« La démocratie est importante. »\\[4pt]
\zh{每个人都可以表达自己的意见。}\\
\emph{Měi ge rén dōu kěyǐ biǎodá zìjǐ de yìjiàn.}\\
« Chacun peut exprimer sa propre opinion. »\\[4pt]
\zh{我们国家有选举。}\\
\emph{Wǒmen guójiā yǒu xuǎnjǔ.}\\
« Notre pays organise des élections. »\\[4pt]
\zh{言论自由是公民的权利。}\\
\emph{Yánlùn zìyóu shì gōngmín de quánlì.}\\
« La liberté d'expression est un droit du citoyen. »
\end{exemplebox}

Observez la structure de la première phrase : \zh{公民} (sujet) + \zh{有} (verbe « avoir ») + \zh{投票的权利} (complément : « le droit de voter »). La particule \zh{的} (de) relie le verbe \zh{投票} au nom \zh{权利} : elle fonctionne comme le « de » français dans « le droit \emph{de} voter ». Cette structure \cle{verbe + 的 + nom} est extrêmement productive : \zh{说话的自由} (la liberté de parler), \zh{学习的权利} (le droit d'étudier).

\begin{propriete}[La structure « Sujet + 有 + verbe + 的 + nom »]
Pour exprimer « avoir le droit / la liberté / la possibilité de faire quelque chose », le chinois utilise le schéma : \cle{Sujet + 有 + verbe + 的 + nom}. Le verbe se place \emph{avant} \zh{的}, et le nom (\zh{权利}, \zh{自由}) après. Exemple : \zh{每个公民都有表达意见的自由。} \emph{Měi ge gōngmín dōu yǒu biǎodá yìjiàn de zìyóu.} « Chaque citoyen a la liberté d'exprimer son opinion. » Attention à l'ordre : on ne dit jamais \zh{有权利的投票} ; le verbe précède toujours \zh{的}.
\end{propriete}

\courssub{Mémorisation active : association mot-image et jeu de devinettes}

\begin{methode}[Mémoriser le lexique en trois techniques]
\begin{enumerate}
\item \textbf{Association mot-image :} pour chaque mot, créez une image mentale. \zh{投票} : une main qui jette un bulletin dans une urne. \zh{自由} : un oiseau qui s'envole. \zh{选举} : des mains levées dans la classe de Terminale A. Plus l'image est personnelle et amusante, plus elle est efficace.
\item \textbf{Carte mentale :} recopiez la carte mentale de la leçon dans votre cahier, de mémoire, puis vérifiez. Le fait de redessiner les branches active la mémoire visuelle.
\item \textbf{Jeu de devinettes :} un élève mime ou définit un mot en français (« c'est ce que font les citoyens le jour du scrutin »), les autres doivent répondre en chinois avec le bon ton : \zh{投票}！
\end{enumerate}
\end{methode}

\begin{experience}[Jeu de devinettes en binômes]
Par deux : l'élève A tire au sort un mot du tableau et le définit en français ou le mime, sans jamais le dire. L'élève B doit répondre en chinois, en prononçant correctement les tons, puis construire une phrase complète avec ce mot. Exemple : A mime un bulletin glissé dans une urne ; B répond \zh{投票！公民有投票的权利。} Chaque binôme joue quatre manches, puis on inverse les rôles. Le professeur circule et corrige les tons.
\end{experience}

\courssub{Exercice résolu et entraînement}

\begin{exoresolu}[Compléter avec le mot qui convient]
Complétez les phrases avec : \zh{民主}, \zh{选举}, \zh{投票}, \zh{意见}, \zh{权利}.\\
1. 每个公民都有\_\_\_\_\_\_的权利。(Chaque citoyen a le droit de voter.)\\
2. \_\_\_\_\_\_很重要，每个国家都需要。(La démocratie est importante, chaque pays en a besoin.)\\
3. 在\_\_\_\_\_\_中，我们选我们的代表。(Lors des élections, nous choisissons nos représentants.)\\
4. 请你说说你的\_\_\_\_\_\_。(Donne ton avis, s'il te plaît.)
\tcblower
\textbf{Solution détaillée :}\\
1. \zh{投票} : la structure « avoir le droit de + verbe » exige un verbe après \zh{有} ; seul \zh{投票} (voter) convient : \zh{每个公民都有投票的权利。} \emph{Měi ge gōngmín dōu yǒu tóupiào de quánlì.}\\
2. \zh{民主} : le sujet de « est important » est un concept, la démocratie : \zh{民主很重要，每个国家都需要。} \emph{Mínzhǔ hěn zhòngyào, měi ge guójiā dōu xūyào.}\\
3. \zh{选举} : après la préposition \zh{在} (dans, lors de), il faut un nom d'événement : \zh{在选举中，我们选我们的代表。} \emph{Zài xuǎnjǔ zhōng, wǒmen xuǎn wǒmen de dàibiǎo.}\\
4. \zh{意见} : « ton avis » se dit \zh{你的意见} ; le verbe \zh{说说} (dire un peu, parler de) appelle un nom de contenu : \zh{请你说说你的意见。} \emph{Qǐng nǐ shuōshuo nǐ de yìjiàn.}
\end{exoresolu}

\begin{exercice}[À vous de jouer]
1. Traduisez en chinois : « Les citoyens ont le droit de voter. » ; « À mon avis, la liberté est importante. »\\
2. Recopiez de mémoire la carte mentale du lexique en y ajoutant le mot \zh{国家}, puis vérifiez avec le manuel.\\
3. Prononcez à voix haute cinq fois : \zh{民主}, \zh{投票}, \zh{意见}, en marquant les tons. Faites-vous corriger par votre binôme.\\
4. Trouvez deux autres mots chinois contenant le caractère \zh{民} et donnez leur traduction.
\end{exercice}

\begin{corrige}[Exercice « À vous de jouer »]
1. \zh{公民有投票的权利。} \emph{Gōngmín yǒu tóupiào de quánlì.} ; \zh{我觉得自由很重要。} \emph{Wǒ juéde zìyóu hěn zhòngyào.} (on accepte aussi \zh{我的意见是，自由很重要。})\\
2. Vérification avec la carte mentale de la leçon : sept branches au total, dont \zh{国家} (guójiā, le pays).\\
3. Auto-correction en binôme ; le professeur vérifie les tons : mínzhǔ (2+3), tóupiào (2+4), yìjiàn (4+4).\\
4. Par exemple : \zh{人民} (rénmín, le peuple) et \zh{农民} (nóngmín, le paysan) ; on accepte aussi \zh{居民} (jūmín, l'habitant, le résident).
\end{corrige}

\begin{aretenir}[Ce qu'il faut retenir de cette étape]
\begin{itemize}
\item Huit mots clés pour parler de démocratie : \zh{民主}, \zh{选举}, \zh{投票}, \zh{意见}, \zh{自由}, \zh{权利}, \zh{公民}, \zh{国家}.
\item \zh{民主} = \zh{民} (peuple) + \zh{主} (maître) : « le peuple est maître ».
\item Structure d'ancrage : \cle{Sujet + 有 + verbe + 的 + nom} : \zh{公民有投票的权利。}
\item La particule \zh{地} (de) après un adjectif forme un adverbe de manière : \zh{自由地表达} (exprimer librement).
\item Les tons font partie du mot : mínzhǔ (2+3), tóupiào (2+4), yìjiàn (4+4, deux chutes franches).
\item Un caractère appris en ouvre d'autres : \zh{民} donne \zh{人民}, \zh{民主}, \zh{公民}.
\end{itemize}
\end{aretenir}

Nous disposons maintenant du vocabulaire et des premières phrases. Dans la prochaine étape, nous entrerons dans un premier dialogue modèle pour apprendre à donner notre avis en chinois.

\coursec{Dialogue modèle 1 : donner son avis}

Dans la section précédente, nous avons découvert le vocabulaire du thème de la démocratie : \zh{民主} (mínzhǔ, « démocratie »), \zh{选举} (xuǎnjǔ, « élection »), \zh{意见} (yìjiàn, « opinion »), \zh{权利} (quánlì, « droit »). Ces mots ne servent à rien si l'on ne sait pas les placer dans une conversation. C'est exactement l'objectif de cette section : apprendre, à partir d'un dialogue modèle, à \cle{donner son avis} en chinois, à le \cle{justifier}, et à prononcer correctement les tons des mots clés. Nous travaillerons dans une situation bien connue des élèves camerounais : l'élection des délégués de classe.

\courssub{Écoute et lecture du dialogue modèle}

\courssubsub*{}
\textbf{Situation.}\par
Deux élèves de Terminale, Awa et Lin, discutent après l'élection des délégués de classe. Awa pose une question d'opinion ; Lin répond en donnant son avis et en le justifiant. Écoutez d'abord le professeur lire le dialogue deux fois, puis lisez-le à voix haute en binôme.

\begin{textebox}[Dialogue modèle : l'élection des délégués]
A : 你觉得民主重要吗？\\
\emph{A : Nǐ juéde mínzhǔ zhòngyào ma ?}\\
« A : Penses-tu que la démocratie est importante ? »\\[0.4em]
B : 我觉得很重要，因为每个人都可以说自己的意见。\\
\emph{B : Wǒ juéde hěn zhòngyào, yīnwèi měi ge rén dōu kěyǐ shuō zìjǐ de yìjiàn.}\\
« B : Je pense que c'est très important, parce que chacun peut exprimer sa propre opinion. »
\end{textebox}

\textbf{Observation guidée.} Avant toute règle, observons ce dialogue et répondons à trois questions :
\begin{enumerate}
\item Quel mot introduit la question d'opinion ? (\zh{觉得}, juéde, « penser, trouver ».)
\item Où se place le point d'interrogation en chinois ? (La particule \zh{吗} en fin de phrase, puis le point d'interrogation pleine chasse \zh{？}.)
\item Quel mot introduit la justification ? (\zh{因为}, yīnwèi, « parce que ».)
\end{enumerate}

\begin{center}
\begin{tabularx}{\linewidth}{|X|X|X|X|}
\hline
\rowcolor{popblueL} Caractères & Pinyin & Nature & Traduction \\
\hline
觉得 & juéde & verbe & penser, trouver, avoir l'impression que \\
\hline
民主 & mínzhǔ & nom & démocratie \\
\hline
重要 & zhòngyào & adjectif & important \\
\hline
因为 & yīnwèi & conjonction & parce que \\
\hline
每 & měi & déterminant & chaque \\
\hline
可以 & kěyǐ & verbe modal & pouvoir (avoir la permission / la possibilité) \\
\hline
说 & shuō & verbe & dire, exprimer \\
\hline
自己 & zìjǐ & pronom & soi-même, son propre \\
\hline
意见 & yìjiàn & nom & opinion, avis \\
\hline
\end{tabularx}
\end{center}

\courssub{Les formules d'opinion : 觉得, 认为, 在我看来}

\textbf{Observation.} Comparons trois phrases qui expriment la même idée avec des nuances différentes :

\begin{exemplebox}[Trois façons de dire « je pense que »]
\zh{我觉得民主很重要。} \emph{(Wǒ juéde mínzhǔ hěn zhòngyào.)} « Je pense que la démocratie est importante. » — registre courant, conversation.\\[0.3em]
\zh{我认为自由很重要。} \emph{(Wǒ rènwéi zìyóu hěn zhòngyào.)} « Je pense que la liberté est importante. » — registre plus formel, exposé, débat.\\[0.3em]
\zh{在我看来，投票是公民的权利。} \emph{(Zài wǒ kànlái, tóupiào shì gōngmín de quánlì.)} « À mon avis, voter est un droit du citoyen. » — introduit un point de vue personnel, souvent en tête de phrase.
\end{exemplebox}

\begin{propriete}[La structure d'opinion : 我觉得 + proposition]
La structure de base pour donner son avis est : \textbf{Sujet + 觉得 / 认为 + proposition complète (sujet + verbe)}.
\begin{itemize}
\item \zh{觉得} (juéde) exprime une opinion personnelle, souvent fondée sur un ressenti : « je trouve que, je pense que ». C'est la formule la plus fréquente à l'oral.
\item \zh{认为} (rènwéi) exprime une opinion réfléchie, argumentée ; il est \textbf{plus formel} et convient parfaitement à l'exposé ou au débat.
\item \zh{在我看来} (zài wǒ kànlái, « à mon avis ») se place en tête de phrase, suivi d'une virgule chinoise \zh{，}.
\end{itemize}
\textbf{Point capital :} en chinois, il n'y a \textbf{aucun mot} qui correspond au « que » français. On enchaîne directement : \zh{我觉得} + phrase complète. La proposition après \zh{觉得} garde l'ordre normal sujet + verbe.
\end{propriete}

\begin{attention}[Ne traduis pas mot à mot « je suis d'avis que »]
Les francophones veulent souvent traduire « je suis d'avis que » ou chercher un équivalent du « que ». En chinois, c'est impossible et inutile : on dit simplement \zh{我觉得} + phrase complète, \textbf{sans « que »}.\\
Erreur typique : dire \zh{我觉得的民主很重要} en insérant un \zh{的} parasite — faute grave. Autre erreur : oublier le sujet dans la proposition subordonnée. On ne dit pas \zh{我觉得很重要} pour « je pense que la démocratie est importante » si l'on veut être précis ; il faut dire \zh{我觉得民主很重要} (sauf si le sujet est évident dans le contexte, comme dans le dialogue modèle où \zh{民主} vient d'être mentionné).
\end{attention}

\courssub{Justifier son avis avec 因为...所以...}

\textbf{Observation.} Dans le dialogue, Lin ne se contente pas de dire « c'est important » : il donne une raison avec \zh{因为}. Comparons :

\begin{exemplebox}[Donner une raison]
\zh{因为每个人都可以说自己的意见。} \emph{(Yīnwèi měi ge rén dōu kěyǐ shuō zìjǐ de yìjiàn.)} « Parce que chacun peut exprimer sa propre opinion. »\\[0.3em]
\zh{因为民主很重要，所以我们要参加选举。} \emph{(Yīnwèi mínzhǔ hěn zhòngyào, suǒyǐ wǒmen yào cānjiā xuǎnjǔ.)} « Parce que la démocratie est importante, nous devons donc participer aux élections. »
\end{exemplebox}

\begin{propriete}[La corrélation 因为...所以...]
Structure : \textbf{因为 + cause，(所以) + conséquence}. « Parce que..., (donc)... ».
\begin{itemize}
\item À l'oral et dans une réponse courte, on peut utiliser \zh{因为} seul, comme dans le dialogue modèle.
\item Dans une phrase complète, le chinois emploie souvent \textbf{les deux} : \zh{因为...所以...}. C'est une différence majeure avec le français, où « parce que... donc... » ensemble est une faute. En chinois, la corrélation double est \textbf{correcte et naturelle}.
\item On peut omettre \zh{因为} et garder seulement \zh{所以} : \zh{民主很重要，所以我们要参加。} « La démocratie est importante, donc nous devons participer. »
\item En revanche, on ne peut pas utiliser \zh{所以} \emph{avant} la cause.
\end{itemize}
\end{propriete}

\begin{center}
\begin{tabularx}{\linewidth}{|X|X|X|}
\hline
\rowcolor{popblueL} Structure & Exemple (caractères + pinyin) & Traduction \\
\hline
因为 seul (réponse courte) & 因为每个人都可以说自己的意见。 Yīnwèi měi ge rén dōu kěyǐ shuō zìjǐ de yìjiàn. & Parce que chacun peut exprimer son opinion. \\
\hline
因为...所以... (phrase complète) & 因为自由很重要，所以我们要保护它。 Yīnwèi zìyóu hěn zhòngyào, suǒyǐ wǒmen yào bǎohù tā. & Parce que la liberté est importante, nous devons la protéger. \\
\hline
所以 seul & 每个人都有意见，所以民主很重要。 Měi ge rén dōu yǒu yìjiàn, suǒyǐ mínzhǔ hěn zhòngyào. & Chacun a une opinion, donc la démocratie est importante. \\
\hline
\end{tabularx}
\end{center}

\begin{popfigure}
\begin{tikzpicture}[node distance=0.5cm and 0.7cm,
box/.style={draw=popblue, thick, rounded corners, fill=popblueL, align=center, minimum height=1cm, inner sep=6pt},
cause/.style={draw=poporange, thick, rounded corners, fill=poporangeL, align=center, minimum height=1cm, inner sep=6pt},
arr/.style={-{Stealth[length=3mm]}, thick, popdark}]
\node[box, align=center] (avis){Formule d'avis\\我觉得 / 我认为};
\node[box, right=of avis, align=center] (op){Opinion\\民主很重要};
\node[cause, right=of op] (yin) {因为};
\node[cause, right=of yin, align=center] (just){Justification\\每个人都可以\\说自己的意见};
\draw[arr] (avis) -- (op);
\draw[arr] (op) -- (yin);
\draw[arr] (yin) -- (just);
\end{tikzpicture}
\legende{Schéma de la structure d'opinion : formule d'avis + opinion + 因为 + justification.}
\end{popfigure}

\courssub{Travail spécifique des tons}

La prononciation des tons est un savoir-faire explicitement visé par cette leçon. Trois mots du dialogue demandent une attention particulière.

\begin{propriete}[Trois mots clés à prononcer juste]
\begin{itemize}
\item \zh{觉得} \emph{juéde} : 2\textsuperscript{e} ton (montant) + \textbf{ton neutre}. Le second caractère \zh{得} se prononce court et léger, sans ton propre. Ne dites pas \emph{juédé} avec deux tons pleins.
\item \zh{重要} \emph{zhòngyào} : 4\textsuperscript{e} ton + 4\textsuperscript{e} ton. Deux tons descendants, francs et décidés : la voix tombe deux fois, comme deux coups secs. C'est le mot idéal pour insister dans un débat.
\item \zh{因为} \emph{yīnwèi} : 1\textsuperscript{er} ton (haut et plat) + 4\textsuperscript{e} ton (descendant). Attention : beaucoup de francophones disent \emph{yìnwéi} par analogie ; la prononciation correcte est bien \emph{yīnwèi}.
\end{itemize}
\end{propriete}

\begin{attention}[Le ton neutre de 得 dans 觉得]
En chinois, le ton neutre n'est pas un « cinquième ton » à inventer : c'est une syllabe prononcée \textbf{courte, légère et sans hauteur fixe}, dont la hauteur dépend du ton précédent. Après le 2\textsuperscript{e} ton de \zh{觉} (jué), le \zh{得} neutre retombe légèrement. Exagérez le contraste au début : \emph{jué-de} (long-court). Le même phénomène existe dans \zh{我们} (wǒmen), \zh{什么} (shénme), \zh{朋友} (péngyou).
\end{attention}

\begin{exoresolu}[Construire une réponse d'opinion complète]
Énoncé : un camarade te demande \zh{你觉得选举重要吗？} (Nǐ juéde xuǎnjǔ zhòngyào ma ? — « Penses-tu que les élections sont importantes ? »). Réponds en une phrase complète avec une justification, en utilisant \zh{觉得} et \zh{因为}.
\tcblower
Solution détaillée : on suit le schéma Formule d'avis + Opinion + 因为 + Justification.\\
1. Formule d'avis : \zh{我觉得} (wǒ juéde).\\
2. Opinion : \zh{选举很重要} (xuǎnjǔ hěn zhòngyào).\\
3. Justification avec \zh{因为} : \zh{因为每个人都可以选自己的代表} (yīnwèi měi ge rén dōu kěyǐ xuǎn zìjǐ de dàibiǎo — « parce que chacun peut choisir son propre représentant »).\\
Réponse complète : \zh{我觉得选举很重要，因为每个人都可以选自己的代表。} \emph{(Wǒ juéde xuǎnjǔ hěn zhòngyào, yīnwèi měi ge rén dōu kěyǐ xuǎn zìjǐ de dàibiǎo.)} « Je pense que les élections sont importantes, parce que chacun peut choisir son propre représentant. »\\
Vérification des tons : juéde (2 + neutre), zhòngyào (4 + 4), yīnwèi (1 + 4).
\end{exoresolu}

\courssub{Méthode et entraînement en binômes}

\begin{methode}[Réussir sa prise de parole : le plan Avis → Raison → Exemple]
Pour toute question d'opinion dans un débat, suis toujours ces trois étapes :
\begin{enumerate}
\item \textbf{Avis} : annonce ton point de vue avec \zh{我觉得} (oral) ou \zh{我认为} (formel). Exemple : \zh{我认为民主很重要。}
\item \textbf{Raison} : justifie avec \zh{因为}. Exemple : \zh{因为每个人都可以说自己的意见。}
\item \textbf{Exemple} : ajoute un cas concret avec \zh{比如} (bǐrú, « par exemple »). Exemple : \zh{比如，我们班每年选班长。} (Bǐrú, wǒmen bān měi nián xuǎn bānzhǎng. — « Par exemple, notre classe élit un délégué chaque année. »)
\end{enumerate}
Ce plan en trois temps donne une réponse claire, complète et naturelle — exactement ce qu'attend l'examinateur à l'oral.
\end{methode}

\begin{experience}[Entraînement en binômes : substitution du thème]
Par groupes de deux, reprenez le dialogue modèle en remplaçant le thème \zh{民主} par l'un des thèmes suivants : \zh{自由} (zìyóu, la liberté), \zh{学校的选举} (xuéxiào de xuǎnjǔ, les élections scolaires), \zh{学习汉语} (xuéxí Hànyǔ, l'apprentissage du chinois).
\begin{enumerate}
\item Élève A : pose la question \zh{你觉得...重要吗？}
\item Élève B : réponds avec le plan Avis → Raison → Exemple.
\item Changez de rôle, puis changez de thème.
\item Le professeur écoute et corrige d'abord les \textbf{tons} (juéde, zhòngyào, yīnwèi), puis la structure.
\end{enumerate}
\end{experience}

\begin{exercice}[À toi de parler]
Réponds oralement, par une phrase complète avec \zh{因为}, aux questions suivantes :\\
1. \zh{你觉得自由重要吗？} (Nǐ juéde zìyóu zhòngyào ma ?)\\
2. \zh{你觉得选举重要吗？} (Nǐ juéde xuǎnjǔ zhòngyào ma ?)\\
3. \zh{你觉得学习汉语重要吗？} (Nǐ juéde xuéxí Hànyǔ zhòngyào ma ?)
\end{exercice}

\begin{corrige}[Exercice « À toi de parler »]
Réponses possibles (toute réponse respectant le schéma Avis + 因为 + justification est acceptée) :\\
1. \zh{我觉得自由很重要，因为每个人都可以做自己喜欢的事。} \emph{(Wǒ juéde zìyóu hěn zhòngyào, yīnwèi měi ge rén dōu kěyǐ zuò zìjǐ xǐhuan de shì.)} « Je pense que la liberté est importante, parce que chacun peut faire ce qu'il aime. »\\
2. \zh{我认为选举很重要，因为公民可以选自己的代表。} \emph{(Wǒ rènwéi xuǎnjǔ hěn zhòngyào, yīnwèi gōngmín kěyǐ xuǎn zìjǐ de dàibiǎo.)} « Je pense que les élections sont importantes, parce que les citoyens peuvent choisir leurs représentants. »\\
3. \zh{我觉得学习汉语很重要，因为中国和喀麦隆有很多合作。} \emph{(Wǒ juéde xuéxí Hànyǔ hěn zhòngyào, yīnwèi Zhōngguó hé Kāmàilóng yǒu hěn duō hézuò.)} « Je pense qu'apprendre le chinois est important, parce que la Chine et le Cameroun ont beaucoup de coopération. »
\end{corrige}

\begin{aretenir}[L'essentiel de la section]
\begin{itemize}
\item Pour donner son avis : \textbf{Sujet + 觉得 / 认为 + proposition complète}, sans mot pour « que ». \zh{认为} est plus formel, \zh{在我看来} se place en tête de phrase.
\item Pour justifier : \zh{因为} seul (réponse courte) ou la corrélation \zh{因为...所以...}, correcte en chinois contrairement au français.
\item Tons à maîtriser : \emph{juéde} (2\textsuperscript{e} + neutre), \emph{zhòngyào} (4\textsuperscript{e} + 4\textsuperscript{e}), \emph{yīnwèi} (1\textsuperscript{er} + 4\textsuperscript{e}).
\item Plan de prise de parole : \textbf{Avis → Raison → Exemple} (avec \zh{比如}).
\end{itemize}
\end{aretenir}

Nous savons maintenant donner un avis simple et le justifier. Mais un débat, ce n'est pas seulement parler : c'est aussi réagir aux autres. La section suivante nous apprendra à exprimer l'accord, le désaccord et à relancer la conversation.

\coursec{Dialogue modèle 2 : accord, désaccord et relance}

Dans la section précédente, nous avons appris à \cle{donner son avis} avec des formules comme \zh{我认为…} (\emph{Wǒ rènwéi…} « je pense que… ») et \zh{我觉得…} (\emph{Wǒ juéde…} « je trouve que… »). Mais un débat ne s'arrête pas à une opinion : il faut ensuite \cle{réagir} à l'avis de l'autre — être d'accord ou pas — et \cle{relancer} la conversation pour que l'échange continue. C'est exactement ce que nous allons construire ici : un mini-débat complet, poli et vivant, sur le thème de la démocratie scolaire.

\courssub{Observation : un mini-débat en classe}

Lisons d'abord ce dialogue entre deux élèves de Terminale qui discutent de la vie de leur lycée à Yaoundé.

\begin{textebox}[Dialogue modèle 2 — 学生的意见 (L'avis des élèves)]
A：我认为学生应该参与学校的决定。\\
\emph{Wǒ rènwéi xuésheng yīnggāi cānyù xuéxiào de juédìng.}\\
« Je pense que les élèves devraient participer aux décisions de l'école. »\\[4pt]
B：你说得有道理，但是学生还太小。\\
\emph{Nǐ shuō de yǒu dàoli, dànshì xuésheng hái tài xiǎo.}\\
« Ce que tu dis se tient, mais les élèves sont encore trop jeunes. »\\[4pt]
A：为什么？你能举个例子吗？\\
\emph{Wèishénme? Nǐ néng jǔ ge lìzi ma?}\\
« Pourquoi ? Tu peux donner un exemple ? »\\[4pt]
B：比如，小学生不知道怎么选班长。\\
\emph{Bǐrú, xiǎoxuéshēng bù zhīdào zěnme xuǎn bānzhǎng.}\\
« Par exemple, les élèves du primaire ne savent pas comment élire un délégué. »\\[4pt]
A：你说得对，不过中学生可以。你觉得呢？\\
\emph{Nǐ shuō de duì, bùguò zhōngxuéshēng kěyǐ. Nǐ juéde ne?}\\
« Tu as raison, mais les collégiens et lycéens, eux, le peuvent. Et toi, qu'en penses-tu ? »\\[4pt]
B：嗯，我同意你的看法。\\
\emph{Ǹg, wǒ tóngyì nǐ de kànfǎ.}\\
« Hum, je suis d'accord avec ton point de vue. »
\end{textebox}

Observons la mécanique du dialogue : A donne son avis (\zh{我认为…}), B réagit avec un désaccord poli (\zh{你说得有道理，但是…}), A relance (\zh{为什么？你能举个例子吗？}), puis le débat se termine par un accord (\zh{我同意你的看法}). Trois mouvements : \textbf{avis → réaction → relance}. C'est le squelette de tout débat guidé.

\begin{definition}[Deux verbes pour « participer » : 参加 et 参与]
\zh{参加} (\emph{cānjiā}) s'emploie pour les activités concrètes : \zh{参加比赛} « participer à un match », \zh{参加活动} « participer à une activité ».\\
\zh{参与} (\emph{cānyù}) s'emploie pour les affaires, les décisions, la vie collective : \zh{参与学校的决定} « participer aux décisions de l'école ». Dans le débat citoyen, c'est 参与 qu'il faut choisir.
\end{definition}

\courssub{Les formules d'accord}

\begin{definition}[Formules d'accord]
\begin{itemize}
\item \zh{我同意。} \emph{Wǒ tóngyì.} « Je suis d'accord. » — formule neutre et directe.
\item \zh{我同意你的看法。} \emph{Wǒ tóngyì nǐ de kànfǎ.} « Je suis d'accord avec ton point de vue. »
\item \zh{你说得对。} \emph{Nǐ shuō de duì.} « Tu as raison (ce que tu dis est juste). »
\item \zh{你说得有道理。} \emph{Nǐ shuō de yǒu dàoli.} « Ce que tu dis se tient / il y a du vrai dans ce que tu dis. »
\item \zh{我也是这么想的。} \emph{Wǒ yě shì zhème xiǎng de.} « Je pense la même chose (moi aussi c'est ce que je pense). »
\end{itemize}
\end{definition}

\begin{propriete}[La structure Verbe + 得 + complément de manière]
Dans \zh{你说得对} et \zh{你说得有道理}, la particule \cle{得} (\emph{de}) introduit un \textbf{complément de manière} : elle dit \emph{comment} l'action est faite.\\
\textbf{Structure :} Verbe + 得 + adjectif / expression qualitative.\\
\zh{你说得对。} \emph{Nǐ shuō de duì.} « Tu parles juste → tu as raison. »\\
\zh{他说得很好。} \emph{Tā shuō de hěn hǎo.} « Il parle très bien. »\\
\zh{她写得很快。} \emph{Tā xiě de hěn kuài.} « Elle écrit très vite. »\\
\textbf{Attention à l'ordre :} le chinois dit littéralement « tu parles-DE-juste ». Le français inverse : « tu as raison ». Ne dites jamais \zh{你说对得} : 得 se place \emph{immédiatement après le verbe}.\\
\textbf{Exception :} si le verbe a un objet, on répète le verbe : \zh{他说汉语说得很好。} \emph{Tā shuō Hànyǔ shuō de hěn hǎo.} « Il parle très bien le chinois. » (litt. « il parle chinois parle-DE-très-bien »).
\end{propriete}

\courssub{Les formules de désaccord poli}

En chinois comme en français, dire brutalement « je ne suis pas d'accord » peut sembler sec. La courtoisie exige d'\textbf{atténuer} le désaccord.

\begin{definition}[Formules de désaccord, du plus direct au plus atténué]
\begin{itemize}
\item \zh{对不起，我不同意。} \emph{Duìbuqǐ, wǒ bù tóngyì.} « Pardon, je ne suis pas d'accord. »
\item \zh{我不这么认为。} \emph{Wǒ bù zhème rènwéi.} « Je ne le pense pas (litt. je ne pense pas ainsi). »
\item \zh{你说得有道理，但是…} \emph{Nǐ shuō de yǒu dàoli, dànshì…} « Ce que tu dis se tient, mais… »
\item \zh{你说得对，不过…} \emph{Nǐ shuō de duì, bùguò…} « Tu as raison, cependant… »
\end{itemize}
\end{definition}

\begin{propriete}[但是 et 不过 : les conjonctions d'atténuation]
\cle{但是} (\emph{dànshì}) et \cle{不过} (\emph{bùguò}) signifient tous deux « mais ».\\
\textbf{Structure :} concession positive + 但是/不过 + argument contraire.\\
\zh{你说得有道理，但是学生还太小。} \emph{Nǐ shuō de yǒu dàoli, dànshì xuésheng hái tài xiǎo.} « Ce que tu dis se tient, mais les élèves sont encore trop jeunes. »\\
\textbf{Nuance :} 但是 est plus fort et plus formel (« mais ») ; 不过 est plus léger et plus oral (« cependant, il n'empêche que »). Dans un débat courtois, 不过 adoucit la contradiction.\\
\textbf{Comparaison avec le français :} le français utilise « mais », « cependant », « par contre » ; le chinois fait de même, mais la concession qui précède (\zh{你说得有道理}) est presque obligatoire dans un échange poli, alors qu'en français on peut attaquer directement par « mais ».
\end{propriete}

\begin{exemplebox}[Exemples traduits]
\zh{对不起，我不同意。} \emph{Duìbuqǐ, wǒ bù tóngyì.} « Pardon, je ne suis pas d'accord. »\\
\zh{我不这么认为。} \emph{Wǒ bù zhème rènwéi.} « Je ne vois pas les choses ainsi. »\\
\zh{你说得有道理，但是我觉得学生可以参与。} \emph{Nǐ shuō de yǒu dàoli, dànshì wǒ juéde xuésheng kěyǐ cānyù.} « Ce que tu dis se tient, mais je trouve que les élèves peuvent participer. »\\
\zh{你说得对，不过我们也要听老师的意见。} \emph{Nǐ shuō de duì, bùguò wǒmen yě yào tīng lǎoshī de yìjiàn.} « Tu as raison, cependant nous devons aussi écouter l'avis des professeurs. »
\end{exemplebox}

\begin{methode}[Réussir un désaccord poli en trois étapes]
\begin{enumerate}
\item \textbf{Reconnaître} d'abord l'avis de l'autre : \zh{你说得有道理。} ou \zh{你说得对。} (« ce que tu dis se tient »).
\item \textbf{Introduire} la contradiction avec \zh{但是} ou \zh{不过}.
\item \textbf{Donner} son propre argument : \zh{但是我觉得…} + raison.
\end{enumerate}
Modèle complet : \zh{你说得有道理，但是我觉得中学生可以选班长。} \emph{Nǐ shuō de yǒu dàoli, dànshì wǒ juéde zhōngxuéshēng kěyǐ xuǎn bānzhǎng.} « Ce que tu dis se tient, mais je pense que les lycéens peuvent élire un délégué. »
\end{methode}

\courssub{Les formules de relance}

Un bon débatteur fait parler les autres. Le chinois dispose de petites questions toutes simples pour relancer la parole.

\begin{definition}[Formules de relance]
\begin{itemize}
\item \zh{你觉得呢？} \emph{Nǐ juéde ne?} « Et toi, qu'en penses-tu ? »
\item \zh{你呢？} \emph{Nǐ ne?} « Et toi ? »
\item \zh{为什么？} \emph{Wèishénme?} « Pourquoi ? »
\item \zh{你能举个例子吗？} \emph{Nǐ néng jǔ ge lìzi ma?} « Tu peux donner un exemple ? »
\item \zh{还有别的想法吗？} \emph{Hái yǒu bié de xiǎngfǎ ma?} « Y a-t-il d'autres idées ? »
\end{itemize}
\end{definition}

\begin{attention}[La particule finale 呢 ne s'oublie pas !]
Les francophones oublient souvent le \cle{呢} final dans \zh{你觉得呢？}. Sans 呢, \zh{你觉得} signifie « tu penses / tu trouves » : ce n'est plus une question, mais une affirmation (ou une phrase incomplète). La particule 呢 transforme l'énoncé en \textbf{relance} : « et toi ? ». Même logique pour \zh{你呢？} : sans 呢, il ne reste que « toi », qui ne veut rien dire seul. Pensez-y : 呢 = le « et… ? » du français.
\end{attention}

\courssub{Tableau récapitulatif des formules du débat}

\begin{center}
\begin{tabularx}{\linewidth}{|X|X|X|}
\hline
\rowcolor{popblueL} Fonction & Formule (caractères + pinyin) & Traduction \\
\hline
Donner son avis & 我认为… \emph{Wǒ rènwéi…} & Je pense que… \\
\hline
Accord neutre & 我同意。 \emph{Wǒ tóngyì.} & Je suis d'accord. \\
\hline
Accord chaleureux & 你说得对。 \emph{Nǐ shuō de duì.} & Tu as raison. \\
\hline
Accord partiel & 你说得有道理。 \emph{Nǐ shuō de yǒu dàoli.} & Ce que tu dis se tient. \\
\hline
Désaccord direct & 对不起，我不同意。 \emph{Duìbuqǐ, wǒ bù tóngyì.} & Pardon, je ne suis pas d'accord. \\
\hline
Désaccord atténué & …但是… / …不过… \emph{dànshì / bùguò} & …mais… / …cependant… \\
\hline
Relance d'opinion & 你觉得呢？ \emph{Nǐ juéde ne?} & Et toi, qu'en penses-tu ? \\
\hline
Relance simple & 你呢？ \emph{Nǐ ne?} & Et toi ? \\
\hline
Demande de raison & 为什么？ \emph{Wèishénme?} & Pourquoi ? \\
\hline
Demande d'exemple & 你能举个例子吗？ \emph{Nǐ néng jǔ ge lìzi ma?} & Tu peux donner un exemple ? \\
\hline
\end{tabularx}
\end{center}

\courssub{La mécanique du débat : organigramme}

\begin{popfigure}
\begin{tikzpicture}[
  node distance=0.55cm and 1.1cm,
  box/.style={draw=popblue, thick, rounded corners=3pt, fill=popblueL, align=center, font=\small, inner sep=5pt},
  box2/.style={draw=poporange, thick, rounded corners=3pt, fill=poporangeL, align=center, font=\small, inner sep=5pt},
  box3/.style={draw=popgreen, thick, rounded corners=3pt, fill=popgreenL, align=center, font=\small, inner sep=5pt},
  box4/.style={draw=poppurple, thick, rounded corners=3pt, fill=poppurpleL, align=center, font=\small, inner sep=5pt},
  fl/.style={-{Stealth[length=2.5mm]}, thick, popdark}
]
\node[box, align=center] (avis){1. Donner son avis\\我认为… / 我觉得…};
\node[box2, right=of avis, align=center] (accord){2a. Accord\\你说得对。\\我同意。};
\node[box2, below=of accord, align=center] (desaccord){2b. Désaccord poli\\你说得有道理，\\但是…};
\node[box3, right=of accord, align=center] (relance){3. Relance\\你觉得呢？\\你能举个例子吗？};
\node[box4, right=of relance, align=center] (conclu){4. Conclusion\\我同意你的看法。};
\draw[fl] (avis) -- (accord);
\draw[fl] (avis.east) -- ++(0.35,0) |- (desaccord.west);
\draw[fl] (accord) -- (relance);
\draw[fl] (desaccord.east) -- ++(0.35,0) |- (relance.west);
\draw[fl] (relance) -- (conclu);
\end{tikzpicture}
\legende{Organigramme du mini-débat : avis, réaction (accord ou désaccord poli), relance, conclusion.}
\end{popfigure}

\courssub{Travail des tons : 同意 et 例子}

Deux mots-clés de cette leçon méritent une attention phonétique particulière.

\begin{propriete}[Prononciation de 同意 et 例子]
\textbf{同意} \emph{tóngyì} = 2\textsuperscript{e} ton + 4\textsuperscript{e} ton. Le 2\textsuperscript{e} ton monte (comme quand on demande « hein ? » en français), le 4\textsuperscript{e} ton descend brusquement (comme un « non ! » sec). Dites : tóng↗ yì↘. Une erreur fréquente est de prononcer \emph{tōngyì} (1\textsuperscript{er} + 4\textsuperscript{e}), ce qui sonne plat et peut troubler l'écoute.\\
\textbf{例子} \emph{lìzi} = 4\textsuperscript{e} ton + ton neutre. Le 4\textsuperscript{e} ton descend fort sur \emph{lì}, puis la syllabe \emph{zi} se dit légèrement, sans ton propre, plus courte et plus basse. Ne mettez pas de ton plein sur \emph{zi} : \emph{lìzǐ} est une faute de prononciation courante chez les francophones, qui accentuent toutes les syllabes comme en français.
\end{propriete}

\begin{exemplebox}[Entraînement oral des tons]
Répétez après le professeur, en exagérant d'abord les tons :\\
\zh{我同意。} \emph{Wǒ tóngyì.} (3\textsuperscript{e} + 2\textsuperscript{e} + 4\textsuperscript{e})\\
\zh{我不同意。} \emph{Wǒ bù tóngyì.} (3\textsuperscript{e} + 4\textsuperscript{e} + 2\textsuperscript{e} + 4\textsuperscript{e})\\
\zh{你能举个例子吗？} \emph{Nǐ néng jǔ ge lìzi ma?} (ton neutre sur \emph{zi} et sur \emph{ma})\\
Astuce : frappez doucement dans vos mains sur chaque 4\textsuperscript{e} ton pour marquer la chute.
\end{exemplebox}

\begin{exoresolu}[Construire un désaccord poli]
Énoncé : votre camarade dit \zh{我认为学生不应该参与学校的决定。} (« Je pense que les élèves ne devraient pas participer aux décisions de l'école. »). Vous n'êtes pas d'accord. Construisez une réponse polie en trois étapes, puis relancez un troisième camarade.
\tcblower
Solution détaillée :
\begin{enumerate}
\item Reconnaissance de l'avis : \zh{你说得有道理，} \emph{Nǐ shuō de yǒu dàoli,} « Ce que tu dis se tient, »
\item Conjonction d'atténuation : \zh{但是} \emph{dànshì} « mais »
\item Argument personnel : \zh{我觉得学生应该参与，因为这是我们的学校。} \emph{wǒ juéde xuésheng yīnggāi cānyù, yīnwèi zhè shì wǒmen de xuéxiào.} « je trouve que les élèves devraient participer, parce que c'est notre école. »
\item Relance : \zh{你觉得呢？} \emph{Nǐ juéde ne?} « Et toi, qu'en penses-tu ? »
\end{enumerate}
Réponse complète : \zh{你说得有道理，但是我觉得学生应该参与，因为这是我们的学校。你觉得呢？}\\
« Ce que tu dis se tient, mais je trouve que les élèves devraient participer, parce que c'est notre école. Et toi, qu'en penses-tu ? »
\end{exoresolu}

\begin{attention}[Pièges fréquents des francophones]
\begin{itemize}
\item \textbf{Calquer « je suis d'accord avec toi » :} ne dites pas \zh{我是同意你}. La bonne formule est \zh{我同意你的看法} (ou simplement \zh{我同意}). Le verbe 同意 suffit, sans 是.
\item \textbf{Placer 但是 au mauvais endroit :} 但是 se met entre les deux idées, jamais à la fin de la phrase : \zh{你说得对，但是学生还太小。} et non \zh{你说得对，学生还太小但是。}.
\item \textbf{Oublier le ton neutre :} dans \zh{例子} (\emph{lìzi}) et dans la particule \zh{吗} (\emph{ma}), la dernière syllabe est légère. Tout accentuer donne un chinois robotique.
\item \textbf{Oublier le classificateur 个 :} dans la formule complète, on dit \zh{举个例子} (\emph{jǔ ge lìzi}) avec le classificateur 个 : \zh{你能举个例子吗？} et non \zh{你能举例子吗？} dans ce contexte de politesse.
\end{itemize}
\end{attention}

\begin{exercice}[À vous de débattre]
Par groupes de trois, jouez le mini-débat suivant : l'élève A donne son avis avec \zh{我认为…} sur le sujet « les élèves devraient choisir leurs délégués de classe » ; l'élève B réagit avec un accord (\zh{你说得对}) ou un désaccord poli (\zh{你说得有道理，但是…}) ; l'élève C relance avec \zh{为什么？} ou \zh{你能举个例子吗？}. Changez ensuite les rôles. Surveillez les tons de \zh{同意} (\emph{tóngyì}) et de \zh{例子} (\emph{lìzi}).
\end{exercice}

\begin{corrige}[Exercice — À vous de débattre]
Production libre : il n'y a pas de réponse unique. Voici un modèle possible.\\
A : \zh{我认为学生应该选班长。} \emph{Wǒ rènwéi xuésheng yīnggāi xuǎn bānzhǎng.} « Je pense que les élèves devraient élire le délégué. »\\
B : \zh{你说得对，不过老师也应该帮助我们。} \emph{Nǐ shuō de duì, bùguò lǎoshī yě yīnggāi bāngzhù wǒmen.} « Tu as raison, cependant le professeur devrait aussi nous aider. »\\
C : \zh{为什么？你能举个例子吗？} \emph{Wèishénme? Nǐ néng jǔ ge lìzi ma?} « Pourquoi ? Tu peux donner un exemple ? »\\
Critères de réussite : la particule 呢 est présente dans les relances ; le désaccord commence par une concession ; les tons 2\textsuperscript{e} + 4\textsuperscript{e} de 同意 sont nets ; le ton neutre de 例子 est léger.
\end{corrige}

\begin{aretenir}[L'essentiel de la section]
\begin{itemize}
\item Accord : \zh{我同意} / \zh{你说得对} / \zh{你说得有道理} / \zh{我也是这么想的}.
\item Désaccord poli en trois temps : reconnaître (\zh{你说得有道理}) + atténuer (\zh{但是} / \zh{不过}) + argumenter.
\item Relance : \zh{你觉得呢？} / \zh{你呢？} / \zh{为什么？} / \zh{你能举个例子吗？} — sans oublier la particule 呢.
\item La particule 得 introduit le complément de manière : Verbe + 得 + adjectif (\zh{你说得对}) ; avec un objet, on répète le verbe (\zh{他说汉语说得很好}).
\item « Participer » aux décisions se dit \zh{参与} (\emph{cānyù}) ; \zh{参加} (\emph{cānjiā}) sert pour les activités concrètes.
\item Tons à soigner : \emph{tóngyì} (2\textsuperscript{e} + 4\textsuperscript{e}) et \emph{lìzi} (4\textsuperscript{e} + neutre).
\end{itemize}
\end{aretenir}

\coursec{Structurer un exposé simple}

Dans la section précédente, nous avons appris à réagir dans un débat : donner son avis, être d'accord ou pas d'accord, relancer la parole. Mais dans un débat comme à l'examen, il arrive un moment où l'on doit parler \textbf{seul, pendant une ou deux minutes, sans être interrompu} : c'est l'\cle{exposé}. Un bon exposé en chinois n'est pas un long texte récité par cœur : c'est un discours \textbf{court, clair et bien structuré}, que l'auditoire peut suivre facilement. Cette section vous donne le plan type, les connecteurs et les formules pour y arriver.

\courssub{Observer d'abord : un mini-exposé modèle}

Lisons d'abord un exposé complet de six phrases, tel que l'enseignant le lit en classe, avant de l'analyser.

\begin{textebox}[Un mini-exposé sur la démocratie]
大家好！今天我要谈谈民主。首先，我认为民主很重要。然后，每个公民都应该投票。最后，民主给我们带来自由。总之，我们要保护民主。我的讲话完了，谢谢大家！\\
\emph{Dàjiā hǎo! Jīntiān wǒ yào tántan mínzhǔ. Shǒuxiān, wǒ rènwéi mínzhǔ hěn zhòngyào. Ránhòu, měi ge gōngmín dōu yīnggāi tóupiào. Zuìhòu, mínzhǔ gěi wǒmen dàilái zìyóu. Zǒngzhī, wǒmen yào bǎohù mínzhǔ. Wǒ de jiǎnghuà wán le, xièxie dàjiā!}\\
« Bonjour à tous ! Aujourd'hui, je vais parler de la démocratie. D'abord, je pense que la démocratie est très importante. Ensuite, chaque citoyen devrait voter. Enfin, la démocratie nous apporte la liberté. En résumé, nous devons protéger la démocratie. Mon exposé est terminé, merci à tous ! »
\end{textebox}

\textbf{Observons.} Combien de phrases ? Six seulement. Pourtant, l'exposé est complet : il salue, annonce le sujet, donne un avis, avance deux arguments, conclut et remercie. Remarquez les mots en début de phrase : \zh{首先}, \zh{然后}, \zh{最后}, \zh{总之}. Ce sont les \cle{connecteurs} : ce sont eux qui donnent sa structure à l'exposé. Sans eux, les phrases se succèdent sans logique ; avec eux, l'auditeur sait toujours où il en est.

\courssub{Le plan type d'un mini-exposé}

\begin{propriete}[Plan type en six étapes]
Un mini-exposé de niveau Terminale suit toujours la même architecture :
\begin{enumerate}
\item \textbf{Salutation} : \zh{大家好！} (\emph{Dàjiā hǎo!}) « Bonjour à tous ! »
\item \textbf{Introduction du sujet} : \zh{今天我要谈谈……。} « Aujourd'hui, je vais parler de… »
\item \textbf{Avis personnel} : \zh{我认为……。} « Je pense que… »
\item \textbf{Arguments} : deux arguments reliés si besoin par \zh{因为} (parce que) et \zh{所以} (donc), introduits par \zh{首先} et \zh{然后}.
\item \textbf{Exemple} : une phrase avec \zh{比如} (par exemple) ou un fait simple.
\item \textbf{Conclusion et clôture} : \zh{总之，……。我的讲话完了，谢谢大家！}
\end{enumerate}
\end{propriete}

Comparaison avec le français : en français on dit « dans un premier temps… ensuite… pour conclure » ; le chinois fait exactement pareil, mais les connecteurs sont \textbf{quasi obligatoires} à l'oral scolaire : un exposé sans \zh{首先} / \zh{然后} / \zh{最后} paraît décousu aux oreilles d'un sinophone. Autre différence : le chinois aime les formules de politesse \textbf{fixes} à l'ouverture (\zh{大家好！}) et à la clôture (\zh{谢谢大家！}), qu'il ne faut ni traduire mot à mot ni oublier.

\begin{popfigure}
\begin{tikzpicture}[node distance=0.35cm]
\node[draw=popblue, fill=popblueL, rounded corners, minimum width=2.6cm, minimum height=0.9cm, align=center, font=\small] (a) {大家好\\ \emph{salutation}};
\node[draw=popgreen, fill=popgreenL, rounded corners, minimum width=2.6cm, minimum height=0.9cm, align=center, font=\small, right=of a] (b) {今天我要谈谈…\\ \emph{sujet}};
\node[draw=poporange, fill=poporangeL, rounded corners, minimum width=2.6cm, minimum height=0.9cm, align=center, font=\small, right=of b] (c) {首先…\\ \emph{argument 1}};
\node[draw=poporange, fill=poporangeL, rounded corners, minimum width=2.6cm, minimum height=0.9cm, align=center, font=\small, below=0.5cm of c] (d) {然后…\\ \emph{argument 2}};
\node[draw=poppurple, fill=poppurpleL, rounded corners, minimum width=2.6cm, minimum height=0.9cm, align=center, font=\small, left=of d] (e) {最后…\\ \emph{exemple}};
\node[draw=poppink, fill=poppinkL, rounded corners, minimum width=2.6cm, minimum height=0.9cm, align=center, font=\small, left=of e] (f) {总之…\\ \emph{conclusion}};
\node[draw=popgold, fill=popgoldL, rounded corners, minimum width=2.6cm, minimum height=0.9cm, align=center, font=\small, below=0.5cm of f] (g) {谢谢大家！\\ \emph{remerciement}};
\draw[-{Stealth}, thick, popdark] (a) -- (b);
\draw[-{Stealth}, thick, popdark] (b) -- (c);
\draw[-{Stealth}, thick, popdark] (c) -- (d);
\draw[-{Stealth}, thick, popdark] (d) -- (e);
\draw[-{Stealth}, thick, popdark] (e) -- (f);
\draw[-{Stealth}, thick, popdark] (f) -- (g);
\end{tikzpicture}
\legende{Frise de l'exposé : les sept jalons à suivre, de la salutation au remerciement.}
\end{popfigure}

\courssub{Les formules d'ouverture et de clôture}

\begin{definition}[Formules d'ouverture]
\begin{itemize}
\item \zh{大家好！} \emph{Dàjiā hǎo!} — « Bonjour à tous ! » (salutation standard devant la classe).
\item \zh{今天我要谈谈民主。} \emph{Jīntiān wǒ yào tántan mínzhǔ.} — « Aujourd'hui, je vais parler de la démocratie. » Structure : \textbf{aujourd'hui + je + vais + parler (un peu) de + sujet}. Le verbe redoublé \zh{谈谈} adoucit le propos : « parler un peu, dire quelques mots ».
\item Variante : \zh{我想谈一谈民主的问题。} \emph{Wǒ xiǎng tán yi tán mínzhǔ de wèntí.} — « Je voudrais parler de la question de la démocratie. »
\end{itemize}
\end{definition}

\begin{definition}[Formule de clôture]
\zh{我的讲话完了，谢谢大家！} \emph{Wǒ de jiǎnghuà wán le, xièxie dàjiā!} — « Mon exposé est terminé, merci à tous ! » Le \zh{了} marque ici l'accomplissement : l'exposé « est arrivé à son terme ». Attention : on ne peut pas dire \zh{谢谢大家的听} (calque du français « merci de votre écoute ») ; on dit simplement \zh{谢谢大家！} ou, dans un registre soutenu, \zh{谢谢大家的聆听！} \emph{Xièxie dàjiā de língtīng!} « Merci à tous de votre écoute attentive. »
\end{definition}

\courssub{Les connecteurs de l'exposé}

\begin{center}
\begin{tabularx}{\linewidth}{|X|X|X|X|}
\hline
\rowcolor{popblueL} Caractères & Pinyin & Nature & Traduction \\
\hline
首先 & shǒuxiān & connecteur & d'abord, en premier lieu \\
\hline
然后 & ránhòu & connecteur & ensuite, puis \\
\hline
最后 & zuìhòu & connecteur & enfin, pour finir \\
\hline
总之 & zǒngzhī & connecteur & en résumé, bref \\
\hline
比如 & bǐrú & connecteur & par exemple \\
\hline
因为 & yīnwei & conjonction & parce que \\
\hline
所以 & suǒyǐ & conjonction & donc, c'est pourquoi \\
\hline
我认为 & wǒ rènwéi & expression & je pense que, j'estime que \\
\hline
应该 & yīnggāi & verbe modal & devoir (il faut, il faudrait) \\
\hline
保护 & bǎohù & verbe & protéger \\
\hline
\end{tabularx}
\end{center}

\begin{propriete}[Emploi des connecteurs]
\begin{itemize}
\item \zh{首先} ouvre le \textbf{premier} argument ; on ne l'utilise qu'une fois.
\item \zh{然后} enchaîne les étapes suivantes ; il peut se répéter.
\item \zh{最后} introduit le dernier point (souvent l'exemple ou la conséquence).
\item \zh{总之} annonce la \textbf{conclusion générale} : il ne s'emploie qu'en toute fin, jamais au milieu.
\item Paire logique : \zh{因为……，所以……} « parce que…, donc… ». En chinois, on peut employer les deux ensemble dans la même phrase, contrairement au français où « parce que… donc… » est incorrect.
\end{itemize}
Exemple : \zh{因为民主给我们带来自由，所以我们要保护民主。} \emph{Yīnwei mínzhǔ gěi wǒmen dàilái zìyóu, suǒyǐ wǒmen yào bǎohù mínzhǔ.} « Parce que la démocratie nous apporte la liberté, nous devons donc la protéger. »
\end{propriete}

\begin{exemplebox}[Exemples traduits]
\zh{总之，民主对我们的生活很重要。}\\
\emph{Zǒngzhī, mínzhǔ duì wǒmen de shēnghuó hěn zhòngyào.}\\
« En résumé, la démocratie est importante pour notre vie. »\\[4pt]
\zh{首先，我认为每个公民都应该投票。}\\
\emph{Shǒuxiān, wǒ rènwéi měi ge gōngmín dōu yīnggāi tóupiào.}\\
« D'abord, je pense que chaque citoyen devrait voter. »\\[4pt]
\zh{然后，年轻人也应该关心国家的事情。}\\
\emph{Ránhòu, niánqīngrén yě yīnggāi guānxīn guójiā de shìqing.}\\
« Ensuite, les jeunes devraient aussi s'intéresser aux affaires du pays. »\\[4pt]
\zh{最后，比如在喀麦隆，很多人参加选举。}\\
\emph{Zuìhòu, bǐrú zài Kāmàilóng, hěn duō rén cānjiā xuǎnjǔ.}\\
« Enfin, par exemple au Cameroun, beaucoup de gens participent aux élections. »
\end{exemplebox}

\begin{attention}[Piège de prononciation : 谈谈]
\zh{谈谈} se prononce \emph{tántan} : le deuxième \zh{谈} passe au \textbf{ton neutre} (reduplication adoucissante), et non deux 2\textsuperscript{e} tons. Dire \emph{tántán} sonne artificiel. Même phénomène avec \zh{看看} (\emph{kànkan}), \zh{想想} (\emph{xiǎngxiang}). Autre piège : ne confondez pas \zh{总之} (\emph{zǒngzhī}, « en résumé ») avec \zh{总是} (\emph{zǒngshì}, « toujours ») : sons presque identiques, sens différents !
\end{attention}

\courssub{Mémoriser son exposé et gérer le stress}

\begin{methode}[Retenir le plan, pas le texte]
\begin{enumerate}
\item Écrivez votre exposé en six phrases, une par étape du plan.
\item Soulignez les quatre connecteurs : \zh{首先}, \zh{然后}, \zh{最后}, \zh{总之}.
\item Mémorisez \textbf{un seul mot-clé par étape} : sujet → avis → voter → liberté → exemple → protéger.
\item Répétez à voix haute en ne regardant que les mots-clés, jamais le texte complet.
\item Le jour J : parlez \textbf{lentement}, marquez bien les tons, regardez l'auditoire (pas vos notes), respirez avant chaque connecteur.
\end{enumerate}
\end{methode}

\begin{exoresolu}[Remettre un exposé dans l'ordre]
Énoncé : remettez ces phrases dans le bon ordre pour former un exposé cohérent.\\
(a) 总之，我们要保护民主。 (b) 大家好！今天我要谈谈民主。 (c) 最后，民主给我们带来自由。 (d) 首先，我认为民主很重要。 (e) 我的讲话完了，谢谢大家！ (f) 然后，每个公民都应该投票。
\tcblower
Solution : l'ordre correct est \textbf{b → d → f → c → a → e}. On reconnaît la salutation et l'annonce du sujet (b) en tête ; \zh{首先} (d) ouvre les arguments, \zh{然后} (f) enchaîne, \zh{最后} (c) donne le dernier point ; \zh{总之} (a) conclut ; la formule de clôture (e) termine. Traduction : « Bonjour à tous ! Aujourd'hui je vais parler de la démocratie. D'abord, je pense qu'elle est très importante. Ensuite, chaque citoyen devrait voter. Enfin, elle nous apporte la liberté. En résumé, nous devons la protéger. Mon exposé est terminé, merci à tous ! »
\end{exoresolu}

\begin{exercice}[À vous de construire]
Rédigez un mini-exposé de six phrases sur le thème « l'école » en suivant le plan type : salutation, \zh{今天我要谈谈学校。}, avis avec \zh{我认为}, deux arguments avec \zh{首先} et \zh{然后}, conclusion avec \zh{总之}, clôture. Puis entraînez-vous à le dire en ne regardant que vos mots-clés.
\end{exercice}

\begin{corrige}[Exercice « À vous de construire »]
Modèle possible : \zh{大家好！今天我要谈谈学校。首先，我认为学校很重要，因为我们在学校学习知识。然后，学校也教我们怎么和别人一起生活。最后，比如在我们的学校，老师帮助我们进步。总之，我们要爱我们的学校。我的讲话完了，谢谢大家！} « Bonjour à tous ! Aujourd'hui je vais parler de l'école. D'abord, je pense que l'école est très importante, parce que nous y apprenons des connaissances. Ensuite, l'école nous apprend aussi à vivre avec les autres. Enfin, par exemple dans notre école, les professeurs nous aident à progresser. En résumé, nous devons aimer notre école. Mon exposé est terminé, merci à tous ! » Tout exposé respectant le plan et les connecteurs est accepté.
\end{corrige}

\begin{savaistu}[Le jour du Bac]
À l'oral du Bac, l'examinateur apprécie bien plus un exposé \textbf{court mais bien structuré} — salutation, connecteurs visibles, conclusion nette — qu'un long texte récité sans logique. En Chine aussi, on enseigne aux élèves que \zh{说得清楚比说得多好} (\emph{shuō de qīngchu bǐ shuō de duō hǎo}) : « parler clairement vaut mieux que parler longtemps ». Six phrases bien construites suffisent pour convaincre !
\end{savaistu}

\coursec{Travail spécifique des tons}

\courssub{Transition et objectifs}

Après avoir appris à \emph{structurer un exposé simple} (section précédente), nous devons maintenant veiller à ce que la \emph{prononciation} serve le propos. En expression orale, une mauvaise tonalité peut transformer « démocratie » en « peuple sensible » ou « élection » en « sélectionner un bureau ». Cette section vise à \textbf{automatiser la mélodie des phrases} du thème grâce à une révision systématique, un travail sur les paires minimales, l'application des règles de \emph{sandhi} (changements de tons en contexte) et un entraînement rythmique des formules fonctionnelles.

\courssub{Révision des quatre tons et du ton neutre — Gestes de la main}

\begin{textebox}[Rappel : les quatre tons du mandarin]
En chinois, la hauteur et le contour de la voix sur une syllabe distinguent le sens des mots. On utilise la main pour \emph{visualiser} la courbe de chaque ton.
\begin{itemize}
    \item \textbf{1er ton (阴平 yīnpíng) :} \emph{Plat et haut.} Main horizontale, paume vers le bas, niveau front. Exemple : \zh{妈 mā} (maman).
    \item \textbf{2e ton (阳平 yángpíng) :} \emph{Montant.} Main partie du bas vers le haut, comme une question surprise. Exemple : \zh{麻 má} (chanvre).
    \item \textbf{3e ton (上声 shǎngshēng) :} \emph{Descendant-montant (en « V »).} Main descend puis remonte. \textbf{Attention :} en pratique courante, il est souvent \emph{demi-3e ton} (seulement descendant bas) quand il n'est pas en fin de groupe rythmique. Exemple : \zh{马 mǎ} (cheval).
    \item \textbf{4e ton (去声 qùshēng) :} \emph{Descendant fort.} Main part du haut et frappe vers le bas, comme un ordre sec. Exemple : \zh{骂 mà} (gronder).
    \item \textbf{Ton neutre (轻声 qīngshēng) :} \emph{Court, léger, sans contour propre.} Main ouverte, geste vif et relâché. Sa hauteur dépend du ton précédent. Exemple : \zh{妈妈 māma} (maman).
\end{itemize}
\end{textebox}

\begin{popfigure}
\begin{tikzpicture}[scale=1.2, >=Stealth]
    % Axes
    \draw[popmuted, thin] (0,0) -- (5,0) node[right] {Temps};
    \draw[popmuted, thin] (0,0) -- (0,4.5) node[above] {Hauteur};
    % Repères
    \foreach \y/\label in {1/bas, 2/moyen, 3/haut, 4/très haut} {
        \draw[popline, dashed, very thin] (0,\y) -- (5,\y);
    }
    % 1er ton
    \draw[popblue, very thick] (0.2,3.5) -- (1.3,3.5);
    \node[popblue, left] at (0.2,3.5) {1er};
    % 2e ton
    \draw[poporange, very thick] (1.7,1.5) .. controls (2.2,2) and (2.8,3.5) .. (3.3,4);
    \node[poporange, left] at (1.7,1.5) {2e};
    % 3e ton (citation form)
    \draw[popgreen, very thick] (3.7,2.5) .. controls (4.0,1.5) and (4.5,1.5) .. (4.8,2.5);
    \node[popgreen, left] at (3.7,2.5) {3e};
    % 4e ton
    \draw[poppink, very thick] (0.2,4) -- (1.3,1);
    \node[poppink, left] at (0.2,4) {4e};
    % Légendes gestuelles
    \node[align=center, font=\small, popdark] at (0.75, -0.5) {Main plate\\haut};
    \node[align=center, font=\small, popdark] at (2.5, -0.5) {Main monte\\bas $\to$ haut};
    \node[align=center, font=\small, popdark] at (4.25, -0.5) {Main descend\\puis remonte (V)};
    \node[align=center, font=\small, popdark] at (0.75, -1.3) {Main tombe\\haut $\to$ bas};
\end{tikzpicture}
\legende{Contours des quatre tons (forme de citation) et gestes de la main associés.}
\end{popfigure}

\begin{methode}[Entraînement gestuel « Main-Ton »]
\begin{enumerate}
    \item L'enseignant prononce une syllabe (ex. \zh{ma}) ; les élèves reproduisent \textbf{en même temps} le geste de la main correspondant au ton entendu.
    \item Inverser : l'enseignant fait le geste ; les élèves chantent la syllabe sur le bon ton.
    \item Appliquer aux mots du thème : \zh{民主 mínzhǔ}, \zh{选举 xuǎnjǔ}, \zh{自由 zìyóu}, \zh{投票 tóupiào}.
\end{enumerate}
\end{methode}

\courssub{Paires minimales du thème — Discrimination auditive}

Le thème « démocratie » oppose des mots dont seul le ton change. L'oreille francophone, peu sensible à la hauteur lexicale, doit s'entraîner à les distinguer.

\begin{center}
\begin{tabularx}{\linewidth}{|X|X|X|X|}
\hline
\rowcolor{popblueL} \textbf{Caractères} & \textbf{Pinyin} & \textbf{Nature} & \textbf{Traduction} \\
\hline
民 & mín & nom & peuple, population \\
\hline
敏 & mǐn & adj & sensible, vif, rapide \\
\hline
义 & yì & nom & justice, sens, signification \\
\hline
一 & yī & num & un \\
\hline
权 & quán & nom & droit, pouvoir, autorité \\
\hline
全 & quán & adj & tout, entier, complet \\
\hline
主 & zhǔ & nom/verbe & maître, principal, diriger \\
\hline
住 & zhù & verbe & habiter, rester, tenir \\
\hline
\end{tabularx}
\end{center}

\begin{exemplebox}[Paires minimales en contexte]
\begin{itemize}
    \item \zh{民主 mínzhǔ} (démocratie) \textbf{vs} \zh{敏主 mǐnzhǔ} (n'existe pas, mais confusion possible : « peuple sensible »).
    \item \zh{意义 yìyì} (signification) \textbf{vs} \zh{一一 yīyī} (un par un).
    \item \zh{权利 quánlì} (droit) \textbf{vs} \zh{全力 quánlì} (toutes ses forces).
    \item \zh{主席 zhǔxí} (président) \textbf{vs} \zh{住西 zhù xī} (habiter à l'ouest).
\end{itemize}
\end{exemplebox}

\begin{experience}[Discrimination auditive : « Levez la carte »]
\begin{enumerate}
    \item Chaque élève prépare 4 cartes numérotées 1, 2, 3, 4 (et une carte Ø pour ton neutre).
    \item L'enseignant lit \textbf{une seule syllabe} ou un \textbf{mot de deux syllabes} (liste ci-dessous). Pas de phrase complète pour isoler la perception tonale.
    \item Les élèves lèvent \textbf{immédiatement} la carte correspondant au ton de la \textbf{première syllabe} (ou de la syllabe unique).
    \item L'enseignant scanne la salle, corrige collectivement, puis fait répéter en chœur avec le geste de la main.
\end{enumerate}
\textbf{Liste pour l'exercice (à lire dans le désordre) :}
\begin{itemize}
    \item \zh{mín} (2), \zh{mǐn} (3), \zh{yì} (4), \zh{yī} (1)
    \item \zh{mínzhǔ} (2-3), \zh{xuǎnjǔ} (3-3 $\to$ 2-3), \zh{zìyóu} (4-2), \zh{tóupiào} (2-4)
    \item \zh{quánlì} (2-4), \zh{zhǔzhāng} (3-1), \zh{píngděng} (2-3), \zh{jiànkāng} (4-1)
\end{itemize}
\end{experience}

\begin{exoresolu}[Discrimination : mín vs mǐn]
\textbf{Énoncé :} L'enseignant prononce : \zh{« 我们要尊重民意。 »} (Wǒmen yào zūnzhòng mínyì.) Puis : \zh{« 他的反应很敏锐。 »} (Tā de fǎnyìng hěn mǐnruì.)
\textbf{Question :} Quel ton sur la première syllabe de \zh{民意} et de \zh{敏锐} ?
\tcblower
\textbf{Solution détaillée :}
\begin{enumerate}
    \item \zh{民意 mín yì} : \zh{民} se prononce \textbf{2e ton (mín)}. La main monte. Sens : « opinion du peuple ».
    \item \zh{敏锐 mǐn ruì} : \zh{敏} se prononce \textbf{3e ton (mǐn)}. La main descend puis remonte (demi-ton bas en contexte). Sens : « vif, perspicace ».
\end{enumerate}
\emph{Astuce :} Associer \zh{民 mín} (peuple) à \zh{人民 rénmín} (2-2, montées successives) et \zh{敏 mǐn} à \zh{敏感 mǐngǎn} (3-3 $\to$ 2-3, sandhi).
\end{exoresolu}

\courssub{Mots pièges de la leçon — Règle du sandhi du 3e ton}

Certains mots-clés du thème subissent des modifications tonales obligatoires en contexte. La règle d'or : \textbf{l'écriture pinyin garde le ton d'origine (citation), la prononciation change.}

\begin{propriete}[Règle du sandhi du 3e ton — Deux 3e tons qui se suivent]
Lorsque \textbf{deux syllabes de 3e ton} se suivent, \textbf{la première se prononce comme un 2e ton} (montant), mais \textbf{s'écrit toujours avec le 3e ton} (caron ǎ).
\begin{center}
\begin{tabularx}{\linewidth}{|X|X|X|}
\hline
\rowcolor{popgreenL} \textbf{Mot (écrit)} & \textbf{Pinyin écrit} & \textbf{Pinyin prononcé (sandhi)} \\
\hline
选举 & xuǎnjǔ & \textbf{xuánjǔ} (2-3) \\
\hline
很好 & hěn hǎo & \textbf{hén hǎo} (2-3) \\
\hline
马上 & mǎshàng & \textbf{má shàng} (2-4) \\
\hline
\end{tabularx}
\end{center}
\emph{Exception :} Si une pause (virgule, fin de groupe rythmique) sépare les deux syllabes, le premier garde son 3e ton complet.
\end{propriete}

\begin{exemplebox}[Application aux mots du thème]
\begin{itemize}
    \item \zh{民主 mínzhǔ} : 2e + 3e $\to$ \textbf{pas de sandhi} (premier n'est pas 3e). Prononcé \textbf{mínzhǔ}.
    \item \zh{选举 xuǎnjǔ} : 3e + 3e $\to$ \textbf{sandhi obligatoire}. Prononcé \textbf{xuánjǔ} (2-3). \emph{N'écrire jamais « xuánjǔ » dans le vocabulaire.}
    \item \zh{自由 zìyóu} : 4e + 2e $\to$ \textbf{pas de sandhi}. Prononcé \textbf{zìyóu}. Attention au 4e ton sec sur \zh{自}.
    \item \zh{投票 tóupiào} : 2e + 4e $\to$ \textbf{pas de sandhi}. Prononcé \zh{tóupiào}. \zh{投} monte, \zh{票} tombe.
    \item \zh{权利 quánlì} : 2e + 4e $\to$ \textbf{pas de sandhi}. Ne pas confondre avec \zh{全力 quánlì} (mêmes tons, sens différent).
\end{itemize}
\end{exemplebox}

\begin{attention}[Erreur fréquente des francophones — Aplatir les tons]
\begin{itemize}
    \item \textbf{Piège :} Parler chinois sur une mélodie « française » (intonation de phrase seulement), en oubliant le ton lexical de chaque syllabe. Une phrase plate devient incompréhensible.
    \item \textbf{Exemple classique :} \zh{妈 mā} (1er ton, maman) vs \zh{麻 má} (2e, chanvre) vs \zh{马 mǎ} (3e, cheval) vs \zh{骂 mà} (4e, gronder) vs \zh{吗 ma} (neutre, particule).
    \item \textbf{Au thème :} Dire \zh{选举 xuǎn jǔ} (3-3 sans sandhi) sonne « lourd, non naturel ». Dire \zh{民主 mǐn zhǔ} (3-3 au lieu de 2-3) change le sens de \zh{民} en \zh{敏}.
    \item \textbf{Conseil :} \emph{Exagérer} les tons au début (théâtre, chant). Mieux vaut un ton trop marqué qu'un ton absent.
\end{itemize}
\end{attention}

\courssub{Point de grammaire phonétique : Le cas de 不 (bù) devant 同意 (tóngyì)}

Les formules fonctionnelles du débat utilisent massivement \zh{不同意} (ne pas être d'accord). La règle de \zh{不} est un sandhi tonal majeur.

\begin{propriete}[Règle de sandhi de 不 (bù)]
\zh{不} s'écrit toujours 4e ton, mais se prononce :
\begin{itemize}
    \item \textbf{2e ton (bú)} devant un \textbf{4e ton} : \zh{不是 bú shì}, \zh{不对 bú duì}, \zh{不要 bú yào}.
    \item \textbf{4e ton (bù)} devant les \textbf{1er, 2e, 3e tons} et ton neutre : \zh{不高 bù gāo}, \zh{不忙 bù máng}, \zh{不好 bù hǎo}, \zh{不去 bù qù}.
\end{itemize}
\end{propriete}

\begin{exemplebox}[Application : 我不同意 / 我同意]
\begin{itemize}
    \item \zh{同意 tóngyì} : 2e ton + 4e ton.
    \item \zh{不} + \zh{同} (2e ton) $\to$ \zh{不} garde le \textbf{4e ton} : \zh{不同意 bù tóngyì}.
    \item \textbf{Erreur fréquente :} Dire \zh{*bú tóngyì} (2e ton sur \zh{不}) par analogie avec \zh{不是 bú shì}. \textbf{Faux !} \zh{同} est 2e ton, pas 4e.
    \item \zh{我同意 Wǒ tóngyì.} (Je suis d'accord.)
    \item \zh{我不同意 Wǒ bù tóngyì.} (Je ne suis pas d'accord.)
\end{itemize}
\end{exemplebox}

\begin{exoresolu}[Choix du ton de 不]
\textbf{Énoncé :} Déterminer la prononciation réelle (pas l'écriture) de \zh{不} dans les phrases suivantes du débat.
\begin{enumerate}
    \item \zh{这不公平。}
    \item \zh{我不赞成。}
    \item \zh{这不是民主。}
    \item \zh{我们不投票。}
\end{enumerate}
\tcblower
\textbf{Solution détaillée :}
\begin{enumerate}
    \item \zh{公平 gōngpíng} (1-2). \zh{不} devant 1er ton $\to$ \textbf{4e ton (bù)}. \zh{Zhè bù gōngpíng.}
    \item \zh{赞成 zànchéng} (4-2). \zh{不} devant 4e ton \zh{赞} $\to$ \textbf{2e ton (bú)}. \zh{Wǒ bú zànchéng.}
    \item \zh{是 shì} (4e ton). \zh{不} devant 4e ton $\to$ \textbf{2e ton (bú)}. \zh{Zhè bú shì mínzhǔ.}
    \item \zh{投票 tóupiào} (2-4). \zh{不} devant 2e ton \zh{投} $\to$ \textbf{4e ton (bù)}. \zh{Wǒmen bù tóupiào.}
\end{enumerate}
\end{exoresolu}

\courssub{Chant rythmé des formules fonctionnelles — Automatisation prosodique}

Pour fluidifier le débat, nous transformons les formules en \emph{ritournelles} (comptines rythmiques). Le rythme (binaire ou ternaire) « porte » les tons : on tape dans les mains sur les temps forts.

\begin{textebox}[Ritournelles du débat — À chanter / taper dans les mains]
\textbf{Rythme binaire (temps forts sur 1 et 3) :}
\begin{itemize}
    \item \zh{我 同 意。} (Wǒ tóng yì.) [clap] Wǒ [clap] tóng yì.
    \item \zh{我 不 同 意。} (Wǒ bù tóng yì.) [clap] Wǒ bù [clap] tóng yì.
    \item \zh{这 是 民 主。} (Zhè shì mín zhǔ.) [clap] Zhè shì [clap] mín zhǔ.
    \item \zh{请 投 票。} (Qǐng tóu piào.) [clap] Qǐng [clap] tóu piào.
\end{itemize}
\textbf{Rythme ternaire (valse : temps fort sur 1) :}
\begin{itemize}
    \item \zh{我 认 为…} (Wǒ rèn wéi…) [CLAP] Wǒ rèn wéi…
    \item \zh{我 觉 得…} (Wǒ jué de…) [CLAP] Wǒ jué de…
    \item \zh{你 怎 么 看？} (Nǐ zěn me kàn?) [CLAP] Nǐ zěn me kàn?
    \item \zh{有 什 么 意 见？} (Yǒu shén me yì jiàn?) [CLAP] Yǒu shén me yì jiàn?
\end{itemize}
\end{textebox}

\begin{methode}[Enregistrer, comparer, exagérer — Protocole « Miroir vocal »]
\begin{enumerate}
    \item \textbf{Modèle :} L'enseignant (ou audio natif) enregistre la ritournelle sur l'application de la classe / WhatsApp / OnBuch+.
    \item \textbf{Enregistrement élève :} Chaque élève s'enregistre seul (casque micro) en imitant \emph{exactement} la mélodie, en exagérant les tons (théâtraliser).
    \item \textbf{Comparaison visuelle :} Utiliser un visualiseur de spectre/pitch (Praat, Audacity, ou appli « Chinese Tones ») : superposer la courbe élève (rouge) et modèle (bleu). Vérifier : 1er ton plat ? 4e ton chute raide ? 3e ton sandhi réalisé ?
    \item \textbf{Correction ciblée :} Reprendre seulement la syllabe fautive 5 fois de suite avec geste main.
    \item \textbf{Production finale :} Enchaîner 3 formules sans hésitation : \zh{我认为… 但是… 你怎么看？}
\end{enumerate}
\end{methode}

\courssub{Synthèse visuelle : Carte mentale des tons du thème}

\begin{popfigure}
\begin{tikzpicture}[
    mindmap,
    grow cyclic,
    every node/.style=concept,
    concept color=popblue,
    level 1/.append style={level distance=4cm, sibling angle=90},
    level 2/.append style={level distance=3cm, sibling angle=45},
    level 3/.append style={level distance=2.5cm, sibling angle=30},
    text=white,
    font=\small
]
\node[align=center]{Tons\\Thème Démocratie} [clockwise from=90]
    child [concept color=poporange] { node[align=center]{Mots-clés\\Sandhi}
        child [concept color=poporangeL] { node[align=center]{选举\\xuǎnjǔ $\to$ xuánjǔ} }
        child [concept color=poporangeL] { node[align=center]{很好\\hěn hǎo $\to$ hén hǎo} }
        child [concept color=poporangeL] { node[align=center]{民主\\mínzhǔ (pas de sandhi)} }
    }
    child [concept color=popgreen] { node[align=center]{Paires\\Minimales}
        child [concept color=popgreenL] { node {mín (2) vs mǐn (3)} }
        child [concept color=popgreenL] { node {yì (4) vs yī (1)} }
        child [concept color=popgreenL] { node[align=center]{quán (2) vs quán (2)\\(diff. sens)} }
    }
    child [concept color=poppink] { node[align=center]{Formules\\Fonctionnelles}
        child [concept color=poppinkL] { node[align=center]{我不 同意\\bù tóngyì (4-2-4)} }
        child [concept color=poppinkL] { node[align=center]{我 赞成\\wǒ zànchéng (4-2)} }
        child [concept color=poppinkL] { node[align=center]{请 投 票\\qǐng tóupiào (3-2-4)} }
    }
    child [concept color=poppurple] { node[align=center]{Règles\\Clés}
        child [concept color=poppurpleL] { node {3e + 3e $\to$ 2e + 3e} }
        child [concept color=poppurpleL] { node {不 + 4e $\to$ 2e (bú)} }
        child [concept color=poppurpleL] { node {不 + 1,2,3 $\to$ 4e (bù)} }
    };
\end{tikzpicture}
\legende{Carte mentale récapitulative : mots-clés, paires minimales, formules et règles de sandhi.}
\end{popfigure}

\courssub{Exercice de consolidation : « Le débat des tons »}

\begin{exercice}[Expression orale guidée — Prononciation ciblée]
\textbf{Consigne :} En binôme. L'élève A lit la phrase \textbf{écrite en caractères}. L'élève B, \textbf{sans regarder le texte}, doit :
\begin{enumerate}
    \item Identifier le ton de la syllabe soulignée (lever carte 1-2-3-4).
    \item Répéter la phrase complète avec les bons tons et le geste de la main.
    \item Échanger les rôles.
\end{enumerate}
\textbf{Corpus :}
\begin{enumerate}
    \item \zh{我们要 \underline{选举} 班长。} (Cible : \zh{选举} xuánjǔ)
    \item \zh{这件事 \underline{不} 公平。} (Cible : \zh{不} bù)
    \item \zh{我 \underline{不同意} 你的看法。} (Cible : \zh{不} bù)
    \item \zh{每个 \underline{公民} 有权利。} (Cible : \zh{民} mín)
    \item \zh{请 \underline{举手} 表决。} (Cible : \zh{举} jǔ — 3e ton isolé $\to$ demi-ton bas)
    \item \zh{民主 \underline{需要} 宽容。} (Cible : \zh{需} xū — 1er ton plat)
\end{enumerate}
\end{exercice}

\begin{corrige}[Exercice : Le débat des tons]
\begin{enumerate}
    \item \zh{选举} : \textbf{2e + 3e} (sandhi 3e+3e). Geste : main monte (xuán) puis V bas (jǔ).
    \item \zh{不公平} : \zh{公} 1er ton $\to$ \zh{不} = \textbf{4e ton (bù)}. Geste : main frappe fort vers le bas.
    \item \zh{不同意} : \zh{同} 2e ton $\to$ \zh{不} = \textbf{4e ton (bù)}. Geste : main frappe (bù), monte (tóng), frappe (yì).
    \item \zh{公民} : \zh{民} = \textbf{2e ton (mín)}. Geste : main monte. (Pas \zh{敏} mǐn !).
    \item \zh{举手} : \zh{举} 3e ton isolé (fin de groupe rythmique ici) $\to$ \textbf{3e ton complet (V bas)} ou \textbf{demi-3e ton} selon débit. Geste : main descend bas.
    \item \zh{需要} : \zh{需} = \textbf{1er ton (xū)}. Geste : main plate haute.
\end{enumerate}
\end{corrige}

\begin{aretenir}[Check-list « Prêt pour le débat oral »]
Avant de passer aux jeux de rôle (section 6), vérifie que tu sais :
\begin{itemize}
    \item[$\square$] Prononcer \zh{民主 mínzhǔ}, \zh{选举 xuánjǔ}, \zh{自由 zìyóu}, \zh{投票 tóupiào} avec les bons tons (et sandhi).
    \item[$\square$] Distinguer \zh{mín} (2) / \zh{mǐn} (3) et \zh{yì} (4) / \zh{yī} (1) à l'oreille.
    \item[$\square$] Appliquer la règle de \zh{不} : \zh{不是 bú shì} vs \zh{不高 bù gāo} vs \zh{不同意 bù tóngyì}.
    \item[$\square$] Chanter/rythmer les 6 formules fonctionnelles sans hésitation tonale.
    \item[$\square$] Utiliser le geste de la main pour t'auto-corriger en temps réel.
\end{itemize}
\end{aretenir}

\coursec{Jeux de rôle et débat guidé}

Après avoir travaillé la prononciation des tons et les structures pour exprimer son avis, nous passons maintenant à la mise en situation réelle : vous allez \textbf{jouer}, \textbf{débattre} et \textbf{observer} vos camarades. C'est l'étape où tout le vocabulaire, la grammaire et la phonétique des sections précédentes deviennent des outils de communication vivants.

\courssub{Jeu de rôle 1 : Élection du délégué de classe (竞选班长)}

\begin{textebox}[Scène : campagne pour délégué de classe]
候选人：同学们好，我是李明。我想竞选班长。我的主张是：第一，加强班级卫生检查；第二，组织更多体育活动；第三，建立学习互助小组。请大家支持我！
同学甲：李明，你的主张很好。但是，卫生检查谁来负责？
候选人：我们可以轮流值日，每个人都负责。
同学乙：如果同学们不配合怎么办？
候选人：我会和大家沟通，听取意见。
同学丙：好，我投你一票！
\end{textebox}

\begin{center}
\begin{tabularx}{\linewidth}{|X|X|X|X|}
\hline
\rowcolor{popblueL} \textbf{Caractères} & \textbf{Pinyin} & \textbf{Nature} & \textbf{Traduction} \\
\hline
竞选 & jìngxuǎn & verbe & se présenter (élection), briguer un poste \\
班长 & bānzhǎng & nom & délégué de classe, chef de classe \\
主张 & zhǔzhāng & nom/verbe & proposition, programme ; revendiquer, proposer \\
加强 & jiāqiáng & verbe & renforcer, intensifier \\
卫生检查 & wèishēng jiǎnchá & nom & inspection d'hygiène \\
组织 & zǔzhī & verbe & organiser \\
体育活动 & tǐyù huódòng & nom & activités sportives \\
建立 & jiànlì & verbe & établir, créer, fonder \\
互助小组 & hùzhù xiǎozǔ & nom & groupe d'entraide, groupe de soutien mutuel \\
支持 & zhīchí & verbe & soutenir \\
轮流 & lúnliú & verbe & alterner, se relayer \\
值日 & zhírì & verbe/nom & faire le service (de classe), être de service \\
配合 & pèihé & verbe & coopérer, collaborer \\
沟通 & gōutōng & verbe & communiquer, dialoguer \\
听取意见 & tīngqǔ yìjiàn & locution & écouter les avis, recueillir les opinions \\
投...一票 & tóu... yī piào & locution & voter pour..., donner sa voix à... \\
\hline
\end{tabularx}
\end{center}

\begin{propriete}[Structure : présenter son programme avec \cle{第一… 第二… 第三…}]
Pour énumérer des points de programme de façon claire et formelle, on utilise la structure :
\zh{第一，……；第二，……；第三，……。} (Dì yī, ……; dì èr, ……; dì sān, …….)
\begin{itemize}
\item \zh{第一} (premier), \zh{第二} (deuxième), \zh{第三} (troisième) sont suivis d'une virgule chinoise \zh{，}.
\item Chaque point se termine par \zh{；} (point-virgule chinois) sauf le dernier \zh{。}.
\item Cette structure s'utilise à l'oral (exposé, débat) comme à l'écrit (rédaction, affiche).
\end{itemize}
\end{propriete}

\begin{exemplebox}[Variantes pour présenter ses idées]
\begin{itemize}
\item \zh{我有三点建议：第一……第二……第三……} (Wǒ yǒu sān diǎn jiànyì…) — J'ai trois suggestions…
\item \zh{我的主张有三点：一是……二是……三是……} (Wǒ de zhǔzhāng yǒu sān diǎn: yī shì… èr shì… sān shì…) — Mon programme a trois points : 1)… 2)… 3)…
\item \zh{我想从三个方面说：首先……其次……最后……} (Wǒ xiǎng cóng sān gè fāngmiàn shuō: shǒuxiān… qícì… zuìhòu…) — Je veux parler de trois aspects : d'abord… ensuite… enfin…
\end{itemize}
\end{exemplebox}

\begin{methode}[Jouer le candidat : étapes]
\begin{enumerate}
\item \textbf{Saluer et se présenter} : \zh{同学们好，我是……} (Tóngxuémen hǎo, wǒ shì……)
\item \textbf{Annoncer l'objectif} : \zh{我想竞选班长。} (Wǒ xiǎng jìngxuǎn bānzhǎng.)
\item \textbf{Énoncer 3 points précis} avec \zh{第一/第二/第三} ou \zh{首先/其次/最后}.
\item \textbf{Appeler au vote} : \zh{请大家支持我！} (Qǐng dàjiā zhīchí wǒ!)
\item \textbf{Répondre aux questions} calmement : \zh{这是个好问题……} (Zhè shì gè hǎo wèntí……)
\end{enumerate}
\end{methode}

\begin{experience}[Jeu de rôle 1 — En binômes ou groupes de 4-5]
\begin{enumerate}
\item Un élève tire au sort le rôle de \textbf{candidat} (准备 1 minute).
\item Les autres sont \textbf{électeurs} : ils préparent 2 questions chacun (sur la faisabilité, le financement, la méthode…).
\item Déroulement : candidat (1 min) → questions-réponses (3 min) → vote à main levée.
\item Changez de candidat et recommencez.
\end{enumerate}
\end{experience}

\courssub{Jeu de rôle 2 : Interview d'un journaliste sur le vote des jeunes (记者采访：青年投票)}

\begin{textebox}[Interview : l'importance du vote des jeunes]
记者：你好，请问你叫什么名字？今年多大？
学生：我叫王芳，今年十八岁，是高三学生。
记者：王芳同学，你觉得年轻人投票重要吗？
学生：我觉得非常重要。因为投票是我们的权利，也是责任。只有参与，才能改变现状。
记者：那你觉得现在的年轻人积极吗？
学生：有的积极，有的不关心。我觉得学校应该加强公民教育，让大家明白投票的意义。
记者：谢谢你的采访。祝你高考顺利！
学生：谢谢记者姐姐！
\end{textebox}

\begin{center}
\begin{tabularx}{\linewidth}{|X|X|X|X|}
\hline
\rowcolor{popgreenL} \textbf{Caractères} & \textbf{Pinyin} & \textbf{Nature} & \textbf{Traduction} \\
\hline
记者 & jìzhě & nom & journaliste \\
采访 & cǎifǎng & verbe/nom & interviewer, interview \\
青年 & qīngnián & nom & jeune, jeunesse \\
投票 & tóupiào & verbe/nom & voter, vote \\
权利 & quánlì & nom & droit \\
责任 & zérèn & nom & responsabilité, devoir \\
参与 & cānyù & verbe & participer \\
改变 & gǎibiàn & verbe & changer \\
现状 & xiànzhuàng & nom & situation actuelle, statu quo \\
积极 & jījí & adjectif & actif, positif, enthousiaste \\
关心 & guānxīn & verbe & se soucier de, s'intéresser à \\
公民教育 & gōngmín jiàoyù & nom & éducation civique \\
意义 & yìyì & nom & sens, signification \\
高考 & gāokǎo & nom & examen d'entrée à l'université (Gaokao) \\
顺利 & shùnlì & adjectif & qui se passe bien, sans encombre \\
\hline
\end{tabularx}
\end{center}

\begin{propriete}[Structure : \cle{只有……，才……} (seulement si… alors…)]
Pour exprimer une \textbf{condition nécessaire} (le seul moyen d'obtenir un résultat), on utilise :
\zh{只有 + condition，才 + résultat.} (Zhǐyǒu ……, cái ……)
\begin{itemize}
\item \zh{只有} (seulement) place la condition en premier.
\item \zh{才} (alors seulement) introduit le résultat, qui ne se réalise \textbf{que si} la condition est remplie.
\item Négation : \zh{不……也……} (même sans…, aussi…) ou \zh{如果不……，就不……}.
\end{itemize}
\end{propriete}

\begin{exemplebox}[Exemples « 只有……才…… »]
\begin{itemize}
\item \zh{只有努力学习，才能考上大学。} (Zhǐyǒu nǔlì xuéxí, cái néng kǎo shàng dàxué.) — Ce n'est qu'en étudiant dur qu'on peut entrer à l'université.
\item \zh{只有大家都参与，才能改变现状。} (Zhǐyǒu dàjiā dōu cānyù, cái néng gǎibiàn xiànzhuàng.) — Ce n'est que si tout le monde participe qu'on peut changer la situation.
\item \zh{只有尊重别人，才能得到尊重。} (Zhǐyǒu zūnzhòng biérén, cái néng dédào zūnzhòng.) — Ce n'est qu'en respectant les autres qu'on obtient du respect.
\end{itemize}
\end{exemplebox}

\begin{attention}[Piège francophone : confusion « seulement » / « 只有 »]
En français, « seulement » peut être adverbe (\textit{Il seulement mange}) ou conjonction (\textit{Seulement si tu viens}). En chinois :
\begin{itemize}
\item \zh{只} (zhǐ) seul = « seulement » (adverbe) : \zh{我只有一张票。} (J'ai \textbf{seulement} un billet.)
\item \zh{只有} (zhǐyǒu) = « seulement si / ce n'est que » (condition) : \zh{只有你去，我才去。} (\textbf{Seulement si} tu vas, j'irai.)
\item Ne dites pas \zh{*只你去，我才去} ✗ — \zh{只} ne peut pas introduire une proposition conditionnelle seul.
\end{itemize}
\end{attention}

\begin{experience}[Jeu de rôle 2 — En trios : Journaliste / Interrogé / Observateur]
\begin{enumerate}
\item \textbf{Journaliste} : prépare 4 questions ouvertes (avec \zh{怎么/为什么/觉得怎么样}).
\item \textbf{Interrogé} : prépare 3 arguments personnels (droits, devoirs, éducation, avenir…).
\item \textbf{Observateur} : note le vocabulaire utilisé, les tons, les formules de politesse.
\item Durée : 3 minutes. Rotation des rôles.
\end{enumerate}
\end{experience}

\courssub{Débat guidé : « Les élèves devraient-ils participer aux décisions de l'école ? » (学生应该参加学校的决定吗？)}

\begin{textebox}[Lancement du débat par l'animateur]
主持人：今天的题目是：学生应该参加学校的决定吗？请大家发表意见。
\end{textebox}

\begin{methode}[Animer un débat en chinois : phrases-clés de l'animateur (主持人)]
\begin{enumerate}
\item \textbf{Ouvrir} : \zh{今天的题目是：……请大家发表意见。} (Sujet : … Merci à tous de donner votre avis.)
\item \textbf{Donner la parole} : \zh{请……发言。} (Qǐng …… fāyán.) — Je donne la parole à…
\item \textbf{Relancer} : \zh{还有谁想说？} (Hái yǒu shéi xiǎng shuō?) — Qui d'autre veut parler ?
\item \textbf{Équilibrer} : \zh{我们听听反方的意见。} (Wǒmen tīngtīng fǎnfāng de yìjiàn.) — Entendons l'avis du camp opposé.
\item \textbf{Gérer le temps} : \zh{时间快到了，最后一位发言。} (Shíjiān kuài dào le, zuìhòu yī wèi fāyán.) — Le temps file, dernier intervenant.
\item \textbf{Conclure} : \zh{总之，今天的讨论很有意义。谢谢大家。} (Zǒngzhī, jīntiān de tǎolùn hěn yǒu yìyì. Xièxie dàjiā.) — En résumé, le débat d'aujourd'hui était très significatif. Merci à tous.
\end{enumerate}
\end{methode}

\begin{exemplebox}[Formules pour les intervenants (Camp POUR / Camp CONTRE)]
\begin{center}
\begin{tabularx}{\linewidth}{|X|X|X|}
\hline
\rowcolor{poppurpleL} \textbf{Fonction} & \textbf{Chinois (Pinyin)} & \textbf{Français} \\
\hline
Donner son avis & \zh{我觉得……} (Wǒ juéde……) & Je pense que… \\
& \zh{我认为……} (Wǒ rènwéi……) & J'estime que… (plus formel) \\
& \zh{依我看来……} (Yī wǒ kànlái……) & À mon avis… \\
Accorder & \zh{我同意。} (Wǒ tóngyì.) & Je suis d'accord. \\
& \zh{赞成。} (Zànchéng.) & Approuvé / Pour. \\
& \zh{说得好！} (Shuō de hǎo!) & Bien dit ! \\
Désaccorder (poliment) & \zh{我不同意。} (Wǒ bù tóngyì.) & Je ne suis pas d'accord. \\
& \zh{我不太赞同。} (Wǒ bù tài zàntóng.) & Je ne suis pas tout à fait d'accord. \\
& \zh{恐怕……} (Kǒngpà……) & Je crains que… / Il me semble que… (adoucit) \\
Ajouter / Nuancer & \zh{我补充一点。} (Wǒ bǔchōng yì diǎn.) & J'ajoute un point. \\
& \zh{还有，……} (Háiyǒu, ……) & De plus, … / En outre, … \\
Demander la parole & \zh{请让我说完。} (Qǐng ràng wǒ shuō wán.) & Laissez-moi finir, s'il vous plaît. \\
& \zh{对不起，我想说……} (Duìbuqǐ, wǒ xiǎng shuō……) & Excusez-moi, je voudrais dire… \\
Relancer l'autre & \zh{你觉得呢？} (Nǐ juéde ne?) & Et toi, qu'en penses-tu ? \\
& \zh{你怎么看？} (Nǐ zěnme kàn?) & Comment tu vois ça ? \\
\hline
\end{tabularx}
\end{center}
\end{exemplebox}

\begin{propriete}[Structure du débat guidé en deux camps]
\begin{enumerate}
\item \textbf{Préparation (5 min)} : la classe est divisée en deux camps \zh{赞成方} (camp POUR) et \zh{反对方} (camp CONTRE). Chaque camp prépare 3 arguments principaux sur une fiche.
\item \textbf{Tour 1 — Exposés initiaux (2 × 2 min)} : un représentant par camp présente les 3 arguments (\zh{第一……第二……第三……}).
\item \textbf{Tour 2 — Échanges libres (8 min)} : l'animateur donne la parole alternativement. Les élèves utilisent \zh{我觉得…}、\zh{我不同意…}、\zh{我补充一点…}、\zh{你觉得呢？}.
\item \textbf{Tour 3 — Synthèse (2 × 1 min)} : chaque camp résume en une phrase son point fort.
\item \textbf{Conclusion de l'animateur (1 min)} : \zh{总之……} + remerciements.
\end{enumerate}
\end{propriete}

\begin{popfigure}
\begin{tikzpicture}[
  node distance=1.2cm and 1.5cm,
  every node/.style={font=\small, align=center},
  box/.style={draw=popdark, thick, rounded corners, fill=popcream, minimum width=3.2cm, minimum height=1.2cm},
  arrow/.style={-Stealth, thick, popdark},
  label/.style={font=\tiny, text=popmuted, midway, sloped}
]
\node[box, fill=popblueL, align=center] (start){\textbf{主持人 ouvre}\\\zh{今天的题目是……}};
\node[box, fill=popgreenL, right=of start, align=center] (pour){\textbf{Camp POUR}\\\zh{我觉得……}\\\zh{因为……}};
\node[box, fill=poporangeL, right=of pour, align=center] (contre){\textbf{Camp CONTRE}\\\zh{我不同意……}\\\zh{理由是……}};
\node[box, fill=poppinkL, below=of pour, align=center] (relance){\textbf{Relances}\\\zh{你觉得呢？}\\\zh{我补充一点。}};
\node[box, fill=poppurpleL, below=of contre, align=center] (conclu){\textbf{主持人 conclut}\\\zh{总之……谢谢大家。}};

\draw[arrow] (start) -- (pour) node[label, above] {donne la parole};
\draw[arrow] (pour) -- (contre) node[label, above] {répond};
\draw[arrow] (contre) -- (relance) node[label, left] {échange};
\draw[arrow] (relance) -- (pour) node[label, left] {rebond};
\draw[arrow] (relance) -- (contre) node[label, right] {rebond};
\draw[arrow] (pour) -- (conclu) node[label, right] {synthèse};
\draw[arrow] (contre) -- (conclu) node[label, left] {synthèse};
\end{tikzpicture}
\legende{Organigramme du débat guidé : flux de parole entre l'animateur, les deux camps et les relances}
\end{popfigure}

\courssub{Grille d'observation par les pairs (观察表)}

Chaque observateur remplit cette grille pendant le débat (un observateur pour 2-3 débateurs). Elle sert de base au débriefing.

\begin{center}
\begin{tabularx}{\linewidth}{|X|c|c|c|X|}
\hline
\rowcolor{popgoldL} \textbf{Critère} & \textbf{Très bien (3)} & \textbf{Correct (2)} & \textbf{À améliorer (1)} & \textbf{Commentaire / Exemple entendu} \\
\hline
Prononciation des tons (声调) & Tons 1-4 clairs, neutre juste & Quelques hésitations sur 1-2 tons & Tons souvent confus, sens altéré & \\
\hline
Formules de débat (辩论用语) & Utilise 4+ formules variées & Utilise 2-3 formules de base & Peu ou pas de formules, français & \\
\hline
Fluidité / Cohérence (流利度) & Enchaîne sans pause longue & Pauses courtes, se corrige seul & Hésitations longues, phrases coupées & \\
\hline
Respect de la parole (尊重发言) & Attend son tour, \zh{请让我说完} & Interrompt 1 fois poliment & Coupe la parole brutalement & \\
\hline
Vocabulaire thème (主题词汇) & 8+ mots du thème (投票、权利…) & 4-7 mots du thème & < 4 mots, vocabulaire basique & \\
\hline
\end{tabularx}
\end{center}

\courssub{Débriefing collectif : points forts et axes de progrès}

\begin{methode}[Débriefing structuré (5-7 minutes)]
\begin{enumerate}
\item \textbf{Retour observateurs} (2 min) : chaque observateur donne \textbf{un point fort} et \textbf{un axe de progrès} pour son binôme/groupe, en citant une phrase entendue.
\item \textbf{Auto-évaluation} (1 min) : les débateurs disent : \zh{我觉得我……下次我要……} (Je pense que j'ai… La prochaine fois je vais…).
\item \textbf{Synthèse prof} (2 min) : relève 2-3 erreurs récurrentes de tons ou de structure (ex. \zh{*我觉得应该} au lieu de \zh{我觉得应该……因为……}), félicite les bonnes formules (\zh{我补充一点}, \zh{只有……才……}).
\item \textbf{Note collective} : la classe vote pour le « Meilleur orateur » (最佳发言人) et le « Meilleur argument » (最佳论点) — symbolique, valorisant.
\end{enumerate}
\end{methode}

\begin{savaistu}[Collocation « 发表意见 » — star de l'oral du Bac]
\zh{发表意见} (fābiǎo yìjiàn) = « exprimer un avis, prendre la parole pour donner son opinion ».
\begin{itemize}
\item \zh{发表} (fābiǎo) = publier, émettre, présenter (formel) — s'utilise avec \zh{演讲} (discours), \zh{声明} (déclaration), \zh{看法} (point de vue).
\item \zh{意见} (yìjiàn) = avis, opinion, suggestion — aussi \zh{建议} (jiànyì, proposition constructive).
\item \textbf{Au Bac oral} : l'examinateur dit souvent \zh{请发表你的意见} ou \zh{谈谈你的看法}. Répondre par \zh{我觉得……因为……} + \zh{发表意见} montre une maîtrise du registre formel.
\item \textbf{Variante} : \zh{提意见} (tí yìjiàn) = faire une remarque / suggérer (plus concret, moins solennel).
\end{itemize}
\end{savaistu}

\begin{exoresolu}[Mettre en ordre les étapes d'un débat]
\textbf{Énoncé :} Remettez dans l'ordre logique (1 à 6) :
\begin{enumerate}
\item \zh{主持人总结：总之……}
\item \zh{双方自由交锋，使用“补充/不同意/你觉得呢”}
\item \zh{主持人宣布题目，请大家发表意见}
\item \zh{每方推选代表，陈述三个论点（第一/第二/第三）}
\item \zh{每方代表作最后总结（一句）}
\item \zh{分组准备：赞成方 / 反对方，各找3 arguments}
\end{enumerate}
\textbf{Solution détaillée :}
L'ordre correct est : \textbf{6 → 3 → 4 → 2 → 5 → 1}.
\begin{itemize}
\item \textbf{6} (Préparation) : avant tout, les camps préparent leurs arguments.
\item \textbf{3} (Ouverture) : l'animateur lance le sujet.
\item \textbf{4} (Exposés initiaux) : chaque camp présente sa position structurée.
\item \textbf{2} (Échanges libres) : débat vif avec relances.
\item \textbf{5} (Synthèses finales) : chaque camp résume son point fort.
\item \textbf{1} (Conclusion) : l'animateur clôture.
\end{itemize}
\end{exoresolu}

\begin{exercice}[Entraînement : complétez les répliques du débat]
Complétez en chinois (caractères + pinyin) :
\begin{enumerate}
\item L'animateur ouvre : « Le sujet d'aujourd'hui est : les élèves devraient-ils participer aux décisions ? \textbf{\_\_\_\_\_\_\_\_\_\_} » (Merci à tous de donner votre avis.)
\item Camp POUR : « \textbf{\_\_\_\_\_\_\_\_\_\_}，投票是我们的权利。 » (Je pense que… / À mon avis…)
\item Camp CONTRE : « \textbf{\_\_\_\_\_\_\_\_\_\_}，学生还不成熟。 » (Je ne suis pas d'accord / Je ne suis pas tout à fait d'accord.)
\item Un élève veut ajouter : « \textbf{\_\_\_\_\_\_\_\_\_\_}，学校应该教我们怎么投票。 » (J'ajoute un point / De plus,…)
\item L'animateur clôture : « \textbf{\_\_\_\_\_\_\_\_\_\_}，今天的讨论很有意义。 » (En résumé / En conclusion,…)
\end{enumerate}
\end{exercice}

\begin{corrige}[Exercice — Compléter les répliques]
\begin{enumerate}
\item \zh{请大家发表意见。} (Qǐng dàjiā fābiǎo yìjiàn.)
\item \zh{我觉得} (Wǒ juéde) / \zh{我认为} (Wǒ rènwéi) / \zh{依我看来} (Yī wǒ kànlái)
\item \zh{我不同意} (Wǒ bù tóngyì) / \zh{我不太赞同} (Wǒ bù tài zàntóng)
\item \zh{我补充一点} (Wǒ bǔchōng yì diǎn) / \zh{还有} (Háiyǒu)
\item \zh{总之} (Zǒngzhī)
\end{enumerate}
\end{corrige}

\begin{attention}[Derniers pièges à éviter en débat oral]
\begin{itemize}
\item \textbf{Ne pas dire} \zh{*你错了！} (Tu as tort !) — trop brutal. Préférez \zh{我不太赞同，因为……} ou \zh{恐怕不太对。}
\item \textbf{Ne pas couper} la parole sans \zh{对不起，我想说……} ou \zh{请让我说完。}
\item \textbf{Ne pas rester silencieux} : même un \zh{我赞同} ou \zh{说得好} compte comme participation.
\item \textbf{Tons — sandhi du 3\textsuperscript{e} ton} : \zh{只有} (zhǐyǒu \rightarrow \textbf{zhí} yǒu), \zh{主张} (zhǔzhāng), \zh{投票} (tóupiào), \zh{很好} (hěn hǎo \rightarrow \textbf{hén} hǎo). Le 3\textsuperscript{e} ton devient 2\textsuperscript{e} ton devant un autre 3\textsuperscript{e} ton.
\item \textbf{Tons — ton neutre} : \zh{觉得} (jué\textbf{de}), \zh{意见} (yì\textbf{jiàn}), \zh{一点} (yì\textbf{diǎn}) — la 2\textsuperscript{e} syllabe est brève et légère, pas accentuée.
\item \zh{才} (cái) dans \zh{只有……才……} est un 2\textsuperscript{e} ton (montant) — ne le lisez pas en 4\textsuperscript{e} ton.
\end{itemize}
\end{attention}

\courssub{Bilan de la leçon 24 — Expression orale : démocratie}

\begin{aretenir}[Ce que vous devez retenir et savoir faire]
\begin{itemize}
\item \textbf{Vocabulaire clé} : \cle{竞选}、\cle{班长}、\cle{主张}、\cle{投票}、\cle{权利}、\cle{责任}、\cle{参与}、\cle{发表意见}、\cle{赞成/反对}、\cle{补充}、\cle{主持人}.
\item \textbf{Structures maîtrisées} :
  \begin{itemize}
  \item Énumérer : \cle{第一……第二……第三……} / \cle{首先……其次……最后……}
  \item Condition nécessaire : \cle{只有……，才……}
  \item Exprimer avis/accord/désaccord : \cle{我觉得……}、\cle{我同意/不同意}、\cle{我不太赞同}
  \item Relancer : \cle{你觉得呢？}、\cle{我补充一点}
  \item Animer : \cle{请……发言}、\cle{还有谁想说？}、\cle{总之……}
  \end{itemize}
\item \textbf{Compétences orales} : présenter un programme (1 min), répondre en interview, débattre en respectant les tours de parole, prononcer les tons 1-4 + neutre dans des phrases longues.
\item \textbf{Culture} : parallèles Cameroun/Chine sur vote des jeunes, délégués de classe, éducation civique.
\end{itemize}
\end{aretenir}

\begin{experience}[Prolongement hors classe — « Journal de débat »]
Chaque élève tient un petit carnet \zh{辩论日记} (Biànlùn rìjì) pendant le trimestre :
\begin{itemize}
\item Après chaque débat : date, sujet, mon camp, mon argument le plus fort, une formule nouvelle apprise, une correction de ton notée par le prof ou un pair.
\item En fin de trimestre : relecture, auto-évaluation de la progression (fluidité, vocabulaire, tons).
\end{itemize}
\end{experience}

\newpage

\coursec{Activité d'intégration}

\courssub{Objectifs de l'activité}

À l'issue de cette activité d'intégration, tu seras capable de :

\begin{itemize}
\item présenter oralement un avis structuré en chinois sur un sujet de citoyenneté (la démocratie et le vote des jeunes) ;
\item utiliser correctement les connecteurs de l'exposé simple : \zh{首先} \emph{(shǒuxiān)} « d'abord », \zh{然后} \emph{(ránhòu)} « ensuite », \zh{最后} \emph{(zuìhòu)} « enfin », \zh{总之} \emph{(zǒngzhī)} « en résumé » ;
\item donner ton opinion, exprimer l'accord et le désaccord poli, et relancer la parole avec des formules fonctionnelles ;
\item prononcer correctement les tons dans un échange spontané, en particulier les tons des mots clés du débat ;
\item écouter activement tes camarades et évaluer une prise de parole à l'aide d'une grille d'observation.
\end{itemize}

\courssub{Situation}

Dans le cadre de la semaine de la citoyenneté de ton lycée à Yaoundé, le club de chinois organise un « forum des jeunes citoyens ». Le proviseur a proposé un thème qui fait débat dans toute la classe : les jeunes de moins de vingt-cinq ans sont très nombreux au Cameroun, mais beaucoup ne s'intéressent pas aux élections. Le club de chinois a décidé de mener ce débat entièrement en chinois, devant les autres classes et un jury d'enseignants. Un élève joue le rôle de l'animateur (\zh{主持人} \emph{zhǔchírén}), deux camps s'affrontent — le camp « pour » et le camp « contre » — et le public participe en posant des questions. À la fin, toute la classe vote à main levée pour élire le meilleur orateur.

Voici l'annonce affichée au tableau du club et le message d'ouverture préparé par l'animateur :

\begin{textebox}[L'annonce du forum et l'ouverture du débat]
青年论坛：年轻人的投票重要吗？\\
\emph{Qīngnián lùntán : Niánqīngrén de tóupiào zhòngyào ma ?}\\
« Forum des jeunes : le vote des jeunes est-il important ? »\\[4pt]
主持人：大家好！欢迎来到青年论坛。今天的题目是：年轻人的投票重要吗？\\
\emph{Zhǔchírén : Dàjiā hǎo ! Huānyíng láidào qīngnián lùntán. Jīntiān de tímù shì : Niánqīngrén de tóupiào zhòngyào ma ?}\\
« L'animateur : Bonjour à tous ! Bienvenue au forum des jeunes. Le sujet d'aujourd'hui est : le vote des jeunes est-il important ? »\\[4pt]
主持人：有的同学说很重要，有的同学说不重要。你们的看法呢？我们先听一听两方的意见。\\
\emph{Zhǔchírén : Yǒu de tóngxué shuō hěn zhòngyào, yǒu de tóngxué shuō bù zhòngyào. Nǐmen de kànfǎ ne ? Wǒmen xiān tīng yi tīng liǎng fāng de yìjian.}\\
« L'animateur : Certains élèves disent que c'est très important, d'autres disent que non. Et vous, qu'en pensez-vous ? Écoutons d'abord l'avis des deux camps. »
\end{textebox}

\begin{definition}[Mots nouveaux du texte]
\begin{itemize}
\item \zh{论坛} \emph{(lùntán)} : forum ; \zh{青年论坛} « forum des jeunes ».
\item \zh{题目} \emph{(tímù)} : sujet, thème (d'un débat, d'un exposé, d'un devoir).
\item \zh{有的…有的…} \emph{(yǒu de… yǒu de…)} : « certains… d'autres… », structure de répartition très utile à l'oral.
\item \zh{听一听} \emph{(tīng yi tīng)} : « écouter un peu » ; la réduplication du verbe (avec ou sans \zh{一}) adoucit et raccourcit l'action.
\item \zh{两方} \emph{(liǎng fāng)} : « les deux camps, les deux parties » ; \zh{方} est ici un nom (« côté »), mesuré par \zh{两} comme les nombres.
\end{itemize}
\end{definition}

\courssub{Consignes}

\begin{enumerate}
\item \textbf{Préparation du lexique.} Relis le tableau de vocabulaire de la leçon. Souligne dans l'annonce du forum les mots du thème « démocratie » et vérifie leur prononciation (tons) avec ton voisin.
\item \textbf{Compréhension.} Réponds oralement en chinois : quel est le sujet du débat ? Qui parle dans le texte ? Que demande l'animateur à la fin ?
\item \textbf{Formation des camps.} La classe se divise en deux camps : le camp « pour » (\zh{赞成} \emph{zànchéng}) et le camp « contre » (\zh{反对} \emph{fǎnduì}). Chaque camp prépare par écrit un mini-exposé de cinq phrases avec le plan \zh{首先} / \zh{然后} / \zh{最后} / \zh{总之}.
\item \textbf{Mini-exposés.} Un porte-parole de chaque camp présente l'exposé devant la classe, en respectant les tons et en utilisant au moins deux formules fonctionnelles (\zh{我认为…}, \zh{我觉得…}, \zh{我的看法是…}).
\item \textbf{Débat ouvert.} Le public relance les orateurs avec les questions apprises : \zh{你觉得呢？} \emph{(Nǐ juéde ne ?)} « Et toi, qu'en penses-tu ? », \zh{你能举个例子吗？} \emph{(Nǐ néng jǔ ge lìzi ma ?)} « Peux-tu donner un exemple ? ». Chaque intervenant réagit avec une formule d'accord (\zh{我同意你的看法。}) ou de désaccord poli (\zh{对不起，我不太同意。}).
\item \textbf{Observation.} Les élèves qui ne parlent pas remplissent la grille d'observation : tons corrects, formules fonctionnelles utilisées, fluidité, clarté des arguments.
\item \textbf{Conclusion et vote.} L'animateur conclut avec \zh{总之…}, puis la classe élit le meilleur orateur par un vote symbolique à main levée : \zh{我们投票吧！} \emph{(Wǒmen tóupiào ba !)} « Votons ! ».
\end{enumerate}

\courssub{Ressources à mobiliser}

\textbf{Lexique du thème.}

\begin{center}
\begin{tabularx}{\linewidth}{|X|X|X|X|}
\hline
\rowcolor{popblueL} Caractères & Pinyin & Nature & Traduction \\
\hline
民主 & mínzhǔ & nom / adj. & démocratie, démocratique \\
\hline
投票 & tóupiào & verbe / nom & voter, vote \\
\hline
选举 & xuǎnjǔ & verbe / nom & élire, élection \\
\hline
意见 & yìjian & nom & opinion, avis \\
\hline
看法 & kànfǎ & nom & point de vue \\
\hline
重要 & zhòngyào & adjectif & important \\
\hline
赞成 & zànchéng & verbe & approuver, être pour \\
\hline
反对 & fǎnduì & verbe & s'opposer, être contre \\
\hline
例子 & lìzi & nom & exemple \\
\hline
公民 & gōngmín & nom & citoyen \\
\hline
权利 & quánlì & nom & droit \\
\hline
责任 & zérèn & nom & responsabilité \\
\hline
表达 & biǎodá & verbe & exprimer \\
\hline
参加 & cānjiā & verbe & participer à \\
\hline
主持人 & zhǔchírén & nom & animateur, présentateur \\
\hline
\end{tabularx}
\end{center}

\textbf{Formules fonctionnelles.}

\begin{center}
\begin{tabularx}{\linewidth}{|X|X|X|}
\hline
\rowcolor{popblueL} Fonction & Formule en caractères + pinyin & Traduction \\
\hline
Donner son avis & 我认为年轻人的投票很重要。Wǒ rènwéi niánqīngrén de tóupiào hěn zhòngyào. & Je pense que le vote des jeunes est très important. \\
\hline
Donner son avis (plus doux) & 我觉得每个人都应该投票。Wǒ juéde měi ge rén dōu yīnggāi tóupiào. & Je trouve que tout le monde devrait voter. \\
\hline
Accord & 我同意你的看法。Wǒ tóngyì nǐ de kànfǎ. & Je suis d'accord avec ton point de vue. \\
\hline
Désaccord poli & 对不起，我不太同意。Duìbuqǐ, wǒ bú tài tóngyì. & Pardon, je ne suis pas tout à fait d'accord. \\
\hline
Relancer & 你觉得呢？Nǐ juéde ne ? & Et toi, qu'en penses-tu ? \\
\hline
Demander un exemple & 你能举个例子吗？Nǐ néng jǔ ge lìzi ma ? & Peux-tu donner un exemple ? \\
\hline
Conclure & 总之，我们应该参加选举。Zǒngzhī, wǒmen yīnggāi cānjiā xuǎnjǔ. & En résumé, nous devrions participer aux élections. \\
\hline
\end{tabularx}
\end{center}

\textbf{Plan de l'exposé simple.} Rappelle-toi la structure en quatre étapes : \zh{首先} (d'abord : annoncer sa position) → \zh{然后} (ensuite : premier argument) → \zh{最后} (enfin : deuxième argument ou exemple) → \zh{总之} (en résumé : conclusion).

\begin{propriete}[La structure hypothétique 如果…就…]
Pour argumenter, on utilise souvent l'hypothèse : \cle{Sujet + 如果 + hypothèse, (Sujet) + 就 + conséquence} « si… alors… ».\\
\zh{如果年轻人不投票，政府就听不到我们的声音。}\\
\emph{Rúguǒ niánqīngrén bù tóupiào, zhèngfǔ jiù tīng bu dào wǒmen de shēngyīn.}\\
« Si les jeunes ne votent pas, le gouvernement ne pourra pas entendre notre voix. »\\
Remarque : \zh{就} se place juste avant le verbe de la conséquence ; \zh{听不到} \emph{(tīng bu dào)} « ne pas pouvoir entendre » est un complément de résultat au négatif.
\end{propriete}

\begin{attention}[Pièges de prononciation et de grammaire]
\begin{itemize}
\item \zh{重要} \emph{zhòngyào} : quatrième ton + quatrième ton, ne pas dire \emph{zhōngyào}. Attention aussi à \zh{选举} \emph{xuǎnjǔ} : deux troisièmes tons, le premier devient un deuxième ton à l'oral (\emph{xuánjǔ}).
\item \zh{觉得} \emph{juéde} : le deuxième caractère est au ton neutre. Même chose pour \zh{看法} \emph{(kànfǎ)}, où \zh{法} garde son troisième ton, mais \zh{例子} \emph{(lìzi)} a un \zh{子} au ton neutre.
\item \zh{不太同意} : \zh{不} devient deuxième ton devant le quatrième ton de \zh{太} : \emph{bú tài tóngyì}.
\item Ne dis pas « \zh{我同意你 »} seul : complète toujours (\zh{我同意你的看法}).
\item L'ordre des mots reste Sujet + Verbe : « \zh{我认为投票很重要 »}, jamais « \zh{我认为很重要投票 »}.
\item \zh{参加} se construit directement avec son complément : \zh{参加选举} « participer aux élections », sans préposition (pas de « \zh{参加到选举 »}).
\end{itemize}
\end{attention}

\courssub{Proposition de production et corrigé commenté}

\begin{exoresolu}[Mini-exposé modèle du camp « pour »]
Présente en cinq phrases l'avis du camp « pour » en suivant le plan 首先 / 然后 / 最后 / 总之.
\tcblower
\textbf{Production modèle :}\\[4pt]
大家好！我认为年轻人的投票非常重要。\\
\emph{Dàjiā hǎo ! Wǒ rènwéi niánqīngrén de tóupiào fēicháng zhòngyào.}\\
« Bonjour à tous ! Je pense que le vote des jeunes est très important. »\\[3pt]
首先，年轻人是国家未来的主人，我们应该表达自己的意见。\\
\emph{Shǒuxiān, niánqīngrén shì guójiā wèilái de zhǔrén, wǒmen yīnggāi biǎodá zìjǐ de yìjian.}\\
« D'abord, les jeunes sont les maîtres de l'avenir du pays, nous devons exprimer nos opinions. »\\[3pt]
然后，投票是每个公民的权利，也是责任。\\
\emph{Ránhòu, tóupiào shì měi ge gōngmín de quánlì, yě shì zérèn.}\\
« Ensuite, voter est un droit de chaque citoyen, et aussi une responsabilité. »\\[3pt]
最后，如果年轻人不投票，政府就听不到我们的声音。\\
\emph{Zuìhòu, rúguǒ niánqīngrén bù tóupiào, zhèngfǔ jiù tīng bu dào wǒmen de shēngyīn.}\\
« Enfin, si les jeunes ne votent pas, le gouvernement ne pourra pas entendre notre voix. »\\[3pt]
总之，我赞成年轻人积极参加选举。谢谢大家！\\
\emph{Zǒngzhī, wǒ zànchéng niánqīngrén jījí cānjiā xuǎnjǔ. Xièxie dàjiā !}\\
« En résumé, je suis pour que les jeunes participent activement aux élections. Merci à tous ! »\\[6pt]
\textbf{Commentaire des choix de langue :}
\begin{itemize}
\item L'orateur ouvre avec la formule d'avis \zh{我认为…} et annonce clairement sa position dès la première phrase : c'est la marque d'un exposé structuré.
\item Les connecteurs \zh{首先} / \zh{然后} / \zh{最后} / \zh{总之} scandent les quatre étapes du plan appris dans la leçon ; ils rendent le propos facile à suivre pour l'auditoire.
\item La phrase avec \zh{如果…就…} (\emph{rúguǒ… jiù…}, « si… alors… ») est une structure hypothétique de niveau Terminale qui donne de la force à l'argumentation ; le complément de résultat \zh{听不到} précise « ne pas pouvoir entendre ».
\item Le mot \zh{公民} \emph{(gōngmín)} et la paire \zh{权利} / \zh{责任} (« droit » / « responsabilité ») enrichissent le lexique de la citoyenneté.
\item Attention aux tons : \zh{非常重要} \emph{fēicháng zhòngyào} enchaîne premier ton, deuxième ton, puis deux quatrièmes tons ; \zh{选举} se prononce \emph{xuánjǔ} par assimilation des deux troisièmes tons.
\end{itemize}
\end{exoresolu}

\begin{methode}[Grille d'observation de l'orateur]
Pendant que tes camarades parlent, note un point pour chaque critère réussi :
\begin{enumerate}
\item \textbf{Tons} : les mots clés (\zh{投票}, \zh{选举}, \zh{重要}, \zh{赞成}, \zh{反对}) sont-ils prononcés avec les bons tons ?
\item \textbf{Formules} : l'orateur utilise-t-il au moins deux formules fonctionnelles (\zh{我认为…}, \zh{我同意…}, \zh{总之…}) ?
\item \textbf{Plan} : entend-on les quatre connecteurs \zh{首先} / \zh{然后} / \zh{最后} / \zh{总之} ?
\item \textbf{Arguments} : y a-t-il au moins deux arguments et un exemple (\zh{例子}) ?
\item \textbf{Fluidité} : le débit est-il naturel, sans longues hésitations ?
\end{enumerate}
Additionne les points sur cinq : cela t'aide à voter pour le meilleur orateur de façon objective.
\end{methode}

\begin{savaistu}[Vote et participation : Cameroun et Chine]
Au Cameroun comme en Chine, les jeunes sont invités à s'intéresser à la vie citoyenne, mais les formes de participation diffèrent. En Chine, les élections des représentants locaux (\zh{选举} \emph{xuǎnjǔ}) concernent d'abord les assemblées populaires de base, et les jeunes débattent beaucoup dans les clubs et les « forums » (\zh{论坛} \emph{lùntán}) des écoles et des universités. Au Cameroun, les débats citoyens dans les lycées, comme ton « forum des jeunes », sont une excellente école de la prise de parole : dire son avis (\zh{意见}), écouter celui des autres et voter (\zh{投票}) à main levée, c'est déjà pratiquer la démocratie (\zh{民主}) au quotidien.
\end{savaistu}

\courssub{Bilan de l'activité}

Cette activité t'a permis de réunir toutes les compétences de la leçon dans une situation de communication authentique et proche de ta vie de lycéen camerounais :

\begin{itemize}
\item \textbf{Compréhension} : tu as lu et compris une annonce et un message d'ouverture de débat en chinois.
\item \textbf{Lexique} : tu as réutilisé le vocabulaire de la démocratie (\zh{民主}, \zh{投票}, \zh{选举}, \zh{赞成}, \zh{反对}) dans un contexte réel.
\item \textbf{Grammaire et structures} : tu as mobilisé les formules fonctionnelles pour donner ton avis, marquer l'accord ou le désaccord poli, relancer la parole et conclure, ainsi que le plan d'exposé \zh{首先} / \zh{然后} / \zh{最后} / \zh{总之} et l'hypothèse \zh{如果…就…}.
\item \textbf{Production orale} : tu as présenté un mini-exposé structuré et participé à un échange spontané, en soignant tes tons.
\item \textbf{Écoute active et évaluation} : grâce à la grille d'observation et au vote final \zh{我们投票吧！}, tu as appris à écouter, à évaluer une prise de parole et à vivre, en chinois, un petit moment de démocratie en classe.
\end{itemize}

Pour aller plus loin, tu peux rejouer ce débat à la maison avec un autre sujet citoyen, par exemple : \zh{学生应该带手机来学校吗？} \emph{(Xuésheng yīnggāi dài shǒujī lái xuéxiào ma ?)} « Les élèves devraient-ils apporter leur téléphone à l'école ? ».

\newpage

\coursec{Méthodes et exercices résolus}

\courssub{Mise en route : le vocabulaire du thème « démocratie »}

Avant de prendre la parole, il faut disposer d'un lexique précis et correctement prononcé. Voici le vocabulaire indispensable du débat sur la démocratie, à apprendre avec les tons.

\begin{center}
\begin{tabularx}{\linewidth}{|X|X|X|X|}
\hline
\rowcolor{popblueL} Caractères & Pinyin & Nature & Traduction \\
\hline
民主 & mínzhǔ & nom / adj. & démocratie ; démocratique \\
\hline
选举 & xuǎnjǔ & nom / verbe & élection ; élire \\
\hline
投票 & tóupiào & verbe & voter \\
\hline
权利 & quánlì & nom & droit \\
\hline
自由 & zìyóu & nom / adj. & liberté ; libre \\
\hline
意见 & yìjiàn & nom & opinion, avis \\
\hline
看法 & kànfǎ & nom & point de vue \\
\hline
同意 & tóngyì & verbe & être d'accord \\
\hline
表达 & biǎodá & verbe & exprimer \\
\hline
参加 & cānjiā & verbe & participer à \\
\hline
参与 & cānyù & verbe & prendre part à \\
\hline
制度 & zhìdù & nom & système, régime \\
\hline
政治 & zhèngzhì & nom & politique \\
\hline
公民 & gōngmín & nom & citoyen \\
\hline
责任 & zérèn & nom & responsabilité, devoir \\
\hline
发展 & fāzhǎn & nom / verbe & développement ; se développer \\
\hline
解决 & jiějué & verbe & résoudre \\
\hline
管理 & guǎnlǐ & verbe / nom & gérer ; gestion \\
\hline
\end{tabularx}
\end{center}

\begin{attention}[Prononciation : les pièges de ce lexique]
\zh{选举} \emph{xuǎnjǔ} porte deux 3e tons : à l'oral, le premier devient un 2e ton, on prononce \emph{xuánjǔ}. \zh{民主} \emph{mínzhǔ} se lit 2e + 3e ton, sans aucun changement. Ne confondez pas \zh{权利} \emph{quánlì} « droit » et \zh{权力} \emph{quánlì} « pouvoir » : mêmes sons, sens différents — au Bac, précisez le caractère à l'écrit.
\end{attention}

\begin{textebox}[Dialogue modèle : parler de la démocratie]
A：你觉得民主是什么？\\
\emph{Nǐ juéde mínzhǔ shì shénme?}\\
« Qu'est-ce que la démocratie, à ton avis ? »

B：我认为民主就是人民可以选举自己的领导人。\\
\emph{Wǒ rènwéi mínzhǔ jiùshì rénmín kěyǐ xuǎnjǔ zìjǐ de lǐngdǎorén.}\\
« Je pense que la démocratie, c'est quand le peuple peut élire ses propres dirigeants. »

A：我同意你的看法。民主也意味着言论自由。\\
\emph{Wǒ tóngyì nǐ de kànfǎ. Mínzhǔ yě yìwèizhe yánlùn zìyóu.}\\
« Je suis d'accord avec toi. La démocratie signifie aussi la liberté d'expression. »

B：对！那你觉得年轻人应该关心政治吗？\\
\emph{Duì! Nà nǐ juéde niánqīngrén yīnggāi guānxīn zhèngzhì ma?}\\
« C'est ça ! Et penses-tu que les jeunes devraient s'intéresser à la politique ? »

A：当然应该，因为政治决定我们的未来。\\
\emph{Dāngrán yīnggāi, yīnwèi zhèngzhì juédìng wǒmen de wèilái.}\\
« Bien sûr qu'ils le devraient, parce que la politique décide de notre avenir. »
\end{textebox}

\courssub{Fiche méthode 1 : Présenter oralement un exposé simple en chinois}

\begin{methode}[Structurer un exposé oral en trois temps]
Un exposé au Bac se construit toujours en trois parties annoncées à voix haute.
\begin{enumerate}
\item \cle{Introduction} : saluer et annoncer le sujet avec la formule \zh{大家好！今天我要谈一谈……} (\emph{Dàjiā hǎo! Jīntiān wǒ yào tán yi tán...}) « Bonjour à tous ! Aujourd'hui, je vais parler de... ».
\item \cle{Développement} : deux ou trois idées, chacune introduite par un connecteur : \zh{首先} (\emph{shǒuxiān}) « d'abord », \zh{其次} (\emph{qícì}) « ensuite », \zh{最后} (\emph{zuìhòu}) « enfin ». Une idée = une phrase courte Sujet + Verbe + Complément.
\item \cle{Conclusion} : donner son avis personnel avec \zh{我认为……} (\emph{wǒ rènwéi...}) « je pense que... » puis remercier : \zh{谢谢大家！} (\emph{Xièxie dàjiā!}) « Merci à tous ! ».
\end{enumerate}
Réflexes : préparer une fiche de 5 à 8 mots-clés (jamais un texte lu), parler lentement, soigner les tons, regarder l'auditoire.
\end{methode}

\begin{exoresolu}[Exposé : « La démocratie dans mon pays »]
\textbf{Énoncé.} Préparez un exposé oral de 6 à 8 phrases sur le thème « 民主 » (\emph{mínzhǔ}, la démocratie) en suivant le plan : introduction, deux arguments, avis personnel, remerciement.
\tcblower
\textbf{Raisonnement.} On applique le plan en trois temps. Pour chaque argument, on utilise une structure simple : Sujet + \zh{可以} (\emph{kěyǐ}, pouvoir) + Verbe, ou Sujet + \zh{很} + Adjectif. On évite les phrases longues calquées du français.

\textbf{Réponse rédigée (modèle).}

大家好！今天我要谈一谈民主。\\
\emph{Dàjiā hǎo! Jīntiān wǒ yào tán yi tán mínzhǔ.}\\
« Bonjour à tous ! Aujourd'hui, je vais parler de la démocratie. »

首先，在民主的国家，人民可以选举领导人。\\
\emph{Shǒuxiān, zài mínzhǔ de guójiā, rénmín kěyǐ xuǎnjǔ lǐngdǎorén.}\\
« D'abord, dans un pays démocratique, le peuple peut élire ses dirigeants. »

其次，每个人都可以自由地表达自己的想法。\\
\emph{Qícì, měi ge rén dōu kěyǐ zìyóu de biǎodá zìjǐ de xiǎngfǎ.}\\
« Ensuite, chacun peut exprimer librement ses idées. »

我认为民主很重要，因为它让人民参与国家的生活。\\
\emph{Wǒ rènwéi mínzhǔ hěn zhòngyào, yīnwèi tā ràng rénmín cānyù guójiā de shēnghuó.}\\
« Je pense que la démocratie est très importante, parce qu'elle permet au peuple de participer à la vie du pays. »

谢谢大家！\\
\emph{Xièxie dàjiā!} « Merci à tous ! »

\textbf{Justifications grammaticales.} \zh{在} + lieu se place avant le verbe ; \zh{可以} précède le verbe principal ; \zh{让} + personne + verbe exprime « faire/permettre que » ; \zh{因为} introduit la cause, comme en français ; \zh{地} (particule adverbiale) relie l'adverbe \zh{自由} au verbe \zh{表达}.

\textbf{Conclusion.} Un exposé réussi = un plan annoncé, des phrases courtes et correctes, des tons soignés.
\end{exoresolu}

\courssub{Fiche méthode 2 : Donner son avis, être d'accord ou non, relancer}

\begin{methode}[Les formules fonctionnelles du débat]
\begin{enumerate}
\item \cle{Donner son avis} : \zh{我觉得……} (\emph{wǒ juéde...}) « je trouve que... » ; \zh{我认为……} « je pense que... » ; \zh{在我看来，……} (\emph{zài wǒ kànlái}) « à mon avis... ».
\item \cle{Être d'accord} : \zh{我同意你的看法。} (\emph{Wǒ tóngyì nǐ de kànfǎ.}) « Je suis d'accord avec ton point de vue. » ; \zh{你说得对！} (\emph{Nǐ shuō de duì!}) « Tu as raison ! ».
\item \cle{N'être pas d'accord} (toujours avec politesse) : \zh{我不太同意，因为……} (\emph{Wǒ bú tài tóngyì, yīnwèi...}) « Je ne suis pas tout à fait d'accord, parce que... ».
\item \cle{Relancer le dialogue} : \zh{你觉得呢？} (\emph{Nǐ juéde ne?}) « Et toi, qu'en penses-tu ? » ; \zh{你的看法呢？} (\emph{Nǐ de kànfǎ ne?}) « Quel est ton avis ? ».
\end{enumerate}
Réflexe : un désaccord se formule en deux temps — d'abord une concession (\zh{你说得有道理，可是……} « ce que tu dis est sensé, mais... »), puis l'argument.
\end{methode}

\begin{exoresolu}[Dialogue : compléter un débat]
\textbf{Énoncé.} Complétez ce dialogue entre deux élèves de Terminale avec les formules étudiées, puis traduisez-le.

A：我觉得民主对国家的发展很重要。 (\emph{Wǒ juéde mínzhǔ duì guójiā de fāzhǎn hěn zhòngyào.})

B：\_\_\_\_\_\_\_\_\_\_\_\_（表达同意），因为人民可以参加选举。

A：可是也有人说，民主不能解决所有问题。\_\_\_\_\_\_\_\_\_\_\_\_（问B的看法）？

B：我不太同意。我认为民主虽然不是完美的，可是它是最好的制度。
\tcblower
\textbf{Raisonnement.} Le premier blanc suit une opinion : B marque son accord, ce que confirme \zh{因为} qui justifie. On choisit donc une formule d'accord. Le deuxième blanc précède un point d'interrogation et demande l'avis de B : on utilise une formule de relance avec la particule \zh{呢}.

\textbf{Réponse rédigée.}

B：我同意你的看法，因为人民可以参加选举。\\
\emph{Wǒ tóngyì nǐ de kànfǎ, yīnwèi rénmín kěyǐ cānjiā xuǎnjǔ.}\\
« Je suis d'accord avec toi, parce que le peuple peut participer aux élections. »

A：可是也有人说，民主不能解决所有问题。你的看法呢？\\
\emph{Kěshì yě yǒu rén shuō, mínzhǔ bù néng jiějué suǒyǒu de wèntí. Nǐ de kànfǎ ne?}\\
« Mais certains disent aussi que la démocratie ne peut pas résoudre tous les problèmes. Quel est ton avis ? »

B：我不太同意。我认为民主虽然不是完美的，可是它是最好的制度。\\
\emph{Wǒ bú tài tóngyì. Wǒ rènwéi mínzhǔ suīrán bú shì wánměi de, kěshì tā shì zuì hǎo de zhìdù.}\\
« Je ne suis pas tout à fait d'accord. Je pense que la démocratie, bien qu'elle ne soit pas parfaite, est le meilleur système. »

\textbf{Justifications grammaticales.} \zh{对……很重要} : la structure \zh{对 + nom + adjectif} signifie « être + adjectif + pour » ; \zh{虽然……可是……} traduit « bien que... mais... » (les deux moitiés s'emploient ensemble en chinois) ; \zh{最 + adjectif + 的 + nom} forme le superlatif.

\textbf{Conclusion.} Dans un débat, chaque prise de parole = une formule fonctionnelle + un argument introduit par \zh{因为}.
\end{exoresolu}

\courssub{Fiche méthode 3 : Échanger avec ses camarades dans un petit débat}

\begin{methode}[Conduire un échange oral en quatre étapes]
\begin{enumerate}
\item \cle{Écouter} le camarade jusqu'au bout ; repérer son idée principale (souvent après \zh{我觉得} ou \zh{我认为}).
\item \cle{Réagir} d'abord avec une formule courte : \zh{对！} « Oui ! », \zh{真的吗？} (\emph{Zhēn de ma?}) « Vraiment ? », \zh{有意思！} (\emph{Yǒu yìsi!}) « Intéressant ! ».
\item \cle{Répondre} en donnant son propre avis + un argument : formule d'accord ou de désaccord poli + \zh{因为}.
\item \cle{Relancer} pour que l'échange continue : \zh{那你呢？} (\emph{Nà nǐ ne?}) « Et toi alors ? », \zh{还有别的想法吗？} (\emph{Hái yǒu bié de xiǎngfǎ ma?}) « Y a-t-il d'autres idées ? ».
\end{enumerate}
Réflexes : ne jamais couper la parole ; si un mot manque, contourner avec des mots connus (\zh{就是……的那个东西} « c'est... le truc qui... ») ; garder des phrases courtes.
\end{methode}

\begin{exoresolu}[Jeu de rôle : le débat en classe]
\textbf{Énoncé.} Sujet : \zh{学生应该不应该参加学校的管理？} (\emph{Xuésheng yīnggāi bù yīnggāi cānjiā xuéxiào de guǎnlǐ?}) « Les élèves devraient-ils participer à la gestion de l'école ? » Rédigez un échange de 4 répliques (2 par élève) avec au moins un accord, un désaccord poli et une relance.
\tcblower
\textbf{Raisonnement.} On distribue les rôles : A est pour, B exprime d'abord un accord partiel puis une réserve. On vérifie la présence des trois éléments exigés. La question \zh{应该不应该} est la forme affirmative-négative de \zh{应该} (\emph{yīnggāi}, devoir).

\textbf{Réponse rédigée (modèle).}

A：我觉得学生应该参加学校的管理，因为我们最了解学生的生活。\\
\emph{Wǒ juéde xuésheng yīnggāi cānjiā xuéxiào de guǎnlǐ, yīnwèi wǒmen zuì liǎojiě xuésheng de shēnghuó.}\\
« Je trouve que les élèves devraient participer à la gestion de l'école, parce que nous connaissons le mieux la vie des élèves. »

B：你说得有道理，我同意。那你觉得学生可以管理什么呢？\\
\emph{Nǐ shuō de yǒu dàolǐ, wǒ tóngyì. Nà nǐ juéde xuésheng kěyǐ guǎnlǐ shénme ne?}\\
« Ce que tu dis est sensé, je suis d'accord. Alors, que pourraient gérer les élèves, à ton avis ? »

A：比如食堂和体育活动。\\
\emph{Bǐrú shítáng hé tǐyù huódòng.}\\
« Par exemple la cantine et les activités sportives. »

B：我不太同意管理食堂，可是体育活动，可以！还有别的想法吗？\\
\emph{Wǒ bú tài tóngyì guǎnlǐ shítáng, kěshì tǐyù huódòng, kěyǐ! Hái yǒu bié de xiǎngfǎ ma?}\\
« Je ne suis pas très d'accord pour gérer la cantine, mais pour les activités sportives, d'accord ! Y a-t-il d'autres idées ? »

\textbf{Justifications grammaticales.} \zh{最 + adjectif} = superlatif ; \zh{比如} (\emph{bǐrú}) introduit un exemple ; \zh{和} relie deux noms (jamais deux phrases) ; la relance \zh{还有……吗？} s'appuie sur \zh{有} + \zh{吗}.

\textbf{Conclusion.} Un bon échange alterne réaction, argument et relance : personne ne monopolise la parole.
\end{exoresolu}

\courssub{Fiche méthode 4 : Prononcer correctement les tons}

\begin{methode}[Travailler les tons du lexique du débat]
\begin{enumerate}
\item \cle{Identifier} le ton de chaque syllabe du vocabulaire clé avant de parler : \zh{民主} \emph{mínzhǔ} (2+3), \zh{选举} \emph{xuǎnjǔ} (3+3), \zh{意见} \emph{yìjiàn} (4+4), \zh{同意} \emph{tóngyì} (2+4).
\item \cle{Appliquer la règle du 3e ton} : deux 3e tons qui se suivent → le premier devient un 2e ton : \zh{选举} se prononce \emph{xuánjǔ}, \zh{很好} \emph{hěn hǎo} se prononce \emph{hén hǎo}.
\item \cle{Soigner les paires à contraste} : 2e ton (montant, comme un point d'interrogation) contre 4e ton (descendant, comme un ordre) : \zh{民} \emph{mín} / \zh{意见} \emph{yìjiàn}.
\item \cle{Enchaîner} : lire le dialogue modèle à voix haute trois fois — lentement, normalement, puis avec l'intonation de la conviction (le débat exige un ton affirmé).
\end{enumerate}
Réflexe : enregistrer sa voix et comparer avec le modèle du professeur ; un ton faux peut changer le sens (\zh{问} \emph{wèn} « demander » / \zh{文} \emph{wén} « écriture »).
\end{methode}

\begin{exoresolu}[Dictée de tons et lecture]
\textbf{Énoncé.} 1) Indiquez les tons de chaque syllabe : 民主、自由、投票、权利。 2) Quelle syllabe change de ton à l'oral dans 选举 et pourquoi ? 3) Lisez à voix haute : \zh{我认为每个人都应该有投票的权利。}
\tcblower
\textbf{Raisonnement et réponses.}

1) 民主 \emph{mínzhǔ} : 2e + 3e ton ; 自由 \emph{zìyóu} : 4e + 2e ton ; 投票 \emph{tóupiào} : 2e + 4e ton ; 权利 \emph{quánlì} : 2e + 4e ton.

2) Dans \zh{选举} \emph{xuǎnjǔ}, les deux syllabes portent le 3e ton ; à l'oral, la première syllabe devient un 2e ton : on prononce \emph{xuánjǔ}. C'est la règle obligatoire de sandhi du 3e ton.

3) Lecture : \emph{Wǒ rènwéi měi ge rén dōu yīnggāi yǒu tóupiào de quánlì.} « Je pense que chacun devrait avoir le droit de vote. » Points de vigilance : \zh{认为} \emph{rènwéi} (4+2), \zh{每} \emph{měi} (3e ton plein devant le classificateur \zh{个} au ton neutre), \zh{应该} \emph{yīnggāi} (1+1, deux tons hauts et plats).

\textbf{Justifications.} Le ton neutre de \zh{个} \emph{ge} ne s'écrit avec aucune marque ; \zh{都} \emph{dōu} « tous » se place toujours avant le verbe et après le sujet auquel il se rapporte.

\textbf{Conclusion.} Au Bac, une prononciation des tons claire et stable vaut autant qu'un vocabulaire riche : entraînez-vous chaque jour à voix haute.
\end{exoresolu}

\courssub{Fiche méthode 5 : Réagir et argumenter avec les bonnes structures}

\begin{methode}[Les structures grammaticales de l'argumentation]
\begin{enumerate}
\item \cle{La cause} : \zh{因为 + cause，所以 + conséquence} (\emph{yīnwèi... suǒyǐ...}) « parce que... donc... ». En chinois, on peut utiliser les deux ensemble, contrairement au français.
\item \cle{La concession} : \zh{虽然 + concession，可是/但是 + argument} (\emph{suīrán... kěshì/dànshì...}) « bien que... mais... ».
\item \cle{L'obligation et la possibilité} : \zh{应该 + verbe} « devoir » ; \zh{可以 + verbe} « pouvoir (autorisation) » ; \zh{能 + verbe} « pouvoir (capacité) ».
\item \cle{Le superlatif} : \zh{最 + adjectif} « le plus... » : \zh{最重要的问题} (\emph{zuì zhòngyào de wèntí}) « le problème le plus important ».
\item \cle{L'exemple} : \zh{比如 + nom} ou \zh{例如 + phrase} « par exemple ».
\end{enumerate}
Réflexe : un argument solide = opinion (\zh{我认为}) + justification (\zh{因为}) + exemple (\zh{比如}).
\end{methode}

\begin{exoresolu}[Composition guidée : un paragraphe d'argumentation]
\textbf{Énoncé.} Rédigez un paragraphe de 4 à 5 phrases répondant à la question : \zh{年轻人应该关心政治吗？} (\emph{Niánqīngrén yīnggāi guānxīn zhèngzhì ma?}) « Les jeunes devraient-ils s'intéresser à la politique ? » Utilisez au moins : \zh{因为}, \zh{虽然……可是……}, \zh{比如}.
\tcblower
\textbf{Raisonnement.} Plan : opinion → justification avec \zh{因为} → concession avec \zh{虽然……可是……} → exemple avec \zh{比如} → conclusion. On contrôle l'ordre des mots : sujet d'abord, compléments de temps et de lieu avant le verbe.

\textbf{Réponse rédigée (modèle).}

我认为年轻人应该关心政治。\\
\emph{Wǒ rènwéi niánqīngrén yīnggāi guānxīn zhèngzhì.}\\
« Je pense que les jeunes devraient s'intéresser à la politique. »

因为政治决定我们的未来，比如教育和工作的问题。\\
\emph{Yīnwèi zhèngzhì juédìng wǒmen de wèilái, bǐrú jiàoyù hé gōngzuò de wèntí.}\\
« Parce que la politique décide de notre avenir, par exemple les questions d'éducation et de travail. »

虽然我们很忙，可是我们可以在网上了解新闻。\\
\emph{Suīrán wǒmen hěn máng, kěshì wǒmen kěyǐ zài wǎngshàng liǎojiě xīnwén.}\\
« Bien que nous soyons très occupés, nous pouvons nous informer sur Internet. »

所以，关心政治也是公民的责任。\\
\emph{Suǒyǐ, guānxīn zhèngzhì yě shì gōngmín de zérèn.}\\
« Donc, s'intéresser à la politique est aussi un devoir de citoyen. »

\textbf{Justifications grammaticales.} \zh{在 + lieu + verbe} : \zh{在网上} précède le verbe \zh{了解} ; \zh{也} se place avant le verbe \zh{是} ; le sujet d'une phrase peut être une proposition verbale (\zh{关心政治} « s'intéresser à la politique » fonctionne ici comme sujet de \zh{是}).

\textbf{Conclusion.} La structure « opinion + cause + concession + exemple + conclusion » est un modèle transférable à tout sujet de débat du Bac.
\end{exoresolu}

\begin{experience}[Grand débat de classe : « La démocratie à l'école »]
Par groupes de quatre : deux élèves défendent la motion \zh{学生应该参加学校的管理} « les élèves devraient participer à la gestion de l'école », deux élèves la combattent. Chaque prise de parole doit contenir : une formule fonctionnelle (\zh{我觉得} / \zh{我同意} / \zh{我不太同意}), un argument avec \zh{因为}, et une relance (\zh{你觉得呢？}). Le reste de la classe note les tons fautifs et vote à la fin : \zh{我们投票！} (\emph{Wǒmen tóupiào!}) « Votons ! ».
\end{experience}

\begin{savaistu}[Cameroun et Chine : deux traditions civiques]
Au Cameroun, les élections et les débats publics rythment la vie citoyenne, et les jeunes y prennent une part croissante, notamment à travers les médias et les réseaux sociaux. En Chine, la participation citoyenne passe par d'autres formes, comme les assemblées de quartier et les consultations locales. Dans les deux pays, l'école joue un rôle central dans la formation du citoyen : apprendre à débattre avec respect — en chinois comme en français — est déjà un exercice de démocratie. Le mot \zh{民主} \emph{mínzhǔ} se lit littéralement « le peuple (\zh{民}) maître (\zh{主}) » : une belle image pour retenir le vocabulaire.
\end{savaistu}

\begin{attention}[Les 5 erreurs qui coûtent des points au Bac]
\begin{enumerate}
\item \cle{Tons négligés} : confondre \zh{民主} \emph{mínzhǔ} (2+3) avec une lecture à tons plats, ou oublier le sandhi du 3e ton dans \zh{选举} (prononcé \emph{xuánjǔ}). À l'oral, un ton faux peut rendre la phrase incompréhensible.
\item \cle{Caractères confondus} : \zh{民} (\emph{mín}, peuple) et \zh{明} (\emph{míng}, clair) ; \zh{权} (\emph{quán}, droit) et \zh{劝} (\emph{quàn}, conseiller) ; \zh{投} (\emph{tóu}, voter/jeter) et \zh{没} (\emph{méi}, ne pas). Vérifiez la clé : \zh{权} a la clé du bois 木, \zh{劝} a la clé de la force 力.
\item \cle{Ordre des mots calqué du français} : ne dites pas \zh{我同意与你} ; l'ordre est \zh{我同意你的看法} (Sujet + verbe + complément). De même, le lieu précède le verbe : \zh{在网上了解新闻}, jamais \zh{了解新闻在网上}.
\item \cle{Mots-outils oubliés ou mal placés} : le \zh{的} du complément du nom (\zh{公民的责任}, pas \zh{公民责任} à l'écrit soigné) ; \zh{都} qui se place avant le verbe (\zh{我们都应该投票}) ; \zh{也} avant le verbe, jamais en fin de phrase comme le « aussi » français.
\item \cle{Verbes modaux mélangés} : \zh{可以} = autorisation, \zh{能} = capacité, \zh{应该} = devoir moral. Ne traduisez pas « pouvoir » automatiquement par \zh{能} : \zh{人民可以投票} (le peuple a le droit de voter), pas \zh{人民能投票} qui suggère la capacité physique.
\end{enumerate}
\end{attention}

\begin{aretenir}[L'essentiel de la leçon]
\begin{itemize}
\item Un exposé oral se construit en trois temps annoncés : introduction (\zh{大家好！今天我要谈一谈……}), développement (\zh{首先}、\zh{其次}、\zh{最后}), conclusion (\zh{我认为……谢谢大家！}).
\item Le débat repose sur quatre gestes : donner son avis (\zh{我觉得} / \zh{我认为}), marquer l'accord (\zh{我同意你的看法}), le désaccord poli (\zh{我不太同意，因为……}), et relancer (\zh{你觉得呢？}).
\item L'argumentation s'appuie sur des structures fixes : \zh{因为……所以……}, \zh{虽然……可是……}, \zh{最 + adjectif}, \zh{比如}.
\item Les tons sont un contenu noté : sandhi du 3e ton (\zh{选举} → \emph{xuánjǔ}), contraste 2e/4e ton, ton neutre de \zh{个}.
\item Les verbes modaux ne sont pas interchangeables : \zh{可以} (autorisation), \zh{能} (capacité), \zh{应该} (devoir).
\end{itemize}
\end{aretenir}

\newpage

\coursec{Exercices}

\courssub{Je vérifie mes connaissances}

\begin{exercice}[QCM : le vocabulaire de la démocratie]
Pour chaque question, choisis la bonne réponse (a, b ou c).

\begin{enumerate}
\item \zh{民主} (mínzhǔ) signifie :
\begin{enumerate}
\item la liberté \item la démocratie \item la loi
\end{enumerate}
\item \zh{投票} (tóupiào) signifie :
\begin{enumerate}
\item voter \item parler \item écouter
\end{enumerate}
\item Le candidat aux élections se dit :
\begin{enumerate}
\item \zh{学生} \item \zh{班长} \item \zh{候选人}
\end{enumerate}
\item \zh{意见} (yìjiàn) signifie :
\begin{enumerate}
\item une réunion \item une opinion \item une question
\end{enumerate}
\end{enumerate}
\end{exercice}

\begin{exercice}[Vrai ou faux ? Justifie ta réponse]
Réponds par V (vrai) ou F (faux) et justifie en une phrase en français.

\begin{enumerate}
\item \zh{我觉得} (wǒ juéde) sert à donner son avis.
\item \zh{我同意} (wǒ tóngyì) exprime le désaccord.
\item \zh{你觉得呢？} (nǐ juéde ne ?) sert à relancer la parole d'un camarade.
\item Dans un exposé, \zh{总之} (zǒngzhī) s'emploie au début, pour introduire le sujet.
\end{enumerate}
\end{exercice}

\begin{exercice}[Vocabulaire : complète le tableau]
Complète le tableau suivant en écrivant le caractère ou la traduction manquante.

\begin{center}
\begin{tabularx}{\linewidth}{|X|X|X|X|}
\hline
\rowcolor{popblueL} Caractères & Pinyin & Nature & Traduction \\
\hline
民主 & mínzhǔ & nom & \ldots \\
\hline
\ldots & tóupiào & verbe & voter \\
\hline
选举 & \ldots & verbe / nom & élire / élection \\
\hline
\ldots & yìjiàn & nom & opinion, avis \\
\hline
权利 & quánlì & nom & \ldots \\
\hline
\ldots & zìyóu & nom / adj. & liberté / libre \\
\hline
\end{tabularx}
\end{center}
\end{exercice}

\begin{exercice}[Pinyin : associe les caractères à leur prononciation]
Associe chaque mot de la colonne A à son pinyin dans la colonne B.

\begin{center}
\begin{tabularx}{\linewidth}{|X|X|}
\hline
\rowcolor{popblueL} Colonne A (caractères) & Colonne B (pinyin) \\
\hline
1. 民主 & a. tóupiào \\
\hline
2. 投票 & b. hòuxuǎnrén \\
\hline
3. 候选人 & c. mínzhǔ \\
\hline
4. 权利 & d. fāyán \\
\hline
5. 发言 & e. quánlì \\
\hline
\end{tabularx}
\end{center}
\end{exercice}

\courssub{Je m'entraîne}

\begin{exercice}[Traduction : du français vers le chinois]
Traduis les phrases suivantes en chinois (caractères et pinyin).

\begin{enumerate}
\item Je pense que la démocratie est très importante.
\item Je suis d'accord avec toi. (utiliser \zh{同意})
\item Pourquoi les élèves doivent-ils voter ? Parce que c'est leur droit.
\item D'abord, je présente le sujet ; ensuite, je donne mon avis ; en conclusion, la démocratie nous concerne tous.
\end{enumerate}
\end{exercice}

\begin{exercice}[Dialogue à trous : les formules fonctionnelles]
Complète le dialogue suivant avec les formules de la boîte : \zh{我觉得} (wǒ juéde) — \zh{我同意} (wǒ tóngyì) — \zh{你觉得呢} (nǐ juéde ne) — \zh{总之} (zǒngzhī).

\begin{textebox}[Dialogue à compléter]
A：我们班要选班长，\ldots\ldots\ldots{}，李明是最好的候选人。\\
A：Wǒmen bān yào xuǎn bānzhǎng, \ldots\ldots\ldots, Lǐ Míng shì zuì hǎo de hòuxuǎnrén.\\
(A : Notre classe va élire un délégué, \ldots, Li Ming est le meilleur candidat.)\\[4pt]
B：\ldots\ldots\ldots{}！他学习好，也爱帮助同学。\\
B：\ldots\ldots\ldots{}！Tā xuéxí hǎo, yě ài bāngzhù tóngxué.\\
(B : \ldots ! Il travaille bien et il aime aider ses camarades.)\\[4pt]
A：王红也不错。\ldots\ldots\ldots{}？\\
A：Wáng Hóng yě búcuò. \ldots\ldots\ldots{}？\\
(A : Wang Hong est bien aussi. \ldots ?)\\[4pt]
B：\ldots\ldots\ldots{}，两个候选人都可以，我们星期五投票吧！\\
B：\ldots\ldots\ldots, liǎng ge hòuxuǎnrén dōu kěyǐ, wǒmen xīngqīwǔ tóupiào ba !\\
(B : \ldots, les deux candidats conviennent, votons vendredi !)
\end{textebox}
\end{exercice}

\begin{exercice}[Remise en ordre : un exposé sur les élections scolaires]
Remets les phrases suivantes dans l'ordre logique d'un exposé (numérote de 1 à 5), puis lis ton exposé à voix haute en faisant attention aux tons.

\begin{enumerate}
\item \zh{总之，选举班长是我们班的一件大事。} (Zǒngzhī, xuǎnjǔ bānzhǎng shì wǒmen bān de yí jiàn dàshì.)
\item \zh{大家好！今天我要谈一谈我们班的选举。} (Dàjiā hǎo ! Jīntiān wǒ yào tán yi tan wǒmen bān de xuǎnjǔ.)
\item \zh{然后，每个候选人要发言，介绍自己的想法。} (Ránhòu, měi ge hòuxuǎnrén yào fāyán, jièshào zìjǐ de xiǎngfǎ.)
\item \zh{首先，我们要知道什么是民主。} (Shǒuxiān, wǒmen yào zhīdào shénme shì mínzhǔ.)
\item \zh{最后，全班同学投票。} (Zuìhòu, quán bān tóngxué tóupiào.)
\end{enumerate}
\end{exercice}

\begin{exercice}[Discrimination des tons : entoure le mot entendu]
Le professeur lit un mot de chaque paire. Entoure le mot entendu, puis écris sa traduction.

\begin{enumerate}
\item mínzhǔ (民主) / mīnzhù
\item tóupiào (投票) / tǒupiāo
\item tóngyì (同意) / tōngyī
\item yìjiàn (意见) / yījiān
\item xuǎnjǔ (选举) / xuānjù
\end{enumerate}
\end{exercice}

\courssub{Je me prépare au Bac}

\begin{exercice}[Compréhension de l'écrit et questions de langue]
Lis le texte suivant, puis réponds aux questions.

\begin{textebox}[我们班的民主日 (Wǒmen bān de mínzhǔ rì)]
上个星期五，我们班举行了一次选举。首先，老师告诉我们什么是民主：民主就是每个人都有权利说话，也有权利投票。然后，三个候选人发言。李明说：「如果我当班长，我会帮助学习有困难的同学。」王红说：「我觉得班里应该更安静，也应该更干净。」最后，全班同学投票。王红得到了二十票，她当选了。总之，这次活动让我们明白：民主不只是投票，也是听别人的意见。\\
Shàng ge xīngqīwǔ, wǒmen bān jǔxíngle yí cì xuǎnjǔ. Shǒuxiān, lǎoshī gàosu wǒmen shénme shì mínzhǔ : mínzhǔ jiùshì měi ge rén dōu yǒu quánlì shuōhuà, yě yǒu quánlì tóupiào. Ránhòu, sān ge hòuxuǎnrén fāyán. Lǐ Míng shuō : « Rúguǒ wǒ dāng bānzhǎng, wǒ huì bāngzhù xuéxí yǒu kùnnan de tóngxué. » Wáng Hóng shuō : « Wǒ juéde bān lǐ yīnggāi gèng ānjìng, yě yīnggāi gèng gānjìng. » Zuìhòu, quán bān tóngxué tóupiào. Wáng Hóng dédàole èrshí piào, tā dāngxuǎn le. Zǒngzhī, zhè cì huódòng ràng wǒmen míngbai : mínzhǔ bù zhǐshì tóupiào, yěshì tīng biérén de yìjiàn.\\
(Vendredi dernier, notre classe a organisé une élection. D'abord, le professeur nous a expliqué ce qu'est la démocratie : la démocratie, c'est que chacun a le droit de parler et le droit de voter. Ensuite, trois candidats ont pris la parole. Li Ming a dit : « Si je deviens délégué, j'aiderai les camarades qui ont des difficultés. » Wang Hong a dit : « Je pense que la classe devrait être plus calme et plus propre. » Enfin, toute la classe a voté. Wang Hong a obtenu vingt voix, elle a été élue. En conclusion, cette activité nous a fait comprendre : la démocratie, ce n'est pas seulement voter, c'est aussi écouter l'avis des autres.)
\end{textebox}

\textbf{Questions de compréhension} (réponds en français) :
\begin{enumerate}
\item Quel événement a eu lieu dans la classe, et quand ?
\item Selon le professeur, quels sont les deux droits liés à la démocratie ?
\item Que promet Li Ming s'il est élu ?
\item Qui a gagné l'élection et avec combien de voix ?
\end{enumerate}

\textbf{Questions de langue} :
\begin{enumerate}
\item Relève dans le texte les trois connecteurs qui structurent l'exposé et traduis-les.
\item Releve la phrase contenant \zh{我觉得} et explique sa fonction.
\item Trouve dans le texte un synonyme de \zh{意见} (avis) : le mot \zh{想法}. Construis une phrase avec \zh{想法}.
\end{enumerate}
\end{exercice}

\begin{exercice}[Thème et version]
\textbf{Thème.} Traduis en chinois (caractères et pinyin) :
\begin{enumerate}
\item Je pense que tous les élèves devraient voter.
\item D'abord, les candidats prennent la parole ; ensuite, nous votons.
\item Parce que la démocratie est importante, nous devons écouter l'avis des autres.
\end{enumerate}

\textbf{Version.} Traduis en français le texte suivant :

\begin{textebox}[Version]
学生应该投票吗？我觉得应该。因为投票是我们的权利，所以每个人都应该参加。\\
Xuésheng yīnggāi tóupiào ma ? Wǒ juéde yīnggāi. Yīnwèi tóupiào shì wǒmen de quánlì, suǒyǐ měi ge rén dōu yīnggāi cānjiā.
\end{textebox}
\end{exercice}

\begin{exercice}[Expression orale et écrite : la démocratie dans ma classe]
\textbf{Partie A — Expression orale (préparation à l'exposé).} Prépare un exposé oral d'une minute sur le sujet : \zh{我们班的民主} (Wǒmen bān de mínzhǔ — La démocratie dans ma classe). Tu dois utiliser au moins trois connecteurs parmi : \zh{首先} (shǒuxiān), \zh{然后} (ránhòu), \zh{最后} (zuìhòu), \zh{总之} (zǒngzhī), et au moins une formule d'opinion : \zh{我觉得} ou \zh{我认为} (wǒ rènwéi).

\textbf{Partie B — Expression écrite.} Rédige ton exposé par écrit en chinois : \textbf{60 à 80 caractères} minimum. Structure attendue : salutation (\zh{大家好！}), introduction du sujet, développement avec connecteurs, conclusion avec \zh{总之}.

\textbf{Partie C — Jeu de rôle.} Par groupes de deux : l'un est journaliste, l'autre candidat aux élections de délégués. Le journaliste pose au moins trois questions (avec \zh{为什么} et \zh{你觉得呢？}) ; le candidat répond en justifiant avec \zh{因为……所以……}.
\end{exercice}

\newpage

\coursec{Corrigés des exercices}

\begin{corrige}[Exercice 1 — QCM : le vocabulaire de la démocratie]
\begin{enumerate}
\item \textbf{b} — \zh{民主} (mínzhǔ) signifie « la démocratie ».
\item \textbf{a} — \zh{投票} (tóupiào) signifie « voter ».
\item \textbf{c} — Le candidat aux élections se dit \zh{候选人} (hòuxuǎnrén). \zh{学生} (xuésheng) = élève ; \zh{班长} (bānzhǎng) = délégué de classe.
\item \textbf{b} — \zh{意见} (yìjiàn) signifie « une opinion, un avis ».
\end{enumerate}
\end{corrige}

\begin{corrige}[Exercice 2 — Vrai ou faux ? Justifie ta réponse]
\begin{enumerate}
\item \textbf{V (Vrai).} \zh{我觉得} (wǒ juéde) signifie « je pense que / à mon avis », c'est la formule standard pour introduire son opinion.
\item \textbf{F (Faux).} \zh{我同意} (wǒ tóngyì) signifie « je suis d'accord », il exprime l'\textbf{accord}. Le désaccord se dit \zh{我不同意} (wǒ bù tóngyì).
\item \textbf{V (Vrai).} \zh{你觉得呢？} (nǐ juéde ne ?) signifie « Et toi, qu'en penses-tu ? », c'est une formule de relance pour inviter l'interlocuteur à s'exprimer.
\item \textbf{F (Faux).} \zh{总之} (zǒngzhī) signifie « en conclusion / en résumé / bref », il s'emploie à la \textbf{fin} d'un exposé pour conclure. Pour introduire, on utilise \zh{首先} (shǒuxiān) ou \zh{今天我要谈一谈…} (Jīntiān wǒ yào tán yi tan…).
\end{enumerate}
\end{corrige}

\begin{corrige}[Exercice 3 — Vocabulaire : complète le tableau]
\begin{center}
\begin{tabularx}{\linewidth}{|X|X|X|X|}
\hline
\rowcolor{popblueL} Caractères & Pinyin & Nature & Traduction \\
\hline
民主 & mínzhǔ & nom & la démocratie \\
\hline
投票 & tóupiào & verbe & voter \\
\hline
选举 & xuǎnjǔ & verbe / nom & élire / élection \\
\hline
意见 & yìjiàn & nom & opinion, avis \\
\hline
权利 & quánlì & nom & le droit, les droits \\
\hline
自由 & zìyóu & nom / adj. & liberté / libre \\
\hline
\end{tabularx}
\end{center}
\end{corrige}

\begin{corrige}[Exercice 4 — Pinyin : associe les caractères à leur prononciation]
\begin{center}
\begin{tabularx}{\linewidth}{|X|X|}
\hline
\rowcolor{popblueL} Colonne A (caractères) & Colonne B (pinyin) \\
\hline
1. 民主 & c. mínzhǔ \\
\hline
2. 投票 & a. tóupiào \\
\hline
3. 候选人 & b. hòuxuǎnrén \\
\hline
4. 权利 & e. quánlì \\
\hline
5. 发言 & d. fāyán \\
\hline
\end{tabularx}
\end{center}
\end{corrige}

\begin{corrige}[Exercice 5 — Traduction : du français vers le chinois]
\begin{enumerate}
\item \zh{我觉得民主很重要。} \emph{(Wǒ juéde mínzhǔ hěn zhòngyào.)} — « Je pense que la démocratie est importante. »
\item \zh{我同意你的意见。} \emph{(Wǒ tóngyì nǐ de yìjiàn.)} \emph{ou simplement} \zh{我同意。} \emph{(Wǒ tóngyì.)} — « Je suis d'accord avec ton avis. / Je suis d'accord. »
\item \zh{为什么学生必须投票？因为这是他们的权利。} \emph{(Wèishéme xuésheng bìxū tóupiào? Yīnwèi zhè shì tāmen de quánlì.)} — « Pourquoi les élèves doivent-ils voter ? Parce que c'est leur droit. »
\item \zh{首先，我介绍题目；然后，我发表意见；总之，民主关乎我们大家。} \emph{(Shǒuxiān, wǒ jièshào tímù; ránhòu, wǒ fābiǎo yìjiàn; zǒngzhī, mínzhǔ guānhu wǒmen dàjiā.)} — « D'abord, j'introduis le sujet ; ensuite, j'exprime mon avis ; en conclusion, la démocratie nous concerne tous. »
\end{enumerate}
\end{corrige}

\begin{corrige}[Exercice 6 — Dialogue à trous : les formules fonctionnelles]
\textbf{Dialogue complété :}
\begin{itemize}
\item A：\zh{我们班要选班长，我觉得，李明是最好的候选人。} \\
\emph{Wǒmen bān yào xuǎn bānzhǎng, wǒ juéde, Lǐ Míng shì zuì hǎo de hòuxuǎnrén.} \\
(A : Notre classe va élire un délégué, \textbf{je pense que}, Li Ming est le meilleur candidat.)
\item B：\zh{我同意！他学习好，也爱帮助同学。} \\
\emph{Wǒ tóngyì! Tā xuéxí hǎo, yě ài bāngzhù tóngxué.} \\
(B : \textbf{Je suis d'accord} ! Il travaille bien et il aime aider ses camarades.)
\item A：\zh{王红也不错。你觉得呢？} \\
\emph{Wáng Hóng yě búcuò. Nǐ juéde ne?} \\
(A : Wang Hong est bien aussi. \textbf{Et toi, qu'en penses-tu ?})
\item B：\zh{总之，两个候选人都可以，我们星期五投票吧！} \\
\emph{Zǒngzhī, liǎng ge hòuxuǎnrén dōu kěyǐ, wǒmen xīngqīwǔ tóupiào ba!} \\
(B : \textbf{En conclusion / Bref}, les deux candidats conviennent, votons vendredi !)
\end{itemize}
\end{corrige}

\begin{corrige}[Exercice 7 — Remise en ordre : un exposé sur les élections scolaires]
L'ordre logique d'un exposé structuré (Introduction → Développement → Conclusion) est le suivant :
\begin{enumerate}
\item \zh{大家好！今天我要谈一谈我们班的选举。} \emph{(Dàjiā hǎo! Jīntiān wǒ yào tán yi tan wǒmen bān de xuǎnjǔ.)} — Salutation + annonce du sujet.
\item \zh{首先，我们要知道什么是民主。} \emph{(Shǒuxiān, wǒmen yào zhīdào shénme shì mínzhǔ.)} — Premier point : définition.
\item \zh{然后，每个候选人要发言，介绍自己的想法。} \emph{(Ránhòu, měi ge hòuxuǎnrén yào fāyán, jièshào zìjǐ de xiǎngfǎ.)} — Deuxième point : déroulement (prises de parole).
\item \zh{最后，全班同学投票。} \emph{(Zuìhòu, quán bān tóngxué tóupiào.)} — Troisième point : vote.
\item \zh{总之，选举班长是我们班的一件大事。} \emph{(Zǒngzhī, xuǎnjǔ bānzhǎng shì wǒmen bān de yí jiàn dàshì.)} — Conclusion.
\end{enumerate}
\textbf{Numérotation attendue :} 1 — 2 — 3 — 4 — 5.
\emph{Attention : l'ordre proposé par le collègue (2-4-3-5-1) était incorrect car il plaçait le vote avant les discours des candidats et la conclusion avant l'introduction.}
\end{corrige}

\begin{corrige}[Exercice 8 — Discrimination des tons : entoure le mot entendu]
Le professeur lit le mot aux tons standards (première proposition). Entourez le premier mot de chaque paire.
\begin{enumerate}
\item \textbf{\zh{民主} (mínzhǔ)} — démocratie \emph{(mīnzhù n'existe pas)}.
\item \textbf{\zh{投票} (tóupiào)} — voter \emph{(tǒupiāo : tons incorrects)}.
\item \textbf{\zh{同意} (tóngyì)} — être d'accord \emph{(tōngyī : tons incorrects)}.
\item \textbf{\zh{意见} (yìjiàn)} — opinion \emph{(yījiān : tons incorrects)}.
\item \textbf{\zh{选举} (xuǎnjǔ)} — élection \emph{(xuānjù : tons incorrects)}.
\end{enumerate}
\emph{Point de vigilance :} En chinois, le changement de ton change le sens du mot (ou donne un non-sens). Il est impératif de mémoriser le ton de chaque syllabe dès l'apprentissage du vocabulaire.
\end{corrige}

\begin{corrige}[Exercice 9 — Compréhension de l'écrit et questions de langue (Type Bac)]
\textbf{Barème indicatif : 10 points au total.}

\textbf{Questions de compréhension (4 pts) :}
\begin{enumerate}
\item Une élection de délégué de classe (\zh{选举班长}) a eu lieu \textbf{vendredi dernier} (\zh{上个星期五}). (1 pt)
\item Selon le professeur, la démocratie c'est : 1) \textbf{le droit de parler} (\zh{有权利说话}) et 2) \textbf{le droit de voter} (\zh{有权利投票}). (1 pt par droit = 2 pts)
\item Li Ming promet : \textbf{« Si je deviens délégué, j'aiderai les camarades qui ont des difficultés d'apprentissage. »} (\zh{如果我当班长，我会帮助学习有困难的同学。}) (1 pt)
\item C'est \textbf{Wang Hong} qui a gagné, avec \textbf{vingt voix} (\zh{二十票}). (1 pt)
\end{enumerate}

\textbf{Questions de langue (6 pts) :}
\begin{enumerate}
\item Les connecteurs structurants du texte sont (3 sur 4 présents) : 
\begin{itemize}
\item \zh{首先} (shǒuxiān) — \textbf{D'abord / Premièrement}
\item \zh{然后} (ránhòu) — \textbf{Ensuite / Puis}
\item \zh{最后} (zuìhòu) — \textbf{Enfin / Finalement}
\item \zh{总之} (zǒngzhī) — \textbf{En conclusion / En résumé / Bref}
\end{itemize}
(1 pt par connecteur identifié + traduction correcte = 3 pts)
\item La phrase est : \zh{「我觉得班里应该更安静，也应该更干净。」} \emph{(Wǒ juéde bān lǐ yīnggāi gèng ānjìng, yě yīnggāi gèng gānjìng.)} — « Je pense que la classe devrait être plus calme, et aussi plus propre. »
\textbf{Fonction :} \zh{我觉得} introduit l'\textbf{opinion personnelle} de la candidate (Wang Hong) dans le cadre de son discours de campagne. C'est une marque d'expression de la subjectivité. (2 pts : 1 pt citation, 1 pt explication)
\item Synonyme trouvé dans le texte : \zh{想法} (xiǎngfǎ) — « idée, pensée, point de vue ».
\textbf{Phrase construite :} \zh{你的想法很有意思，我同意。} \emph{(Nǐ de xiǎngfǎ hěn yǒu yìsi, wǒ tóngyì.)} — « Ton idée est très intéressante, je suis d'accord. » (1 pt)
\end{enumerate}
\end{corrige}

\begin{corrige}[Exercice 10 — Thème et version (Type Bac)]
\textbf{Barème indicatif : 10 points (5 pts thème, 5 pts version).}

\textbf{Thème — Traduction en chinois (caractères + pinyin) :}
\begin{enumerate}
\item \zh{我觉得所有学生都应该投票。} \emph{(Wǒ juéde suǒyǒu xuésheng dōu yīnggāi tóupiào.)}
\item \zh{首先，候选人发言；然后，我们投票。} \emph{(Shǒuxiān, hòuxuǎnrén fāyán; ránhòu, wǒmen tóupiào.)}
\item \zh{因为民主很重要，所以我们必须听别人的意见。} \emph{(Yīnwèi mínzhǔ hěn zhòngyào, suǒyǐ wǒmen bìxū tīng biérén de yìjiàn.)}
\end{enumerate}
\emph{Points de vigilance :} 
\begin{itemize}
\item Structure \zh{因为……所以……} (cause → conséquence) respectée.
\item \zh{所有} (suǒyǒu) pour « tous les », placé avant le nom.
\item Ponctuation chinoise ； pour séparer les propositions coordonnées dans la phrase 2.
\end{itemize}

\textbf{Version — Traduction en français :}
\begin{quote}
\textbf{« Les élèves doivent-ils voter ? Je pense que oui. Parce que voter est notre droit, donc tout le monde devrait participer. »}
\end{quote}
\emph{Analyse linguistique :}
\begin{itemize}
\item \zh{学生应该投票吗？} : Question à \zh{吗} (particule interrogative) = « Les élèves doivent-ils voter ? ».
\item \zh{我觉得应该。} : \zh{应该} (yīnggāi) reprend le verbe modal de la question, ellipsis du verbe \zh{投票} = « Je pense que oui / qu'ils le doivent. ».
\item \zh{因为……所以……} : Structure cause-conséquence standard.
\item \zh{每个人都应该参加} : \zh{每个人} (chacun) + \zh{都} (dōu, tous) + \zh{应该参加} (devrait participer).
\end{itemize}
\end{corrige}

\begin{corrige}[Exercice 11 — Expression orale et écrite : la démocratie dans ma classe (Type Bac)]
\textbf{Barème indicatif : 15 points (5 pts oral, 5 pts écrit, 5 pts interaction).}

\textbf{Partie A — Expression orale (Préparation) :} Pas de correction unique, voir modèle écrit ci-dessous qui sert de support.

\textbf{Partie B — Expression écrite (Modèle de rédaction — environ 85 caractères) :}
\begin{itemize}
\item \zh{大家好！今天我要谈一谈我们班的民主。} \emph{(Dàjiā hǎo! Jīntiān wǒ yào tán yi tan wǒmen bān de mínzhǔ.)}
\item \zh{首先，民主就是每个人都有权利说话，也有权利投票。} \emph{(Shǒuxiān, mínzhǔ jiùshì měi ge rén dōu yǒu quánlì shuōhuà, yě yǒu quánlì tóupiào.)}
\item \zh{然后，我们选举班长，候选人发言，介绍想法。} \emph{(Ránhòu, wǒmen xuǎnjǔ bānzhǎng, hòuxuǎnrén fāyán, jièshào xiǎngfǎ.)}
\item \zh{最后，全班同学投票。} \emph{(Zuìhòu, quán bān tóngxué tóupiào.)}
\item \zh{总之，民主不只是投票，也是听别人的意见，我觉得很重要。} \emph{(Zǒngzhī, mínzhǔ bù zhǐshì tóupiào, yěshì tīng biérén de yìjiàn, wǒ juéde hěn zhòngyào.)}
\item \zh{谢谢大家！} \emph{(Xièxie dàjiā!)}
\end{itemize}
\textbf{Vérification des contraintes :}
\begin{itemize}
\item Salutation : \zh{大家好！} ✓
\item Connecteurs (4/4) : \zh{首先}、\zh{然后}、\zh{最后}、\zh{总之} ✓
\item Formule d'opinion : \zh{我觉得} ✓
\item Longueur : ~85 caractères chinois ✓
\item Structure complète : Intro — Développement — Conclusion ✓
\end{itemize}

\textbf{Partie C — Jeu de rôle (Modèle de dialogue) :}
\textbf{Journaliste (J) / Candidat (C)}
\begin{itemize}
\item J：\zh{你好，你是候选人吗？为什么想当班长？} \emph{(Nǐ hǎo, nǐ shì hòuxuǎnrén ma? Wèishéme xiǎng dāng bānzhǎng?)}
\item C：\zh{你好，我是候选人。因为我想为同学服务，所以想当班长。} \emph{(Nǐ hǎo, wǒ shì hòuxuǎnrén. Yīnwèi wǒ xiǎng wèi tóngxué fúwù, suǒyǐ xiǎng dāng bānzhǎng.)}
\item J：\zh{你觉得班里最大的问题是什么？} \emph{(Nǐ juéde bān lǐ zuì dà de wèntí shì shénme?)}
\item C：\zh{我觉得班里不够安静。因为吵闹影响学习，所以我们要制定规则。} \emph{(Wǒ juéde bān lǐ bù gòu ānjìng. Yīnwèi chǎonào yǐngxiǎng xuéxí, suǒyǐ wǒmen yào zhìdìng guīzé.)}
\item J：\zh{如果你当选，你会做什么？你觉得呢？} \emph{(Rúguǒ nǐ dāngxuǎn, nǐ huì zuò shénme? Nǐ juéde ne?)}
\item C：\zh{我会组织学习小组，帮助有困难的同学。因为互助让班级更团结，所以我会努力做。} \emph{(Wǒ huì zǔzhī xuéxí xiǎozǔ, bāngzhù yǒu kùnnan de tóngxué. Yīnwèi hùzhù ràng bānjí gèng tuánjié, suǒyǐ wǒ huì nǔlì zuò.)}
\end{itemize}
\textbf{Points de vigilance pour l'évaluation :}
\begin{itemize}
\item \textbf{Prononciation des tons :} Vérifier particulièrement \zh{民主} (mínzhǔ, 2-3), \zh{投票} (tóupiào, 2-4), \zh{候选人} (hòuxuǎnrén, 4-3-2), \zh{权利} (quánlì, 2-4), \zh{觉得} (juéde, 2-neutre/léger 3), \zh{同意} (tóngyì, 2-4).
\item \textbf{Fluidité :} Liaison naturelle des connecteurs (\zh{首先}、\zh{然后}…).
\item \textbf{Interaction :} Le journaliste utilise \zh{为什么} et \zh{你觉得呢？} ; le candidat utilise \zh{因为……所以……}.
\item \textbf{Caractères :} Écriture correcte des clés : \zh{言} (yán, parole) dans \zh{意见}、\zh{发言} ; \zh{示} (shì, montrer) dans \zh{票} ; \zh{人} (rén) dans \zh{候选人}.
\end{itemize}
\end{corrige}

\newpage

\coursec{Fiche bilan}

\courssub{Fiche Bilan — Leçon 24 : Expression orale — Démocratie}

\begin{popfigure}
\begin{tikzpicture}[
  mindmap,
  grow cyclic,
  every node/.style={concept, circular drop shadow, execute at begin node=\strut},
  root concept/.append style={concept color=popblue, font=\bfseries\large, text width=3.8cm, align=center},
  level 1 concept/.append style={level distance=4.8cm, sibling angle=360/7, font=\small, text width=2.4cm, align=center},
  level 2 concept/.append style={level distance=3.2cm}
]
\node[root concept, align=center]{Leçon 24\\ Démocratie\\(Expression orale)}
  child[concept color=popblue] { node[align=center]{Vocabulaire\\ clé} }
  child[concept color=poporange] { node[align=center]{Donner\\ son avis} }
  child[concept color=popgreen] { node {Accord /\textbackslash Désaccord} }
  child[concept color=poppurple] { node[align=center]{Relancer\\ le débat} }
  child[concept color=poppink] { node[align=center]{Structurer\\ un exposé} }
  child[concept color=popgold] { node[align=center]{Prononciation\\ Tons \& Sandhi} }
  child[concept color=popdark] { node[align=center]{Contexte\\ Cameroun / Chine} };
\end{tikzpicture}
\legende{Carte mentale récapitulative de la Leçon 24 (Module 3 : Citoyenneté)}
\end{popfigure}

\begin{definition}[Lexique clé (HSK 3--4) — Thème : Démocratie \& Citoyenneté]
\begin{center}
\begin{tabularx}{\linewidth}{|X|X|X|X|}
\hline
\rowcolor{popblueL} \textbf{Caractères} & \textbf{Pinyin} & \textbf{Nature} & \textbf{Traduction} \\ \hline
民主 & mínzhǔ & n. & démocratie \\ \hline
选举 & xuǎnjǔ & v./n. & élire / élection \\ \hline
投票 & tóupiào & v./n. & voter / vote \\ \hline
权利 & quánlì & n. & droit \\ \hline
义务 & yìwù & n. & devoir, obligation \\ \hline
公民 & gōngmín & n. & citoyen \\ \hline
宪法 & xiànfǎ & n. & constitution \\ \hline
言论自由 & yánlùn zìyóu & n. & liberté d'expression \\ \hline
参与 & cānyù & v. & participer \\ \hline
代表 & dàibiǎo & n./v. & représentant / représenter \\ \hline
制度 & zhìdù & n. & système, régime politique \\ \hline
透明 & tòumíng & adj. & transparent \\ \hline
平等 & píngděng & adj. & égalitaire, égal \\ \hline
监督 & jiāndū & v./n. & superviser / contrôle, supervision \\ \hline
决策 & juécè & n./v. & décision / décider \\ \hline
\end{tabularx}
\end{center}
\end{definition}

\begin{propriete}[Structures fonctionnelles à connaître par cœur (Grammaire de l'oral)]
\begin{center}
\begin{tabularx}{\linewidth}{|X|X|X|}
\hline
\rowcolor{poporangeL} \textbf{Fonction} & \textbf{Formule type (Structure)} & \textbf{Exemple (Caractères + Pinyin + Traduction)} \\ \hline
\bfseries Présenter son exposé & 今天我想谈谈…… / 我的主题是…… & \zh{今天我想谈谈民主。} (Jīntiān wǒ xiǎng tántan mínzhǔ.) « Aujourd'hui je veux parler de la démocratie. » \\ \hline
\bfseries Donner son avis & 我认为…… / 就我而言…… / 据我所知…… & \zh{我认为投票是公民的权利。} (Wǒ rènwéi tóupiào shì gōngmín de quánlì.) « Je pense que voter est un droit du citoyen. » \\ \hline
\bfseries Exprimer l'accord & 我同意…… / 完全正确 / 很有道理 & \zh{我同意你的看法，很有道理。} (Wǒ tóngyì nǐ de kànfǎ, hěn yǒu dàolǐ.) « Je suis d'accord avec ton avis, c'est très sensé. » \\ \hline
\bfseries Exprimer le désaccord (politesse) & 我不同意…… / 我不太赞成…… / 话虽这么说，可是…… & \zh{话虽这么说，可是权利要伴随义务。} (Huà suī zhème shuō, kěshì quánlì yào bànsuí yìwù.) « Certes, mais les droits doivent s'accompagner de devoirs. » \\ \hline
\bfseries Relancer / Interroger & 你怎么看？ / 为什么呢？ / 还有别的意见吗？ & \zh{你怎么看这个问题？} (Nǐ zěnme kàn zhège wèntí?) « Qu'en penses-tu de cette question ? » \\ \hline
\bfseries Structurer (connecteurs) & 首先……；其次……；最后…… / 总之…… & \zh{首先，民主是权利；其次，需要参与；总之，人人有责。} (Shǒuxiān, mínzhǔ shì quánlì; qícì, xūyào cānyù; zǒngzhī, rénrén yǒu zé.) « Premièrement, la démocratie est un droit ; deuxièmement, elle requiert la participation ; en somme, chacun a sa part de responsabilité. » \\ \hline
\end{tabularx}
\end{center}
\end{propriete}

\begin{methode}[Méthodes express pour l'oral (Bac \& Contrôle continu)]
\begin{enumerate}
  \item \textbf{Préparation (5 min) :} Notez des \cle{mots-clés} (caractères + pinyin) sur une fiche, \textbf{pas des phrases entières}. Structure imposée : \textit{Introduction → 2 arguments → Conclusion}. Exemple de fiche : \zh{民主、权利、义务、参与、选举} + connecteurs \zh{首先、其次、总之}.
  \item \textbf{Prononciation — Sandhi des tons (ciblé) :}
        \begin{itemize}
          \item \cle{3e ton + 3e ton → 2e ton + 3e ton} : \zh{选举} (xuǎnjǔ $\rightarrow$ \textbf{xuán}jǔ) ; \zh{监督} (jiāndū, 1+1 inchangé) ; \zh{代表} (dàibiǎo, 4+3 inchangé). \textbf{Attention :} \zh{民主} (mínzhǔ) = 2e + 3e \textbf{pas de sandhi} (民 est 2e ton).
          \item \cle{\zh{不} (bù)} : devient \zh{bú} (2e ton) devant un 4e ton (ex: \zh{不赞成} bú zànchéng) ; reste \zh{bù} (4e ton) sinon.
          \item \cle{\zh{一} (yī)} : 1e ton seul / à la fin / avant chiffre ; 2e ton (\zh{yí}) devant 4e ton ; 4e ton (\zh{yì}) devant 1e, 2e, 3e ton.
          \item \cle{Tons neutres} : particules \zh{了、吗、呢、吧} ; suffixes \zh{们、子、头} (ex: \zh{代表们} dàibiǎomen).
        \end{itemize}
  \item \textbf{Débit et prosodie :} Parlez par \cle{groupes respiratoires} (感群, gǎnqún). Marquez une pause nette après les connecteurs structurants : \zh{首先 / 其次 / 最后 / 总之}. Évitez le débit syllabe par syllabe.
  \item \textbf{Interaction et stratégies :} Regardez l'interlocuteur. Pour relancer : \zh{你呢？} / \zh{你怎么看？}. Pour gagner du temps : \zh{这个问题很好……} / \zh{让我想想……}. Pour reformuler : \zh{也就是说……} / \zh{换句话说……}.
  \item \textbf{Ancrage contextuel (Cameroun / Chine) :} Intégrez \textbf{un fait précis} par exposé.
        \begin{itemize}
          \item \textbf{Cameroun} : \zh{喀麦隆的市政选举} (élections municipales), \zh{雅温得} (Yaoundé), \zh{投票年龄20岁} (âge de vote 20 ans), \zh{多党制} (multipartisme).
          \item \textbf{Chine} : \zh{全国人民代表大会} (Assemblée populaire nationale), \zh{人民代表} (députés du peuple), \zh{基层民主协商} (concertation démocratique à la base), \zh{全过程人民民主} (démocratie populaire tout au long du processus).
        \end{itemize}
\end{enumerate}
\end{methode}

\begin{aretenir}[Checklist avant le Bac — Leçon 24]
$\square$ Je connais par cœur le vocabulaire clé (caractères, pinyin, sens) : \zh{民主、选举、投票、权利、义务、公民、宪法、言论自由、参与、代表、制度、透明、平等、监督、决策}.\par
$\square$ Je sais présenter mon exposé : \zh{今天我想谈谈…… / 我的主题是……}.\par
$\square$ Je sais donner mon avis avec \zh{我认为…… / 就我而言…… / 据我所知……}.\par
$\square$ Je sais exprimer l'accord : \zh{我同意…… / 很有道理 / 完全正确}.\par
$\square$ Je sais exprimer le désaccord poliment : \zh{我不太赞成…… / 话虽这么说，可是……}.\par
$\square$ Je sais relancer le débat : \zh{你怎么看？ / 为什么呢？ / 还有别的意见吗？}.\par
$\square$ Je structure mon discours avec \zh{首先…… 其次…… 最后…… / 总之……}.\par
$\square$ Je maîtrise les \textbf{sandhi de tons} : 3e+3e $\rightarrow$ 2e+3e (ex: \zh{选举} xuánjǔ) ; \zh{不} bù $\rightarrow$ bú devant 4e ton ; \zh{一} yī selon contexte.\par
$\square$ Je distingue les paires phonétiques pièges (tons / finales / initiales) : \zh{参与} (cānyù) vs \zh{残余} (cányú) ; \zh{制度} (zhìdù) vs \zh{质度} (zhìdù) ; \zh{监督} (jiāndū) vs \zh{简单} (jiǎndān) ; \zh{权利} (quánlì) vs \zh{全力} (quánlì).\par
$\square$ Je peux citer un fait \textbf{Cameroun} (ex: \zh{雅温得市政选举}, vote à 20 ans) et un fait \textbf{Chine} (ex: \zh{全国人大代表}, concertation à la base).\par
$\square$ Je parle à un débit naturel, sans lire mes notes, en regardant mes camarades (contact visuel).\par
$\square$ Je réagis en chinois aux interventions des autres : \zh{对 / 不对 / 我也这么觉得 / 我不太确定 / 有道理}.\par
\end{aretenir}

\findecours

\end{document}
